%------------------------------------------------------------------------------%
% LINEAR OPERATORS AND SPECTRAL THEORY                                         %
%------------------------------------------------------------------------------%
% File  : Linear_Operators_and_Spectral_Theory.tex                             %
% Author: Klaus Hoeflinger, Beethovenstrasse 11, D-73117 Wangen                %
% Email : khoeflinger@t-online.de                                              %
%------------------------------------------------------------------------------%

%------------------------------------------------------------------------------%
%                       DOCUMENT CLASS and INCLUDE FILES                       %
%------------------------------------------------------------------------------%
\documentclass[11pt,twoside]{article}
\usepackage{geometry}
\geometry{paperheight=241mm,
          paperwidth=152mm,
          top=22mm,
          bottom=17mm,
          left=20mm,
          right=20mm,
          textwidth=112mm,
          textheight=202mm,
          headheight=10mm,
          headsep=3mm,
          footskip=9mm,
          includehead,
          includefoot,
          marginparsep=0mm,
          marginparwidth=0mm}

\renewcommand{\baselinestretch}{0.945}\normalsize

\AtBeginDocument{\setlength\abovedisplayskip{2mm}}
\AtBeginDocument{\setlength\belowdisplayskip{2mm}}

%------------------------------------------------------------------------------%
% define title formats                                                         %
%------------------------------------------------------------------------------%
\usepackage{titlesec}

\renewcommand{\thesection}{\arabic{section}}
\renewcommand{\sfdefault}{FiraSans}
\renewcommand{\ttdefault}{pcr}
\renewcommand{\rmdefault}{antiqua}

\titleformat{\part}[display]
{\normalfont \large \rmfamily \bfseries \centering}
{\small{\thepart.}}{2pt}{}

\titleformat{\section}[runin]
{\normalfont \rmfamily \bfseries}
{\thesection}{1pt}{.\;}[.]

%\titleformat{\section}
%{\normalfont \bfseries \small \filcenter}
%{\thesection.\;}{0mm}{}

\titleformat{\subsection}[runin]
{\normalfont \itshape}
{}{0mm}{}

\titleformat{\subsubsection}[runin]
{\normalfont \itshape}
{}{}{}{}

\titleformat{\paragraph}[runin]
{\normalfont}
{}{}{}{}

\titleformat{\subparagraph}[runin]
{\normalfont}
{}{}{}{}

\titlespacing*{\part}          { 0pt}{ 0pt}{ 8pt}
\titlespacing*{\section}       { 0pt}{ 7pt}{ 4pt}
\titlespacing*{\subsection}    { 0pt}{14pt}{ 6pt}
\titlespacing*{\subsubsection} { 0pt}{ 6pt}{12pt}
\titlespacing*{\paragraph}     { 0pt}{ 6pt}{12pt}
\titlespacing*{\subparagraph}  { 0pt}{ 0pt}{12pt}

%------------------------------------------------------------------------------%
% define enumerate parameters                                                  %
%------------------------------------------------------------------------------%
\usepackage{enumitem}
\setlength{\parindent}{3pt}
\setlength{\parskip}  {0pt}

%\setlist{noitemsep}

\setlist[enumerate]{
         leftmargin=12pt,
         rightmargin=0pt,
         itemsep=1pt,
         itemindent=3pt,
         topsep=3pt,
         labelsep=4pt,
         partopsep=0pt,
         labelwidth=0pt,
         listparindent=0pt,
         parsep=0pt}

\setlist[itemize]{
         leftmargin=12pt,
         rightmargin=0pt,
         itemsep=0pt,
         itemindent=1pt,
         topsep=1pt,
         labelsep=4pt,
         partopsep=1pt,
         labelwidth=0pt,
         listparindent=0pt,
         parsep=0pt}

%------------------------------------------------------------------------------%
% define packages                                                              %
%------------------------------------------------------------------------------%
\usepackage{upref}
\usepackage{amssymb}
\usepackage{slantsc}
\usepackage{amsmath}
\usepackage{amsthm}
\usepackage{thmtools}
\usepackage{amscd}
\usepackage{verbatim}
\usepackage{latexsym}
\usepackage{multicol}
\usepackage{graphicx}
\usepackage{setspace}
\usepackage[ngerman,english]{babel}
\usepackage[protrusion=true,expansion=true]{microtype}
\usepackage{tikz}
\usetikzlibrary{arrows,arrows.meta,decorations.markings}
\usepackage{mathrsfs}
\usepackage{amsfonts}
\usepackage{datetime2}
\usepackage{textgreek}
\usepackage{hyperref}
\usepackage{pifont}
\usepackage{relsize}
%\usepackage{biblatex}
\relsize{-0.05}
\usepackage[skip=0mm]{parskip}
\usetikzlibrary{decorations.pathmorphing}

\usepackage{mlmodern}
\usepackage[T5,T1]{fontenc}
\usepackage{ragged2e}

%\usepackage{mathabx}
%\usepackage{times}

%\usepackage{fontspec}
%\setmainfont{Nimbus Roman No9 L}
%\usepackage[T1]{fontenc}

%------------------------------------------------------------------------------%
% define page layout                                                           %
%------------------------------------------------------------------------------%
\usepackage{fancyhdr}
\pagestyle{fancy}
\setlength{\headheight}{30.0pt}
\renewcommand{\headrulewidth}{0pt}

\AtBeginDocument{
  \addtolength\abovedisplayskip{-0.0\baselineskip}
  \addtolength\belowdisplayskip{-0.0\baselineskip}
}

\fancyhead{}
\fancyfoot{}
\addtolength{\headwidth}{0.02\marginparwidth}

\makeatletter
\let\ps@plain\ps@fancy
\makeatother

\renewcommand\sectionmark[1]
{
 \def\sectionname{#1}
 \markright{\thesection #1}
}

\makeatletter
\@addtoreset{section}{part}
\makeatother

%------------------------------------------------------------------------------%
% define footnote without number                                               %
%------------------------------------------------------------------------------%
\newcommand{\myfootnote}[1]{%
\renewcommand{\thefootnote}{}%
\footnotetext{#1}%
\renewcommand{\thefootnote}{\arabic{footnote}}%
}

%------------------------------------------------------------------------------%
% define numbering in equations                                                %
%------------------------------------------------------------------------------%
\numberwithin{equation}{section}

%------------------------------------------------------------------------------%
% define newtheorem                                                            %
%------------------------------------------------------------------------------%
\declaretheoremstyle[
headfont=\scshape, 
headindent = 0mm, %\parindent,
%postheadhook = {\hspace{0mm}\newline},
%postheadspace = \newline, 
spaceabove = 3mm, 
spacebelow = 3mm,
numberwithin=section,
name=Theorem]{khMATH}
\declaretheorem[style=khMATH]{theorem}

\declaretheoremstyle[
headfont=\bfseries, 
headindent = 0mm, %\parindent,
%postheadhook = {\hspace{0mm}\newline},
%postheadspace = \newline, 
spaceabove = 3mm, 
spacebelow = 3mm,
numberwithin=part,
name=Theorem]{khMATH1}
\declaretheorem[style=khMATH1]{theorem1}

\declaretheoremstyle[
headfont=\scshape, 
headindent = 0mm, %\parindent,
%postheadhook = {\hspace{0mm}\newline},
%postheadspace = \newline, 
spaceabove = 3mm, 
spacebelow = 3mm,
numberwithin=section,
name=Corollary]{khMATH1}
\declaretheorem[style=khMATH1]{corollary}

%            name         counter     label              numeration-mode
\theoremstyle{plain}
\newtheorem*{introduction}            {\sc Introduction}
\newtheorem {standard}                {}                 [section]
\newtheorem*{definition}              {\sc Definition}
\newtheorem{satz1}                    {\sc Satz}
\newtheorem{lemma}                    {\sc Lemma}
\newtheorem*{proposition}             {\sc Proposition}
%\newtheorem{corollary}               {\sc Corollary}
\newtheorem*{corollar}                {\sc Korollar}
\newtheorem*{remark}                  {\sc Remark}
\newtheorem*{remarks}                 {\sc Remarks}
\newtheorem*{example}                 {Example}
\newtheorem*{examples}                {Examples}
\newtheorem*{theo}                    {\sc Theorem}

%------------------------------------------------------------------------------%
% define tabular                                                               %
%------------------------------------------------------------------------------%
\usepackage{tabularx}
\newcolumntype{L}[1]{>{\raggedright\arraybackslash}p{#1}}
\newcolumntype{C}[1]{>{\centering\arraybackslash}  p{#1}} 
\newcolumntype{R}[1]{>{\raggedleft\arraybackslash} p{#1}} 

%------------------------------------------------------------------------------%
% define \sh, \zB                                                              %
%------------------------------------------------------------------------------%
\def\dsh{{d.\,h.}}
\def\ie{{i.\,e.}}
\def\zB{z.\,B.}
\renewcommand{\abstractname}{ }

%------------------------------------------------------------------------------%
% change default spacing for cases                                             %
%------------------------------------------------------------------------------%
\makeatletter
\renewcommand*\env@cases[1][1.6]{%
  \let\@ifnextchar\new@ifnextchar
  \left\lbrace
  \def\arraystretch{#1}%
  \array{@{}l@{\quad}l@{}}%
}
\makeatother

\newenvironment{rcases}
  {\left.\begin{aligned}}
  {\end{aligned}\right\rbrace}

%------------------------------------------------------------------------------%
% change default for \predisplaypenalty                                        %
%------------------------------------------------------------------------------%
\predisplaypenalty=990
\allowdisplaybreaks \numberwithin{equation}{section}

%------------------------------------------------------------------------------%
% general notations                                                            %
%------------------------------------------------------------------------------%
\def\*{\star}
\def\={\,=\,}
\def\+{\,+\,}
\def\xsubset{{\,\subset\,}}
\def\xtimes{{\,\times\,}}
\def\xeof{{\,\in\,}}
\def\eq{{\,=\,}}
\def\xto{{\,\to\,}}
\def\eqdef{\,:=\,}
\def\:{{\,:\,}}
\def\ele{{\,\in\,}}
\def\exists{\mbox{ exists }}
\def\and{\mbox{ and }}
\def\for{\mbox{ for }}
\def\fa{\mbox{ for all }}
\def\fe{\mbox{ for each }}
\def\qfa{\quad \mbox{ for all }}
\def\df{\,\dotfill}
\def\iff{if and only if }
\def\wrt{with respect to}

\makeatletter
\newcommand \Df {\leavevmode \cleaders \hb@xt@ .70em{\hss .\hss }\hfill \kern \z@}
\makeatother

\def\sql{\mathfrak{a}}               % sesquilinear form a
\def\DF{B}                           % Dirichlet form a
%------------------------------------------------------------------------------%
% number sets and special elements                                             %
%------------------------------------------------------------------------------%
\def\P{\mathbb{H}}                   % half plane of $\C$
\def\J{I}                            % interval I
\def\N{{\mathbb{N}}}                 % unsigned integers
\def\Q{{\mathbb{Q}}}                 % rational integers
\def\Z{{\mathbb{Z}}}                 % integers
\def\R{{\mathbb{R}}}                 % real numbers
\def\Rn{{\mathbb{R}}^n}              % real numbers in Rn
\def\Rd{{\mathbb{R}}^d}              % real numbers in Rn
\def\C{{\mathbb{C}}}                 % complex numbers
\def\F{{\mathbb{K}}}                 % arbitrary field
\def\Torus{{\mathbb{T}}}             % complex circle line
\def\T{\tau}                         % upper bound of interval (0,t)
\def\sign#1{\operatorname{sign}(#1)} % sign of a number

%------------------------------------------------------------------------------%
% set theory                                                                   %
%------------------------------------------------------------------------------%
\def\G{G}                            % domain $G$
\def\clG{\overline{\G}}              % closure of $G$
\def\bndG{\partial \G}               % boundary of $G$
\def\setA{\mathcal{A}}               % set A
\def\setB{\mathcal{B}}               % set B
\def\setC{\mathcal{C}}               % set C
\def\setM{M}                         % set M
\def\setN{N}                         % set N
\def\setS{\mathcal{S}}               % set S
\def\famF{\mathfrak{F}}              % family of sets
\def\indI{\mathcal{I}}               % index set I
\def\pot#1{\mathfrak{P}(#1)}         % power set

\def\relR{\mathbf{R}}                % relation R
\def\mapf{f}                         % mapping f
\def\mapg{g}                         % mapping g

\def\set#1#2{\{ \, #1 \,: \, #2 \, \}}
\def\tensor#1#2{#1 \otimes #2}
\def\Tens{\oplus}

\def\cal#1{{\mathcal{#1}}}
\def\aut#1{\operatorname{Aut}(#1)}

\def\cross#1#2{#1 {\,\times\,} #2}
\def\multicross#1#2{#1 \times  \ldots \times  #2}
\def\multitens#1#2 {#1 \otimes \ldots \otimes #2}

%------------------------------------------------------------------------------%
% group theory                                                                 %
%------------------------------------------------------------------------------%
\def\1{{\mathsf{1}}}                % unit element of a monoid
\def\0{{\mathsf{0}}}                % unit element of an Abelian monoid
\def\kernel#1{\mathfrak{N}(#1)}     % kernel of a homomorphism
\def\cokernel#1{\mathfrak{C}{(#1)}} % cokernel of a homomorphism
\def\Hom#1{\operatorname{Hom}(#1)}

%------------------------------------------------------------------------------%
% linear algebra                                                               %
%------------------------------------------------------------------------------%
\def\vecV{\mathit{V}}               % vector space V
\def\vecH{\mathit{H}}               % vector space H
\def\vecW{\mathit{W}}               % vector space W
\def\vecU{\mathit{U}}               % subspace U
\def\convC{\mathcal{C}}             % convex set C
\def\basH{\mathfrak{B}}             % Hamel base B

\def\LO#1{\mathscr{L}(#1) }         % algebra of linear transformations

\def\scalar#1#2{ \langle #1,#2 \rangle }
\def\trace#1{\operatorname{tr}(#1)}
\def\dim#1{{\operatorname{dim} \, #1 }}
\def\codim#1{{\operatorname{codim} \, #1 }}
\def\dim{\operatorname{dim}}
\def\codim{\operatorname{codim}}
\def\ind{\operatorname{ind}}

\def\genG{\cal{G}}

\def\algA{\cal{A}}
\def\algB{\cal{B}}
\def\algU{\cal{U}}

\def\idI{\cal{I}}

%------------------------------------------------------------------------------%
% topology                                                                     %
%------------------------------------------------------------------------------%
\def\compK{{K}}          % compact set C
\def\openO{\mathcal{O}}  % open set O
\def\openU{\mathcal{U}}  % open set U
\def\U{\mathcal{U}}      % neighbourhood U
\def\V{\mathcal{V}}      % neighbourhood V
\def\d{\mathit{d}}       % metric function d

\def\interior#1{\mathring{#1}}
\def\closure#1{{\overline{#1}}}

\def\topT{\mathcal{T}}   % topology T
\def\topX{X}             % topological space X
\def\topY{Y}             % topological space Y
\def\topZ{Z}             % topological space Z
\def\metrS{M}            % metric space M

\def\grpG{G}             % topological group G
\def\grpA{A}             % topological group A
\def\grpB{B}             % topological group B
\def\grpC{C}             % topological group C
\def\grpH{H}             % topological group H
\def\grpU{U}             % topological group U
\def\grpV{V}             % topological group V
\def\topG{G}             % topological group G
\def\topH{H}             % topological group H
\def\topU{U}             % topological subgroup U

\def\subS{\cal{S}}
\def\riR{R}              % ring R

%------------------------------------------------------------------------------%
% calculus                                                                     %
%------------------------------------------------------------------------------%
\def\supp#1{\operatorname{supp}(#1)}
\def\singsupp#1{\operatorname{sing \, supp}(#1)}

\def\re{\operatorname{Re}}
\def\im{\operatorname{Im}}
\def\grad{\operatorname{grad}}

%\def\abs#1 {\left|#1\right|}
\def\abs#1 {\lvert #1 \rvert}
%\def\abs#1 {\left|\,#1\,\right|}

%------------------------------------------------------------------------------%
% measure theory                                                               %
%------------------------------------------------------------------------------%
\def\salgA{\Sigma}               % sigma-algebra S
\def\salgB{\cal{B}}              % sigma-algebra B
\def\Borel#1{\cal{B}_{#1}}       % Borel-algebra 
\def\LB{\lambda}                 % Lebesgue-Borel measure

%------------------------------------------------------------------------------%
% function spaces                                                              %
%------------------------------------------------------------------------------%
\def\Lp{L^p(\G)}                 % spaces L_p
\def\Lq{L^q(\G)}                 % spaces L_q
\def\Lh{L^2(\G)}                 % spaces L_2
\def\Li{L^{\infty}(\G)}          % spaces L_infty
\def\Lc{L_{loc}^1(\G)}           % spaces L_1^{loc}
\def\Ck{C^k(\G)}                 % spaces C_k

\def\BV#1{BV(#1) }               % set of functions of bounded variation
\def\AC#1{AC(#1) }               % set of absolutely continuous functions

\def\MP#1{\mathscr{M}^+(#1) }    % set of positive measures
\def\MS#1{\mathscr{M}^-(#1) }    % vector space of finite signed measures
\def\MC#1{\mathscr{M}(#1) }      % vector space of complex measures
\def\MCR#1{\mathscr{M}_{r}(#1) } % vector space of regular complex measures

\def\SobW{W^{k,p}(\G)}                % Sobolev space W
\def\SobO{\mathring{W}^{k,p}(\G)}     % Sobolev space W
%\def\WN#1#2{\mathring{W}^{#1}(#2) }   % Sobolev space W
\def\WN#1#2{W_0^{#1}(#2) }   % Sobolev space W
%\def\HN#1#2{\mathring{H}^{#1}(#2) }   % Sobolev space H
\def\HN#1#2{H_0^{#1}(#2) }   % Sobolev space H
%------------------------------------------------------------------------------%
% functional analysis                                                          %
%------------------------------------------------------------------------------%
\def\snorm#1{\lVert #1 \rVert}
\newcommand{\norm}[1]{\left\lVert #1 \right\rVert}

\def\topV{V}                       % topological vector space V
\def\topW{W}                       % topological vector space W
\def\normS{X}                      % normed space

\def\banX{\mathit{X}}              % Banach space X
\def\banY{\mathit{Y}}              % Banach space Y
\def\banZ{\mathit{Z}}              % Banach space Z
\def\dbanX{\mathit{X^*}}           % dual space of Banach space X
\def\dbanY{\mathit{Y^*}}           % dual space of Banach space Y
\def\dbanZ{\mathit{Z^*}}           % dual space of Banach space Z

\def\banA{\mathit{A}}              % Banach algebra A
\def\banB{\mathit{B}}              % Banach algebra B
\def\banC{\mathit{C}}              % C*-algebra C


\def\domain#1{\mathfrak{D}(#1)}    % domain
\def\range#1{\mathfrak{R}(#1)}     % range

\def\isoJ{\mathfrak{J}}            % isometry J
\def\repA{{\cal{A}_{\sql}}}        % representation operator A
\def\linS{{S}}                     % linear operator S
\def\linG{{G}}                     % linear operator G
\def\linL{{L}}                     % linear operator L
\def\linT{{T}}                     % linear operator A
\def\linA{\mathit{A}}              % linear operator A
\def\linB{{B}}                     % linear operator B
\def\linM{{M}}                     % linear operator M
\def\linU{{U}}                     % linear operator U
\def\linI{{I}}                     % linear operator I
\def\A{\cal{A}}                    % linear differential operator A
\def\BDO{\cal{B}}                  % linear boundary differential operator B
\def\linU{{U}}                     % unitary operator U
\def\linP{{P}}                     % projection P
\def\compA{k}                      % compact operator k
\def\linFR{f}                      % finite rank operator f
\def\fredA{A}                      % Fredholm operator F

\def\HAM{\cal{H}}                  % Hamilton operator
\def\PA{\partial^{\alpha}}         % elementary differential operator
\def\DO{\cal{A}}                   % linear differential operator
\def\BC{\cal{B}}                   % boundary operator
\def\SLO{\cal{L}_{p,q}}            % Sturm-Liouville operator
\def\MO{\cal{L}_{M}}               % momentum operator

\def\HS#1{S(#1)           }        % algebra of Hilbert-Schmidt operators
\def\FR#1{F(#1)           }        % algebra of finite rank operators
\def\TC#1{N(#1)           }        % algebra of trace class operators
\def\BU#1{\mathscr{U}(#1) }        % algebra of unitary linear operators
\def\BO#1{\mathscr{L}(#1) }        % algebra of bounded linear operators
\def\CO#1{\mathscr{C}(#1) }        % algebra of compact linear operators
\def\CL#1{\mathscr{C}(#1) }        % algebra of compact linear operators
\def\FO#1{\Phi(#1) }               % set of Fredholm operators
\def\FK#1#2{\Phi_{#2}(#1) }        % set of Fredholm operators with index k

\def\hilH{\mathit{H}}              % Hilbert space H
\def\clU{\mathit{U}}               % closed subspace U

\def\orthO{\mathfrak{O}}           % orthonormal system O
\def\orthS{\mathfrak{S}}           % orthonormal system S
\def\basB{\mathfrak{B}}            % basis B of a Hilbert space

\def\GL#1 {\mathbf{GL}(#1)}        % linear group
\def\OG#1 {\mathbf{O}(#1)}         % linear group
\def\SOG#1 {\mathbf{SO}(#1)}       % linear group

%------------------------------------------------------------------------------%
% distribution theory                                                          %
%------------------------------------------------------------------------------%
\def\disF{F}                  % distribution F
\def\disG{G}                  % distribution G
\def\disH{H}                  % distribution H
\def\disU{U}                  % distribution U
\def\disT{T}                  % distribution T
\def\SF{C^{\infty}(\G)}       % space of smooth functions
\def\TF{C_c^{\infty}(\G)}     % space of compactly supported smooth functions
\def\TFRd{C_c^{\infty}(\R^d)} % space of test functions
\def\SRd{\mathscr{S}(\R^d)}   % Schwartz space
\def\SR1{\mathscr{S}(\R)}     % Schwartz space
\def\B1#1{C^1(#1) }           % space of continuously differentiable functions
\def\Bm#1{C^m(#1) }           % space of continuously differentiable functions
\def\Ck{C^k(\G)}
\def\TD{\mathscr{S}'(\R^d)}   % space of temperated distributions
\def\E{\cal{E}}               % space of distributions with compact support
\def\D{\mathscr{D}}           % D
\def\FT{\mathcal{F}}          % Fourier transformation F
\def\FT{\mathscr{F}}          % Fourier transformation F

%------------------------------------------------------------------------------%
% other definitions                                                            %
%------------------------------------------------------------------------------%
\def\Proof{\textsc{Proof}}
\def\FRECHET{Fr\'{e}chet}
\def\ARZELA{Arzel\'{a}}
\def\GARDING{G\r{a}rding}
\def\HOELDER{H\"older}
\def\SCHROEDINGER{Schr\"odinger}
\def\JOERGENS{J\"orgens}
\def\WUEST{W\"uest}
\def\KAEHLER{K\"aehler}
\def\HOEFLINGER{H\"oflinger}
\def\POINCARE{Poincar\'{e}}
\def\FA{Functional Analysis}

\def\Author   {K.\,{\HOEFLINGER}}
\def\AuthorUC {K.\,HOEFLINGER}
\def\Version  {1.0}
\def\Email    {khoeflinger@t-online.de}


%\renewcommand\thesection{\Roman{section}}

\def\Title{Linear Operators and Spectral Theory}

\titleformat{\section}
{\normalfont \bfseries \filcenter}
{\thesection.\;}{0mm}{}

\titleformat{\subsection}[runin]
{\normalfont \filcenter}
{}{0mm}{}

\titleformat{\subsubsection}[runin]
{\itshape \filcenter}
{}{0mm}{}

\titlespacing*{\section}      {5pt}{12pt}{ 4pt}
\titlespacing*{\subsection}   {0pt}{ 6pt}{12pt}
\titlespacing*{\subsubsection}{0pt}{ 6pt}{ 3pt}

\def\SG{$\{S(t)\}_{t \geq 0}$}

%------------------------------------------------------------------------------%
%                                BEGIN OF DOCUMENT                             %
%------------------------------------------------------------------------------%
\begin{document}
\thispagestyle{empty}

%\centerline{\fontfamily{lmr}\selectfont \small{\textbf{\Title}}}
%\centerline{\fontfamily{lmss}\selectfont \small{\textbf{\Title}}}
%\centerline{\fontfamily{bodoni}\selectfont {\textbf{\Title}}}
%\centerline{\fontfamily{lmtt}\selectfont \large{\textbf{\Title}}}
%\centerline{\fontfamily{bch}\selectfont {\textbf{\Title}}}
%\centerline{\fontfamily{pcr}\selectfont \large{\textbf{\Title}}}
\centerline{\large{\textbf{\Title}}}

\vspace{4pt}
\centerline{\small{{\textsc{\Author}}}}
\vspace{10pt}

\fancyhead[C] {\small{\textsc{\Title}}}
\fancyhead[OL]{\small{\thepage}}
\fancyhead[ER]{\small{\thepage}}

\myfootnote
{
  Version {\Version} from \today;
  Email:  \textit{khoeflinger@\,t-online.de}
}

%------------------------------------------------------------------------------%
%                                 INTRODUCTION                                 %
%------------------------------------------------------------------------------%
This elaboration is a review on the theory of linear operators
acting in complex Banach or Hilbert spaces and their spectral properties.
The subject belongs to the mathematical area of \textit{functional analysis}
and has many useful applications,
especially to linear partial differential operators
of elliptic and parabolic type
and their associated boundary and initial value problems.

\subparagraph{}
By use of spectral theory one can study resonance phenomena appearing
in mathematical physics and engineering science.
Here we are led to so-called \textit{eigenvalue problems},
that is,
equations of the form $\A \, u \eq \lambda\,u$
with $\A$ being a linear-elliptic partial differential operator,
where the question arises,
for which complex parameters $\lambda \xeof \C$ this equation
has non-trivial solutions.
But spectral theory also forms the foundation for \textit{semigroup theory},
which is a decicional tool to solve evolution equations
of the form $\dot{u} + \A \, u \eq f$ and $\ddot{u} + \A \, u \eq f$,
respectively.

\subparagraph{} %- honoration to Prof. Dr. Werner
I wrote this treatise in honour of Prof.\;Dr.\;P.\;Werner (1932-2018),
who conveyed my fascination for these beautiful mathematical topics
by his lessons hold in 1984 and 1985
at Mathematical Institute A of University Stuttgart.

%------------------------------------------------------------------------------%
\section{Banach and Hilbert Spaces}                                            %
%------------------------------------------------------------------------------%
%-- 
%-- 1. BANACH AND HILBERT SPACES
%--    -------------------------
%-- normed linear spaces
Throughout this text let $\F$ be one of the fields $\R$ and $\C$.
By a \textit{normed space} $\banX$ we mean a $\F$-linear space
that differs from other abstract linear spaces
by the presence of a particular mapping $\norm{\,\cdot\,}\: \banX \xto \R$,
denoted the \textit{norm} on $\banX$,
which fulfills for all $x,y \xeof \banX$ and $\lambda \xeof \F$

\begin{itemize}[leftmargin=40pt,itemsep=2pt]
\item[(N-1)]
\,$\norm{x} \, \geq \, 0$,

\item[(N-2)]
\,$\norm{\lambda \, x} \= \abs{\lambda} \cdot \norm{x}$,

\item[(N-3)]
\,$\norm{x+y} \, \leq \, \norm{x} + \norm{y}$
\; \textit{(triangle inequality)},

\item[(N-4)]
\,$\norm{x} \= 0$ \iff $x \eq \0$.
\end{itemize}

\paragraph{} %- norm topology
%------------------------------------------------------------------------------%
Linear spaces may of course be furnished with different norms.
We shall call arbitrary norms
$\norm{\,\cdot\,}_{\alpha}$ and $\norm{\,\cdot\,}_{\beta}$
on a linear space $\banX$ to be \textit{equivalent} to each another,
if there are constants $c_1,c_2 > 0$ such that
$
c_1 \norm{x}_{\alpha} \leq
    \norm{x}_{\beta}  \leq 
c_2 \norm{x}_{\alpha}$
for all $x \xeof \banX$.

\subparagraph{}
Each normed linear space $\banX$ can be furnished
with a very natural topology $\topT$,
called the \textit{norm topology} on $\banX$.
Its definition is given either by the assertion
that $\topT := \topT_{\norm{\cdot}}$ is the smallest topology on $\banX$
such that the norm mapping remains continuous,
or by the metric topology $\topT_d$,
when $\banX$ is considered as a metric space with metric function
$d(x,y) := \norm{x-y} \fa x,y \xeof \banX$.
In fact,
both definitions lead to the same collection of so-called
"open" subsets in $\banX$,
so we have $\topT := \topT_{\norm{\cdot}} \eq \topT_d$,
which justifies to speak of \textit{the} topology of a normed linear space.
Hence a multiplicity of topological concepts,
like the \textit{closure} of a set,
the \textit{neighbourhood} of a point,
the \textit{convergence} of a series and many others are meaningful
in the setting of normed linear spaces.
Notice
that equivalent norms on $\banX$ are creating one and the same topology.
For example,
if $\banX$ is a finite-dimensional linear space,
then all norms are equivalent
and consequently there is exactly one topology on $\banX$ induced by them.

\subsection{} %-- sequentially complete normed spaces, Banach spaces
Recall from metric space theory
that a sequence $(x_n)$ in a normed linear space $\banX$
is called \textit{Cauchy},
if to each $\epsilon > 0$ there is $N_{\epsilon} \xeof \N$
such that $m,n > N_{\epsilon}$ implies $\norm{x_m{-}x_n} < \epsilon$.
The normed space $\banX$ is called \textit{sequentially complete}
(or shortly \textit{complete}),
if for any Cauchy sequence $(x_n)$ %in $\banX$
there is $x \xeof \banX$
such that $\lim_{\,n \xto \infty} \norm{x_n{-}x} \eq 0$.
Equivalently,
$\banX$ is complete
iff
for every absolutely convergent series
$\sum_{n=1}^{\infty} \norm{x_n} < \infty$
there is $x \xeof \banX$ such that $\norm{x{-}x_n} \xto 0$ for $n \xto \infty$.
A complete normed linear space is called a \textit{Banach space}.

\subparagraph{}
It should be mentioned that each uncomplete normed linear space can be regarded
as a dense linear subspace of a certain Banach space.
The general procedure readily follows the way
how to construct the set of real numbers out of that of the rational ones.

\paragraph{} %- embeddings
%------------------------------------------------------------------------------%
We go ahead with the notion of \textit{embedding},
which is a special kind of mappings between normed linear spaces
and leads to the concept of {norm-isomorphism}.
Preliminary,
a mapping $f \: \banX \xto \banY$ is called continuous
if for every open set $\openO \xsubset \banY$
the pre-image $f^{-1}(\openO)$ is open in $\banX$.
Then by an embedding $\iota$ is meant
a continuous and injective mapping from a normed linear space $\banX$
into a further normed linear space $\banY$,
which is \textit{linear},
that is,
\[
\iota(x + y) \= \iota x + \iota y \;\and\; \iota(\lambda x) = \lambda \iota x
\]
holds for $x,y \xeof \banX$ and $\lambda \xeof \F$.
If $\iota$ is even surjective,
it is said to be a linear isomorphism between the spaces $\banX$ and\,$\banY$.
Then \textsc{Banach}\textit{'s isomorphism theorem} claims that % in this case
the inverse mapping $\iota^{-1}$ is also an embedding,
so $\iota$ and $\iota^{-1}$ both constitute linear homeomorphisms
between $\banX$ and $\banY$.
Furthermore,
the embedding $\iota$ is said to be a \textit{norm-isomorphism},
if it is surjective and $\norm{\iota x} \eq \norm{x}$ holds
for all $x \xeof \banX$ (notation:\,$\banX \cong \banY$).
It should be clear that
in this case the inverse mapping $\iota^{-1}$ is a norm-isomorphism, too.

\subsection{} %-- inner products and Hilbert spaces
%------------------------------------------------------------------------------%
A complex (real) \textit{Hilbert space} differs
from an arbitrary Banach space by the presence of an \textit{inner product},
which is nothing but a symmetric and positive sesquilinear (bilinear) form
on such a space.
This form defines the norm and the metric topology of the Hilbert space,
but also determines the inner geometry to be close
to that of an Euclidean space.

\subparagraph{} %- inner products
We want to make these assertions precise.
Given a real or complex linear space $\hilH$,
an \textit{inner product} in $\hilH$ is defined to be a mapping
$\scalar{\,\cdot\,}{\,\cdot\,} \: \hilH \xtimes \hilH \xto \F$,
which is

\begin{itemize}[leftmargin=35pt]
\item[(I-1)]
linear in the first argument,

\item[(I-2)]
either linear or antilinear in the second argument (conditional upon,
whether $\F$ is the real or the complex number field),
\end{itemize}
together with the additional properties

\begin{itemize}[leftmargin=35pt]
\item[(I-3)]
$\scalar{y}{x} \eq           \scalar{x}{y}$  ($\F=\R$) rsp.
$\scalar{y}{x} \eq \overline{\scalar{x}{y}}$ ($\F=\C$),
\item[(I-4)]
$\scalar{x}{x} \geq 0$ for all $x \neq 0$,
\item[(I-5)]
$\scalar{x}{x} \eq 0$ implies $x \eq 0$.
\end{itemize}

\subparagraph{}
Property (I-3) is denoted the \textit{Hermitean symmetry}
of $\scalar{\,\cdot\,}{\,\cdot\,}$,
while (I-4) and (I-5) together mean \textit{positive definiteness}.
The presence of the inner product gives rise to define a norm in $\hilH$
according to
\begin{equation}\label{DERIVED_NORM}
\norm{x} \eqdef \scalar{x}{x} ^{1/2} \quad \quad (x \xeof \hilH),
\end{equation}
so each linear space endowed with an inner product is normed
in a natural manner.
The inner product and the norm are coupled by
the \textit{Cauchy-Schwarz} \textit{inequality}
\[
\abs{\scalar{x}{y}} \, \leq \, \norm{x} \cdot \norm{y}
\quad \quad (x,y \xeof \hilH),
\]
which implies that the mapping
$\scalar{\,\cdot\,}{\,\cdot\,}\:\hilH \xtimes \hilH \xto \F$ is continuous.
An inner product space $\hilH$ is called a \textit{Hilbert space}
or shortly \textit{Hilbertian},
if it is complete {\wrt} the norm defined by \eqref{DERIVED_NORM}.

\subparagraph{} %- parallelogram equation and polarization formula
The pleasant geometric properties of inner product spaces
may be attributed to the validness of the \textit{parallelogram equation}
\begin{equation}\label{PARALLELOGRAM_EQUATION}
2 \norm{x}^2 + 2 \norm{y}^2 \= \norm{x+y}^2 + \norm{x-y}^2
\quad \quad (x,y \xeof \hilH).
\end{equation}
In general,
the norm $\norm{\,\cdot\,}$ on a linear space $\banX$
can be induced by an inner product
\iff \eqref{PARALLELOGRAM_EQUATION} holds.
In this case and if $\banX$ is a complex linear space,
the inner product can be recovered from the norm
by the \textit{polarization formula}
\[
\scalar{x}{y}
\=
\frac{1}{4} \,
(\,
\norm{x + y}^2 - \norm{x-y}^2 +i \norm{x+iy}^2 - i \norm{x-iy}^2
\,).
\]
If $\banX$ is a real space,
however,
the polarization formula simplifies to
\[
\scalar{x}{y} \=
\frac{1}{4} \, (\, \norm{x + y}^2 - \norm{x-y}^2 \,).
\]

\subsection{} %-- the inner geometry of a Hilbert space
%------------------------------------------------------------------------------%
In the following we want to point out
that Hilbert spaces differ from general Banach spaces
by a couple of desirable properties:

\begin{itemize}[leftmargin=4mm,itemsep=6pt]
\item[1.] \textit{Orthogonality}.
Two elements $x,y \xeof \hilH$ are called \textit{orthogonal} to each another,
if $\scalar{x}{y} \eq 0$ (notation: $x {\bot} y$).
Any subset $\orthO \xsubset \hilH$ is said to be \textit{orthogonal}
or an \textit{orthogonal system},
if all elements in $\orthO$ are orthogonal to each other.
Finally,
the \textit{orthogonal space} $\setA^{\bot}$
of a subset $\setA \xsubset \hilH$
contains exactly the elements $y \xeof \hilH$,
which are orthogonal to all $x \xeof \setA$,
and forms a closed linear subspace of $\hilH$.

\item[2.] \textit{Projection Theorem}.
Let $\setA \xsubset \hilH$ be a \textit{closed} linear subspace.
Then for each $x \xeof \hilH$ there is exactly one $x_0 \xeof \setA$
such that
\[
\norm{ x - x_0 } \= \inf \, \{ \, \norm {x-y} \: y \in \setA \,\}.
\]
The element $x_0$ is called the \textit{projection} of $x$ onto $\setA$
and minimizes the distance between $\setA$ and $x$.
The associated linear mapping
$\pi_{\setA} \: x \xeof \hilH \mapsto x_0 \xeof \setA$
is also denoted the \textit{projection} of $\hilH$ onto $\setA$.

\item[3.] \textit{Decomposition Theorem}.
Using the notation of orthogonal space $\setA^{\bot}$,
the projection of an element $x \xeof \hilH$ onto $\setA$
is characterized as follows:
$x_0 \eq \pi_{\setA} x$ iff $x {-} x_0 \xeof \setA^{\bot}$.
Hence each $x \xeof \hilH$ can uniquely be represented by the formula
$x \eq x_0 {+} x^{\bot}$,
where $x_0 \xeof \setA$ and $x^{\bot} \xeof \setA^{\bot}$,
so the entire space $\hilH$ is norm-isomorphic
to the direct sum $\setA \, \oplus \, \setA^{\bot}$
of the closed linear subspaces $\setA$ and $\setA^{\bot}$.
This fact is known as the \textit{decomposition theorem} in Hilbert space.

\subparagraph{}
A closed linear subspace $\setA$ of a Banach space $\banX$
is said to be \textit{topologically complementable},
if there is a closed subspace $\setB \xsubset \banX$
such that $\setA \cap \setB \eq \{\0\}$,
          $[\setA \cup \setB] \eq \banX$ and $\banX \cong \setA \oplus \setB$,
that is,
the space $\banX$ is linear-homeomorphic
to the product space $\setA \xtimes \setB$.
Using this notation, the decomposition theorem quotes that
in the Hilbert space setting every closed subspace $\setA$ is topologically
complemented by the orthogonal space $\setA^{\bot}$.

\item[4.] \textit{Existence of Orthonormal Bases}.
An orthogonal system $\orthO \xsubset \hilH$ is denoted
an \textit{orthonormal system},
if $\scalar{x}{x} \eq 1$ for all $x \xeof \orthO$.
However,
if $\orthO$ contains only \textit{finitely} many elements,
then the span $[\orthO]$ forms a \textit{closed} subspace in $\hilH$,
and the projection $\pi_{[\orthO]}$ of any element $x$ onto $[\orthO]$
can be expressed by the formula
\[
\pi_{[\orthO]} \, x \= \sum_{i=1}^n \, \scalar{x}{e_i} \, e_i.
\]
But if $\orthO$ consists of arbitrarily many elements,
then the inequality $\scalar{x}{e_{\alpha}} \neq 0$ can only be valid
for mostly countable many elements $e_{\alpha} \xeof \orthO$,
which in conclusion enables to add the coefficients $\scalar{x}{e_{\alpha}}$
together about all elements of $\orthO$.

\subparagraph{} %- orthonormal bases
An orthonormal system $\orthO$ is called
an \textit{orthonormal base},
if its span $[ \orthO ]$ lies dense in $\hilH$.
In this case it is possible to represent each element $x \xeof \hilH$
by the formal sum
\begin{equation}\label{FOURIER_COEFFICIENTS}
x \= \sum_{e_{\alpha} \in \, \orthO} \, \scalar{x}{e_{\alpha}} \, e_{\alpha}.
\end{equation}
The inner products $\scalar{x}{e_{\alpha}}$ are called
the \textit{Fourier coefficients} of $x$ {\wrt} $\orthO$.
Notice that the right hand side of \eqref{FOURIER_COEFFICIENTS}
consists of \textit{mostly countable many} non-zero summands,
and this series proves to be invariant
by realignments of the elements $e_{\alpha} \xeof \orthO$.

\subparagraph{}
There is of course the need to determine,
whether a given orthonormal system $\orthO$ forms an orthonormal base.
In fact,
it can be shown that
the following statements are equivalent:

\begin{itemize}[leftmargin=9mm]
\item[(a)]
$\orthO$ is an orthonormal base in $\hilH$.

\item[(b)]
$\orthO$ is \textit{complete},
that is,
$\scalar{x}{e} \eq 0$ for all $e \xeof \orthO$ implies $x \eq \0$.

\item[(c)]
$\sum_{\, e_{\alpha} \xeof \orthO} \,
\abs{ \scalar{x}{e_{\alpha}} } ^2 \eq \norm{x}^2$ for all $x \xeof \hilH$
(\textit{Parseval}\textit{'s equation}).
\end{itemize}

\paragraph{} %- existence of orthonormal bases
%------------------------------------------------------------------------------%
It is also nearby to ask for the existence of orthonormal bases.
First of all,
the collection of all orthonormal systems in a Hilbert space $\hilH$
forms a partial ordered set,
and under the assumption that \textsc{Zorn}{'s lemma}\,\footnote{\,
\,The \textit{axion of choice} claims that
if $\cal{F}$ is a disjoint family $F_{\alpha}$ ($\alpha \xeof \J$)
of non-empty sets,
then there is a set $A$ that shares exactly one element
with each other set $F_{\alpha} \xeof \cal{F}$.}
is valid,
it contains a \textit{maximal} element $\orthO$
that cannot be increased,
hence each Hilbert space $\hilH \neq \{\0\}$ owns an orthonormal base.

\subparagraph{} %- Hilbert space dimension
Furthermore,
if $\orthO_1$ and $\orthO_2$ are orthonormal bases in $\hilH$,
then there is a bijective mapping between them,
thus they own identical cardinal numbers.
This legitimates us to introduce the notation
of \textit{Hilbert space dimension} $\abs{\hilH} $
being the cardinal number of any orthonormal base.

\subparagraph{} %- separable Hilbert spaces
A Hilbert space $\hilH$ is called \textit{separable}
if there is a finite or mostly countable subset lying dense in $\hilH$.
It turns out that $\hilH$ is separable
\iff
$\abs{ \hilH} \leq \abs{\N} $.
Indeed,
in this case
every orthonormal base in $\hilH$ is finite or mostly countable,
and the space $\hilH$ itself proves to be norm-isomorphic
to the Hilbert space $\ell^2$ of square-summarizable sequences
$f \: \N \xto \F$.
\end{itemize}

%------------------------------------------------------------------------------%
\section{Basic Theory of Linear Operators}                                     %
%------------------------------------------------------------------------------%
%-- 
%-- 2. BASIC THEORY OF LINEAR OPERATORS
%--    --------------------------------
%--    generalities on linear operators
In this section we want to sketch the fundamentals of operator theory.
If we assume $\banX,\,\banY$ to be normed linear spaces,
then by a \textit{linear operator} between $\banX$ and $\banY$
is meant a mapping
$\linA \: \domain{\linA} \xsubset \banX \xto \banY$,
where the \textit{domain of definition} $\domain{\linA}$
forms a linear subspace of $\banX$ and
$\linA$ is \textit{linear},
that is,
\[
\linA(\lambda x + \mu y) \= \lambda \, \linA x + \mu \, \linA y
\]
holds for all $x,y \xeof \domain{\linA}$ and $\lambda,\mu \xeof \F$.
The space $\banX$ is denoted the \textit{domain space} of $\linA$,
whereas $\banY$ is said to be the \textit{range space} of $\linA$.

\subparagraph{} %- choices of the domain space
The domain of definition $\domain{\linA}$ may of course
coincide with the entire space $\banX$,
but it is often sufficient to require that 
$\domain{\linA}$ lies dense in $\banX$,
so each $x \xeof \banX$ can be approximated
by a sequence of elements in $\domain{\linA}$.
In this case it is usual to say
that $\linA$ is \textit{densely defined} in $\banX$.
However,
if the set $\domain{\linA}$ is \textit{not} dense in $\banX$,
it is nearby to consider the closure $\banX_0$ of $\domain{\linA}$ in $\banX$
and to pay attention to the slightly modified but densely defined operator
$\linA_0\:\domain{\linA} \xsubset \banX_0 \xto \banY$.

\paragraph{} %- sums and compositions of linear operators
%------------------------------------------------------------------------------%
Rather the domains of definition for operators $\linA_1$ and $\linA_2$
with same domain space $\banX$ and range space\,$\banY$
may be proper linear sub\-spaces of $\banX$,
the \textit{sum} $\linA_1 + \linA_2$ and
the \textit{composition} $\linA_1 \circ \linA_2$ are well-defined,
if their domains of definition are given by
\[
\begin{cases}[1.2]
\quad 
\domain{\linA_1 + \linA_2} \eqdef
\domain{\linA_1} \, \cap \, \domain{\linA_2},\\
\quad
\domain{\linA_1 \circ \linA_2} \eqdef
\{\, x \in \domain{\linA_2} \: \linA_2 \, x \in \domain{\linA_1}\,\}.\\
\end{cases}
\]
Furthermore,
for arbitrary $\lambda \xeof \F$
the operator $\lambda \linA$ is well-defined by
$\domain{\lambda \linA} := \domain{\linA}$ and
$(\lambda \linA) x := \lambda (\linA x)$
for $x \xeof \domain{\lambda \linA}$
and is called a \textit{dilatation} in $\banX$.

\subparagraph{} %- extensions of linear operators
If the domain of definition is a proper subset of the domain space,
the concept of \textit{extension} is meaningful.
Let $\linA_1$ and $\linA_2$ be linear operators
with domain of definitions $\domain{\linA_1} \xsubset \banX$
and $\domain{\linA_2} \xsubset \banX$,
respectively.
We introduce the partial ordered relation
\[
\linA_1 \prec \linA_2
: \Longleftrightarrow 
\domain{\linA_1} \xsubset \domain{\linA_2}
\, \and \,
\linA_1 \, x = \linA_2 \, x
\fa x \in \domain{\linA_1}
\]
and denote $\linA_2$ an \textit{extension} of $\linA_1$
if $\linA_1 \prec \linA_2$.
The operators $\linA_1$ and $\linA_2$ are said to be \textit{equal}
(briefly $\linA_1 \eq \linA_2$)
iff $\linA_1 \prec \linA_2$ and $\linA_2 \prec \linA_1$,
or equivalently,
if $\domain{\linA_1} \eq \domain{\linA_2}$ and
$\linA_1\,x \eq \linA_2\,x$ for all $x \xeof \domain{\linA_1}$.

\subparagraph{} %- kernels, ranges and cokernels
Finally,
to every linear operator $\linA \: \domain{\linA} \xsubset \banX \xto \banY$
there are basically assigned three linear subspaces,
namely

\begin{itemize}[leftmargin=4mm]
\item[1.]
the \textit{kernel} $\kernel{\linA} :=
\{\,x \xeof \domain{\linA} \: \linA x \eq \0\,\} \xsubset \banX$;

\item[2.]
the \textit{range} $\range{\linA} :=
\{\, \linA x: x \xeof \domain{\linA}\,\} \xsubset \banY$;

\item[3.]
the \textit{cokernel} $\cokernel{\linA} \xsubset \banY$,
that is,
the linear-algebraic complement of $\range{\linA}$ within the space $\banY$,
hence
$\range{\linA} \oplus \cokernel{\linA} \eq \banY$.
\end{itemize}

\subsection{} %-- classification of linear operators by topological properties
%------------------------------------------------------------------------------%
A first classification of linear operators
between a domain space $\banX$ and a range space\,$\banY$
can be performed by checking certain properties,
which are defined by means of the norm topologies in these spaces.
It leads in conclusion to the concepts of \textit{continuous}, \textit{compact},
\textit{closed} and \textit{closable} operator.

\begin{itemize}[leftmargin=12pt,itemsep=3pt]
\item[1.] %--  - continuous and bounded operators
%------------------------------------------------------------------------------%
Remember that a linear operator
$\linA \: \domain{\linA} \xsubset \banX \xto \banY$
is \textit{continuous},
if the pre-image $\linA^{-1}(\openO)$ of each open set $\openO \xsubset \banY$
is open in the domain space $\banX$.
However,
in the setting of normed linear spaces
we may replace continuity by \textit{sequential continuity},
which means that $\linA$ is continuous in $x \xeof \banX$
if for each convergent series $(x_n)$ with limit element $x \xeof \banX$
we have $\lim_{n \to \infty} \linA x_n  \eq \linA x$.

\subparagraph{} %- bounded linear operators and operator norm
Moreover,
a linear operator $\linA$ between normed linear spaces $\banX$
and\,$\banY$ is denoted \textit{norm-bounded} or simply \textit{bounded},
if there is $c \geq 0$
such that $\norm{\linA x} \leq c \cdot \norm{x}$
for each $x \xeof \domain{\linA}$.
It is easy to check that $\linA$ is continuous iff $\linA$ is bounded,
hence these concepts are equivalent to each other.
In this case the \textit{operator norm} of $\linA$ is well-defined according to
\begin{equation}\label{OPERATOR_NORM}
\norm{\linA} \eqdef \sup_{x \in \domain{\linA}, \norm{x}=1} \norm{\linA x}.
\end{equation}

\paragraph{} %- the B.L.T. theorem
%------------------------------------------------------------------------------%
If a linear operator $\linA$ is continuous and 
$\domain{\linA}$ is a dense subset of the domain space $\banX$,
then $\linA$ can uniquely be extended to a linear mapping
defined on the closure $\overline{\domain{\linA}} \eq \banX$,
which proves to be continuous as well.
This is the assertion of the \textit{B.L.T.\,theorem}
(written out the \textit{bounded linear transformation theorem},
see \cite{RS1}, page 6).
It is founded by the fact that
any continuous linear operator $\linA$ is even \textit{uniformly} continuous,
that is,
for each $\epsilon > 0$ there exists $\delta > 0$,
such that
\[
\norm{x-y}_{\banX} < \delta
\; \mbox{implies} \;
\norm{\linA x - \linA y}_{\banY} < \epsilon
\]
for $x,y \xeof \domain{\linA}$,
hence $\delta$ can be chosen independently of $x,y$.
Thus we may basically assume
that if $\linA$ is continuous,
then $\domain{\linA} \eq \banX$ by applying this canonical extension,
and the operator $\linA$ finally appears as a mapping
$\linA \: \banX \xto \banY$.

\subparagraph{} %- the linear space B(X,Y)
Let $\BO{\banX,\banY}$ be the set
of \textit{bounded} linear operators $\linA \: \banX \xto \banY$.
Endowed with addition $\linA{+}\linB$
and scalar multiplication $\lambda \linA$,
the set $\BO{\banX,\banY}$ forms a linear space,
and furnished with the norm given by \eqref{OPERATOR_NORM}
it is even a normed linear space.
This space is complete,
if $\banY$ is so,
and in this case $\BO{\banX,\banY}$ forms a Banach space in its own right.

\subparagraph{} %- the operator algebra B(X)
If $\banY \eq \banX$,
the set $\BO{\banX} := \BO{\banX,\banX}$ constitutes a \textit{normed algebra}
with composition $(\linA_1,\linA_2) \mapsto \linA_2 \circ \linA_1$
as its multiplicative operation
and owns the property
\[
\norm{\linA_1 \circ \linA_2} \, \leq \, \norm{\linA_1} \cdot \norm{\linA_2}.
\]
It is called the \textit{operator algebra} associated with $\banX$
and actually forms a Banach algebra if\,$\banY$ is complete.

\paragraph{} %- main theorem on continuous linear operators
%------------------------------------------------------------------------------%
Let us collect the obtained results 
to the following main theorem on continuous linear operators,
where $\banX$ and\,$\banY$ are assumed to be normed linear spaces:

\begin{theorem}\label{CONTINUITY_THEOREM}
For $\linA \: \banX \xto \banY$ being a linear operator,
the subsequent assertions are equivalent to each other:

\begin{itemize}[leftmargin=9mm]
\item[(a)] $\linA$ is continuous;
\item[(b)] $\linA$ is uniformly continuous;
\item[(c)] $\linA$ is sequentially continuous in every $x \xeof \banX$;
\item[(d)] $\linA$ is sequentially continuous in the origin $\0 \xeof \banX$;
\item[(e)] $\linA$ is bounded.
\end{itemize}
\end{theorem}

\item[2.] %--   - compact operators
%------------------------------------------------------------------------------%
A linear operator $\compA \: \banX \xto \banY$ between Banach spaces
$\banX,\,\banY$ is called \textit{compact},
if the image of each bounded set under $\compA$ is relative compact,
or equivalently,
if $\{x_n\}_{n \in \N}$ is a bounded sequence in $\banX$,
then the sequence $\{\compA \, x_n\}_{n \in \N}$ in $\banY$
owns a convergent subsequence.

\subparagraph{}
For example,
a linear operator $\compA$ is compact,
if at least one of the spaces $\banX$ and $\banY$ is finite-dimensional.
In case of $\banY \eq \banX$,
the identity operator $\1_{\banX} \: \banX \xto \banX$ is compact
iff $\banX$ is finite-dimensional.

\paragraph{} %- properties of compact operators
%------------------------------------------------------------------------------%
Compact operator own a couple of useful properties:

\begin{itemize}[leftmargin=9mm]
\item[(a)]
every compact linear operator is continuous;

\item[(b)]
the composition $\linA_1 \circ \linA_2$ of bounded linear operators
is compact,
if one of them is so;

\item[(c)]
the first and second property together imply that
the set $\CO{\banX}$ of compact linear operators $\compA\:\banX \xto \banX$
is a \textit{closed} ideal within the operator algebra $\BO{\banX} $,
hence it is complete and forms a Banach algebra in its own right;

\item[(d)]
finite linear combinations of compact linear operators are compact, too;

\item[(e)]
the limit of a norm convergent sequence of compact operators is compact.
\end{itemize}

\paragraph{} %- finite-rank operators
%------------------------------------------------------------------------------%
A linear operator $\linFR \: \banX \xto \banY$ is said
to have \textit{finite rank},
if the range $\range{\linFR}$ is finite-dimensional.
It is obvious that $\linFR$ is compact.

\subparagraph{}
Suppose that $(\linFR_n)$ is a sequence of finite-rank operators
converging to the operator $\linA$ {\wrt} the norm of $\BO{\banX,\banY}$.
Then $\linA$ is compact
and the set $\CO{\banX,\banY}$ of compact linear operators
proves to be a subset of the closure of the set $\FR{\banX,\banY}$
of finite-rank operators,
{\ie}
$\CO{\banX,\banY} \xsubset \overline{\FR{\banX,\banY}}$.

\subparagraph{}
Reversely,
in the setting of Hilbert spaces 
each compact linear operator $\compA$ appears in this way,
that is,
$\compA$ is the norm-limit of a sequence of finite-rank operators
(but this is not necessarily true in the general setting of Banach spaces).

\paragraph{} %- Hilbert Schmidt operators
%------------------------------------------------------------------------------%
Let $\basB \eq \{e_1,e_2,\ldots\}$ be any orthonormal base
of a separable Hilbert space $\hilH$ and let $\linA \xeof \BO{\hilH}$.
We define the number $S(\linA)$ according to
\[
S(\linA) \eqdef \sum_{i \in \basB} \, \norm{ \linA \, e_i}^2.
\]
Then it turns out this definition for $S(\linA)$
does not depend on the choice of the orthonormal base.
In case of $S(\linA) < \infty$
the operator $\linA$ is called a \textsc{Hilbert-Schmidt} \textit{operator},
and the number
\[
\norm{\linA}_{HS} \eqdef \sqrt{S(\linA)}
\]
is denoted the \textit{Hilbert-Schmidt} \textit{norm} of $\linA$,
since the mapping $\linA \mapsto \norm{\linA}_{HS} $
owns the properties of a norm.

\subparagraph{}
The set $\HS{\hilH}$ of {Hilbert-Schmidt} operators in $\hilH$
is topologically and algebraically closed in $\BO{\hilH}$
{\wrt} the operations $\linA_1 + \linA_2$ and $\lambda \linA$,
hence $\HS{\hilH}$ endowed with norm
$\norm{\,\cdot\,}_{HS}$ forms a Banach space in its own right.
Each {Hilbert-Schmidt} operator is compact,
so $\HS{\hilH}$ forms a two-sided ideal within the space $\CO{\hilH}$
of compact linear operators in $\hilH$.

\item[3.] %--   - graph-closed operators
%------------------------------------------------------------------------------%
A linear operator $\linA \: \domain{\linA} \xsubset \banX \xto \banY$
is called \textit{graph-closed} (or shortly \textit{closed}) if the set
\[
\Gamma_{\linA} \eqdef
\{(x,\linA x) \: x \xeof \domain{\linA}\} \, \subset \, \banX \xtimes \banY,
\]
denoted the \textit{graph} of $\linA$,
is closed in the space $\banX \xtimes \banY$
{\wrt} the product topology $\topT_{\,\banX \xtimes \banY}$.
This property especially implies that
the kernel $\kernel{\linA}$ of $\linA$ forms a closed set in $\banX$.

\subparagraph{} %- sequentially closed operators
It is a matter of fact
that an operator $\linA$ is closed \iff it is \textit{sequentially closed},
that is,
$x_n \xto x$ and $\linA x_n \xto y$ implies
$x \xeof \domain{\linA}$ and $\linA x \eq y$,
briefly written
\[
\begin{rcases}
 x_n       \xto x \quad \\
 \linA x_n \xto y \quad \\
\end{rcases}
\quad \Longrightarrow \quad
x \xeof \domain{\linA} \, \and \, \linA x \eq y.
\]
This can easily proved by the fact
that a subset of a normed linear space is closed
\iff it is sequentially closed.
Checking the sequentially closeness is the usual method
to verify that a given linear operator is graph-closed.
But notice that the equivalence of these both notations is in general not valid,
if $\banX$ and $\banY$ are non-metrizable topological linear spaces.

\subparagraph{} %- graph norm
It is often helpful to make the domain of definition $\domain{\linA}$
of a linear operator $\linA$ into a normed linear space by its own.
This can be done by introducing the so-called \textit{graph norm}
\[
\norm{x}_{\linA} \eqdef \norm{x}_{\banX} \+ \norm{\linA x}_{\banY}
\]
for $x \xeof \domain{\linA}$.
In fact,
the pair $(\domain{\linA},\norm{\,\cdot\,}_{\linA})$
forms a normed linear space,
which constitutes a Banach space \iff 
$\linA$ is sequentially closed.
This equivalence provides a further method
to verify closeness of a given linear operator.

\paragraph{} %- closeness of operators based on Schauder estimates
%------------------------------------------------------------------------------%
The following theorem
uses the possible validness of an abstract \textit{Schauder} \textit{estimate}
to ensure sequentially closenss:

\begin{theorem}\label{BROWDER_THEOREM}
Let $\linA_0 \: \banX \xto \banY$ be a continuous linear operator,
where the Banach space $\banX$ is continuously and densely embedded
into the Banach space $\banY$ by a mapping $\iota \: \banX \xto \banY$.
If the estimate
\begin{equation}\label{SCHAUDER_ESTIMATE}
\norm{x}_{\banX} \, \leq \,
c \cdot
\{ \,\norm{\iota x}_{\banY} + \norm{\linA_0 \, x}_{\banY} \, \}
\end{equation}
is fulfilled for some $c {\,>\,} 0$ and all $x \xeof \banX$,
then the linear operator $\linA$ acting in $\banY$ with
$\domain{\linA} := \iota(\banX)$ and
$\linA\, \iota x := \linA_0 \, x$ for $x \xeof \banX$ is closed.
\end{theorem}

\subparagraph{} %- proof
In fact,
let $(x_n)_{n \in \N}$ be a sequence in $\domain{\linA}$
with $\norm{x_n-x}_{\banX} \xto 0$
and $\norm{\linA \,x_n-y}_{\banY} \xto 0$.
It is to verify that $x \xeof \domain{\linA}$ and $\linA x \eq y$.
Inequality \eqref{SCHAUDER_ESTIMATE} implies that
$(x_n)_{n \in \N}$ is a Cauchy sequence in $\banX$
with limit element $x \xeof \banX$.
Since $\linA_0$ is continuous,
the element $y$ is also the limit
of the sequence $(\linA \, \iota x_n)_{n \in \N}$,
hence $\linA$ is closed.

\item[4.] %- pre-closed operators
%------------------------------------------------------------------------------%
Many linear operators acting between Banach or Hilbert spaces are not closed,
but can be extended to operators being closed.
We shall call such operators \textit{closable} or \textit{pre-closed}.
If a linear operator $\linA$ is closable,
then there is a \textit{minimal closed extension} $\linA_c$ of $\linA$,
called the \textit{closure} of $\linA$.
This mapping $\linA_c$ is given by the fact
that each other closed extension of $\linA$ extends $\linA_c$,
and $\linA_c$ is nothing but the intersection
of all closed extensions of $\linA$.
Clearly,
each closed operator is closable,
since it constitutes a closed extension of itself.
Furthermore,
given a closed linear operator $\linA$
with domain of definition $\domain{\linA}$,
any linear subspace $\cal{L} \xsubset \domain{\linA}$
is called a \textit{core} of $\linA$,
if the closure of $\linA$ restricted to $\cal{L}$ resembles $\linA$.

\paragraph{} %- criterion for closability
%------------------------------------------------------------------------------%
There is a well-known criterion for closability:

\begin{theorem}
A linear operator $\linA$ with domain of definition $\domain{\linA}$
is closable,
if each sequence $(x_n,\linA x_n)_{n \in \N}$
with $x_n \xeof \domain{\linA}$ for $n \xeof \N$ and converging to $(\0,y)$
implies $y \eq \0$,
briefly written
\[
\begin{rcases}
 x_n       \xto \0 \quad \\
 \linA x_n \xto  y \quad \\
\end{rcases}
\quad \Longrightarrow \quad
y \= \0.
\]
\end{theorem}

\begin{comment}
\subparagraph{}
\Proof:
If $\linA$ is closable,
then $\overline{\Gamma_{\linA}}$ is the graph of the closure $\overline{\linA}$
and forms a closed subspace in $\banX \xtimes \banY$.
Let $x_n \xto 0$ and $\linA x_n \xto y$,
so $(x_n,\linA x_n)$ converges to $(\0,y)$,
which implies $y=\0$.

\subparagraph{}
On the other hand,
let $\linA$ own the following property:
$x_n \xto \0$ with $x_n \xeof \domain{\linA}$
and $\linA x_n \xto y$ implies $y=\0$.
We introduce the subspace $D$ consisting of all
$x \xeof \banX$,
for which there is $x_n \xeof \domain{\linA}$ and $y \xeof \banX$
such that $x_n \xto x$ and $\linA x_n \xto y$.
Define $\overline{\linA} x := y$ for $x \xeof D$,
then $\domain{\linA} \xsubset D$ and $\linA \prec \overline{\linA}$.
It is still to show that

\begin{itemize}[leftmargin=7mm]
\item[(1)] $\overline{\linA}$ is well-defined on $D$ and
\item[(2)] $\overline{\linA}$ is closed, hence $\linA$ is closable.
\end{itemize}

\subparagraph{}
To verify (1),
let $(\tilde{x}_n)$ be a further sequence converging to $x$ with
$\linA \tilde{x}_n \xto \tilde{y}$.
But since $x_n - \tilde{x}_n \xto \0$
and $\linA (x_n - \tilde{x}_n) \xto y-\tilde{y}$,
we have $y \eq \tilde{y}$.
Moreover,
it is easy to verify that $\overline{\linA}$ is linear on $D$.

\subparagraph{}
To verify (2),
let $\{x_n\}_{n \in \N}$ be a sequence in $D$ such that $x_n \xto x$
and $\overline{\linA} \xto y$.
We have to verify that $x \xeof D$ and $\overline{\linA} x \eq y$.
Since for arbitrary $\epsilon > 0$ and $x_n \xeof D$
there is $z_n \xeof \domain{\linA}$ such that
$\norm{z_n - x_n} < \frac{\epsilon}{2}$ and
$\lVert \overline{\linA} x_n - \linA z_n \rVert < \frac{\epsilon}{2}$,
we obtain
\[
\norm{x_n - z_n} \, + \,
\lVert \overline{\linA} x_n - \linA z_n \rVert
\, < \, \epsilon.
\]
Hence $z_n \xto x$ and $\linA z_n \xto y$,
that is,
$x \xeof D$ and $\overline{\linA} x \eq y$.
\qed
\end{comment}
\end{itemize}

\subsection{} %-- comparison of continuous and closed linear operators
%------------------------------------------------------------------------------%
Let us emphasize that in general 
the concepts of closeness and continuity of a linear operator $\linA$
are not comparable.
However,
if $\banX$ and $\banY$ are Banach or Hilbert spaces
and $\domain{\linA}$ is a \textit{closed} set in $\banX$,
then closeness of $\linA$ is an easy consequence of continuity of $\linA$.
This is especially true in the following situation:
if $\linA$ is densely defined and continuous,
then $\linA$ can be extended to a continuous linear operator
$\linA_{ext} \: \banX \xto \banY$,
and since $\banX$ itself is closed,
the extension $\linA_{ext}$ is a closed operator.
Hence continuous linear operators of the form $\linA \: \banX \xto \banY$
fit in a natural way
into the class of densely defined and closed linear operators.

\subparagraph{} %-- closed graph theorem
Reversely,
the \textit{closed graph theorem} claims that if $\domain{\linA}$ is closed,
then closeness implies continuity,
so both concepts are equivalent to each other
under this restrictive condition.

\paragraph{} %- bounded inverse theorem
%------------------------------------------------------------------------------%
The following application of the closed graph theorem is of outstanding meaning:
if a linear operator $\linA$ is closed and bijective,
then its inverse $\linA^{-1}$ is also linear and closed,
and since the domain of definition $\domain{\linA^{-1}}$
is the entire space $\banY$,
the closed graph theorem implies that $\linA^{-1} \: \banY \xto \banX$
is continuous and bounded,
respectively.
This fact is known as the \textit{bounded inverse theorem}
for closed linear operators.
But notice that the quotation is also valid for continuous linear operators
$\linA \: \banX \xto \banY$,
which is an immediate consequence
of \textsc{Banach}'s \textit{open mapping theorem}.

\paragraph{} %- summary
%------------------------------------------------------------------------------%
In conclusion,
if the domain of definition $\domain{\linA}$ of a linear operator $\linA$
is a \textit{closed} subspace of the domain space $\banX$,
then the following chain of implications is valid:
\[
\linA \mbox{ compact }    \Longrightarrow
\linA \mbox{ bounded }    \Longleftrightarrow
\linA \mbox{ continuous } \Longleftrightarrow
\linA \mbox{ closed}. %  \Longrightarrow
%\linA \mbox{ closable}.
\]

%------------------------------------------------------------------------------%
\section{Adjoint Spaces and Operators}                                         %
%------------------------------------------------------------------------------%
%-- 
%-- 3. ADJOINT SPACES AND OPERATORS
%--    ----------------------------
%--    classification of linear operators by choice of the range space
So far we have described and classified linear operators
by various topological properties.
Another approach to setup a classification
in the totality of linear operators between a domain space $\banX$
and a range space $\banY$ is to consider specific choices of the range space.
For example,
one can chose

\begin{itemize}[leftmargin=4mm]
\item[1.]
$\banY \eq \banX$ to obtain operators of the form
$\linA \: \domain{\linA} \xsubset \banX \xto \banX$,
and in this case the operator $\linA$ is said to \textit{act in} $\banX$.
In fact,
if the domain space and the range space are the same,
one can study perturbations of $\linA$
of the form $\linA {+} \lambda \1_{\banX}$
(shortly denoted $\linA {+} \lambda$),
where $\lambda \xeof \F$
and $\1_{\banX}$ denotes the \textit{identity operator} in $\banX$,
and this is performed in what is called \textit{spectral theory}
(see section 4).

\item[2.]
$\domain{\linA} \eq \banX$ and $\banY \eq \F$,
leading to \textit{linear functionals} $\linA \: \banX \xto \F$.
By the \textit{adjoint space} $\banX^*$ of a normed linear space $\banX$
is meant the linear space of continuous linear functionals
$\phi \: \banX \xto \F$,
which itself forms a Banach space with norm
\[
\norm{\phi} \eqdef \sup_{\norm{x}=1} \abs{f(x)} \,.
\]
The \textit{Hahn-Banach theorem} claims
that $\banX^* \neq \{\0\}$,
so there are non-trivial continuous linear functionals on $\banX$.
Any Banach space $\banX$ is called \textit{reflexive},
if $\banX$ is norm-isomorphic to the \textit{bi-adjoint space}
$\banX^{**} := (\banX^*)^*$.

\item[3.]
$\domain{\linA} \eq \banX$ and $\banY \eq \banX^*$,
the \textit{adjoint space} of $\banX$,
so the linear operator adopts the form $\linA \: \banX \xto \banX^*$.

\subparagraph{}
In case of $\banX \eq \hilH$, where $\hilH$ is a complex Hilbert space,
it is more convenient to consider linear operators of the form
$\linA \: \hilH \xto \hilH_a^*$,
where $\hilH_a^*$ is the linear space of continuous and
\textit{conjugate linear} functionals $\phi \: \hilH \xto \C$,
which distinguish by the property
\[
\phi(\lambda \cdot x) \eq \overline{\lambda} \, \phi(x)
\; \fa x \in \hilH \and \lambda \in \C.
\]
The linear space $\hilH_a^*$ is denoted the \textit{conjugate space} of $\hilH$
and can be endowed with the norm
\[
\norm{\phi} \eqdef \sup_{\,\norm{x} \eq 1} \, \abs{\phi(x)} \,.
\]

\subparagraph{} %- sesquilinear forms
For each operator $\linA$ with domain space $\hilH$ and
range space $\hilH_a^*$ we can associate a \textit{sesquilinear form}
$\sql_{\linA} \: \hilH \xtimes \hilH \xto \C$ by
\begin{equation}\label{REPRESENTATION_OPERATOR}
\sql(x,y) \eqdef (\linA x)y
\quad \quad (x,y \in \hilH).
\end{equation}
Reversely,
if a sesquilinear form $\sql$ is given,
we may introduce a linear operator
$\linA_{\sql} \: \hilH \xto \hilH_a^*$
by the same formula \eqref{REPRESENTATION_OPERATOR},
called the \textit{representation operator} of $\sql$.

\paragraph{} %- main theorem on continuity of sesquilinear forms
%------------------------------------------------------------------------------%
Similar to theorem \ref{CONTINUITY_THEOREM},
the following theorem characterizes continuity of sesquilinear forms
by equivalent assertions:

\begin{theorem}
Let $\sql \: \hilH \xtimes \hilH \xto \C$ be a sesquilinear form.
Then:

\begin{itemize}[leftmargin=9mm]
\item[(a)]
$\sql$ is continuous;
\item[(b)]
$\sql$ is sequentially continuous in $(x,y)$ for every $x,y \xeof \hilH$;
\item[(c)]
$\sql$ is sequentially continuous in $(\0,\0) \xeof \hilH \xtimes \hilH$;
\item[(d)]
there is $\Theta \geq 0$ such that
$\abs{\sql(x,y)} \leq \Theta \norm{x} \cdot \norm{y}$ for $x,y \xeof \hilH$;
\item[(e)]
The representation operator $\linA_{\sql}$ is continuous.
\end{itemize}
\end{theorem}
\end{itemize}

\subsection{} %-- Riesz's representation theorem
Let us especially consider the case $\sql(x,y) \= \scalar{x}{y}$.
In this case the representation operator $\Phi$
of the inner product is called the \textit{Riesz map} of $\hilH$.
Since inner products are positive definite by definition,
the mapping $\Phi$ proves to be injective.
But \textsc{Riesz}\textit{'s representation theorem} claims
that $\Phi$ is also surjective,
so there is a canonical correspondence
between the elements of $\hilH$ and $\hilH_a^*$:

\begin{theorem}\label{RIESZ_REPRESENTATION_THEOREM}
If $\hilH$ is a complex Hilbert space,
then $\Phi\: \hilH \xto \hilH_a^*$ is bijective and 
consequently a linear homeomorphism.
\end{theorem}

\subparagraph{}
The proof of this theorem can be found in nearly every textbook
on functional analysis.

\paragraph{} %- inner product in the conjugate space
%------------------------------------------------------------------------------%
Theorem \ref{RIESZ_REPRESENTATION_THEOREM} implies that
for $\phi \xeof \hilH_a^*$ there is a unique element $x_{\phi} \xeof \hilH$,
which fulfills the equation
$\scalar{x_{\phi}}{y} \eq \phi(y)$ for every $y \xeof \hilH$.
Moreover,
the Riesz map gives rise to introduce an inner product
on the conjugate space $\hilH_a^*$ by
\[
\scalar{\phi}{\psi} \eqdef \scalar{\Phi^{-1}(\phi)}{\Phi^{-1}(\psi)}
\]
for all $\phi,\psi \xeof \hilH_a^*$,
hence $\Phi$ is not only norm-preserving,
but also \textit{isometric},
that is,
$\scalar{\Phi x}{\Phi y} \eq \scalar{x}{y}$ for $x,y \xeof \hilH$.

\subparagraph{} %- the Riesz map in real Hilbert spaces
On the other hand,
if $\hilH$ is a \textit{real} Hilbert space,
we may replace $\hilH_a^*$ by $\hilH^*$ and verify
that the Riesz map $\Phi$ again is a \textit{norm-isomorphism}
between the spaces $\hilH$ and $\hilH^*$,
since in this case
the inner products in $\hilH$ and $\hilH^*$
are also linear in the second argument.

\subparagraph*{}
The representation theorem shows that
real and complex Hilbert spaces are reflexive.

\subsection{} %-- adjoint operators
%------------------------------------------------------------------------------%
In operator theory it is often advantageous to combine
the study of a densely defined linear operator
with that of its \textit{adjoint}.
The rest of this section is a brief review
on important concepts {\wrt} adjoint operators.

\subparagraph{} %- definition of the adjoint operator
Let $\banX,\banY$ be arbitrary Banach spaces and
$\banX^*,\banY^*$ be the adjoint spaces.
If $\linA \: \domain{\linA} \xsubset \banX \xto \banY$
is a linear operator with domain of definition $\domain{\linA}$
dense in $\banX$,
then the domain of definition $\domain{\linA^*}$
of the {adjoint operator} $\linA^*$ is defined by
\[
\domain{\linA^*} \eqdef
\{\, \phi \xeof \banY^*:
\phi \circ \linA \: \domain{\linA} \xto \F \mbox{ is continuous}\,\}.
\]
For $\phi \xeof \domain{\linA^*} $ let
$\linA^* \phi := \overline{\phi \circ \linA}$,
the unique extension of the continuous linear functional $\phi \circ \linA$
to the whole of $\banX$,
which implies $\linA^* \phi \xeof \banX^*$.
Thus the linear operator $\linA^*$
has domain space $\banY^*$ and range space $\banX^*$,
and the domain of definition $\domain{\linA^*}$
forms a linear subspace of $\banY^*$,
on which $\linA^*$ is linear.

\subparagraph{}
We emphasize that
$\linA$ is supposed to be densely defined in order to introduce its adjoint.
But the adjoint $\linA^*$ is not necessarily densely defined.
Indeed,
it can even happen that $\domain{\linA^*} = \{0\}$.
The question however,
whether $\domain{\linA^*}$ is dense in $\banX^*$,
depends essentially on the closability of the operator $\linA$,
see assertion (d) in theorem \ref{MAIN_THEOREM_ON_ADJOINT_OPERATORS} below.

\subparagraph{} %- orthogonality relationships
The usage of the adjoint operator is founded by the validness
of two relationships between the kernels and ranges
of the operators $\linA$ and $\linA^*$,
namely
\begin{equation}\label{ORTHOGONALITY_RELATIONS}
\overline{\range{\linA^*} } \= \kernel{\linA}^{\bot}
\quad \and \quad
\overline{\range{\linA} } \= ^{\bot}\kernel{\linA^*},
\end{equation}
where for any subset $M \xsubset \banX$ and $N \xsubset \banX^*$
the \textit{annihilator}     $M^{\bot}$ and
the \textit{pre-annihilator} $^{\bot}N$
are defined by
\[
\begin{cases}[1.2]
\quad M^{\bot} \eqdef \{\,x^*\xeof \banX^* \: \scalar{x}{x^*} = 0
      \fa x   \xeof M \,\},\\
\quad ^{\bot}N \eqdef \{\,x  \xeof \banX   \: \scalar{x}{x^*} = 0
      \fa x^* \xeof N \,\}\\
\end{cases}
\]
and $\scalar{x}{x^*} := x^*(x)$ denotes the pairing of elements in $\banX$
and $\banX^*$.
Checking the validness of the relations \eqref{ORTHOGONALITY_RELATIONS}
is elementary,
hence we do not want to elaborate this in detail.

\paragraph{} %- properties of the adjoint
%------------------------------------------------------------------------------%
Adjoint operators probably enjoy a couple of properties:

\begin{theorem}\label{MAIN_THEOREM_ON_ADJOINT_OPERATORS}
Let $\banX,\banY$ be Banach spaces and
$\linA,\linB$ be densely defined linear operators with domain space
$\banX$ and range space\,$\banY$.

\begin{itemize}[leftmargin=9mm]
\item[(a)] $\linA \prec \linB$ implies $\linB^* \prec \linA^*$.
\item[(b)] The adjoint operator $\linA^*$ is closed.
\item[(c)] If $\linA$ is closable, then $\linA^* \eq \overline{\linA}\,^*$.
\item[(d)] If $\linA$ is closable, then $\linA^*$ is densely defined,
           which implies that
           the bi-adjoint $\linA^{**} := (\linA^*)^*$ exists.
\end{itemize}
\end{theorem}

\begin{comment}
\Proof:\\
Assertion (a) follows from the definition of the adjoint operator:
let $y \xeof \domain{\linB^*}$,
which means that there is $y^* \xeof \hilH$ such that
\[
\scalar{\linB x}{y} \= \scalar{\linA x}{y} \= \scalar{x}{y^*}
\quad \fa x \xeof \domain{\linB},
\]
from which we conclude
that $y \xeof \domain{\linA^*}$ and $\linA^* y = \linB^* y$.

To prove assertion (b),
let $(y_n)$ be a sequence in $\domain{\linA^*}$
such that $y_n \xto y \xeof \banY^*$ and $\linA^* y_n = z \xeof \banX^*$.
It is to show that $y \xeof \domain{\linA^*}$ and $\linA^* y = z$.
In fact,
the pairing with $\linA x$ for any element $x \xeof \domain{\linA}$ yields
\[
\scalar{\linA x}{y} = \lim_{n \to \infty} \scalar{\linA x}{y_n}
                    = \lim_{n \to \infty} \scalar{x}      {\linA^* y_n}
                    =                     \scalar{x}      {z}.
\]

To prove assertion (c),
let us first mention that
since $\overline{\linA}$ is an extension of $\linA$,
we immediately obtain $\overline{\linA}\,^* \prec \linA^*$ by (a).
To show that $\linA^* \prec \overline{\linA}\,^*$,
let $y \xeof \domain{\linA^*}$.
If $x \xeof \domain{\overline{\linA}}$,
then there is a sequence $(x_n)$ in $\domain{\linA}$ such that
$x_n \xto x$ and $\linA x_n \xto \overline{\linA} x$.
Since $\scalar{\linA x_n}{y} \eq \scalar{x_n}{\linA y}$ for all $n \xeof \N$,
we obtain
\[
\scalar{\overline{\linA} x}{y} \= \scalar{x}{\linA^* y}
\]
by letting $n \xto \infty$,
which shows that $y \xeof \domain{\overline{\linA}\,^*}$ and
$\overline{\linA}\,^* y \eq \linA^* y$.

Next we show that $\linA^*$ is densely defined.
Assume that $\domain{\linA^*}$ is not a dense subset in $\hilH$.
Then there exists $0 \neq z \xeof \domain{\linA^*} ^{\bot}$
fulfilling the equations
\[
\scalar{y}{z} \eq 0 \fa y \in \domain{\linA^*},
\]
which means nothing but $(z,0) \xeof (\Gamma_{\linA^*})^{\bot}$.
But then the set $ \Psi \, (\Gamma_{\linA^*}^{\bot})$
contains $(0,z)$ and cannot be the graph of an operator.
If $\linA$ is closable,
we necessarily would have
$\Gamma_{\overline{\linA}} = \Psi \, (\Gamma_{\linA^*})^{\bot}$
in opposite to the conclusion above.
Hence $\linA^*$ must be densely defined.
\qed
\end{comment}

\paragraph{} %- formally adjoint pairs
%------------------------------------------------------------------------------%
We also introduce the concept of \textit{adjoint pair} of linear operators.
The both densely defined linear operators
$\linA \: \domain{\linA} \xsubset \banX   \xto \banY$ and
$\linB \: \domain{\linB} \xsubset \banY^* \xto \banX^*$
are said to be \textit{formally adjoint} to each other
or to be a \textit{formally adjoint pair} $(\linA,\linB)$ if
\[
\scalar{\linA \,x}{f}_{\banY \xtimes \banY^*} \=
\scalar{x}{\linB f}_{\banX \xtimes \banX^*}
\]
holds for all $x \xeof \domain{\linA}$ and $f \xeof \domain{\linB}$.

\subparagraph{} %- restricted adjoint
In addition,
we define the \textit{restricted adjoint} $\linB_r$
of the operator $\linB$ according to $\domain{\linB_r} :=$
\[
\{ x \xeof \banX \: \mbox{ there is } y \xeof \banY \mbox{ such that }
(\linB f)\,x \eq f(y) \fa f \xeof \domain{\linB} \}
\]
and $\linB_r \, x := y$.
The mapping $\linB_r$ acts between the Banach spaces $\banX$ and $\banY$
and proves to be linear.
We may reformulate the notation of {formally adjoint pair} by saying that
$\linA$ and $\linB$ are formally adjoint to one another
if $\linB_r$ is an extension of $\linA$.

\subparagraph{} %- realization of a formally adjoint pair
To the end,
by a \textit{realization} of a formally adjoint pair $(\linA,\linB)$
we mean any densely defined and closed linear operator $R$
lying between $\linA$ and $\linB_r$,
so $\linA \prec R \prec \linB_r$ must hold.

\paragraph{} %- compactness of the adjoint
%------------------------------------------------------------------------------%
Recall that the adjoint $\linA^*$ of a linear operator $\linA$ forms
a mapping from the adjoint space $\banY^*$ to the adjoint space $\banX^*$.
With respect to the question of compactness of $\linA$ and $\linA^*$,
\textsc{Schauder}\textit{'s theorem} holds,
whose proof can be found for example in \cite{Yo}, pp. 282-283:

\begin{theorem}
A continuous linear operator $\linA \: \banX \xto \banY$
between Banach spaces $\banX$ and $\banY$ is compact
\iff $\linA^*$ is compact.
\end{theorem}

\subsection{} %-- the Hilbert adjoint
%------------------------------------------------------------------------------%
In the Hilbert space setting,
where $\banX \eq \banY \eq \hilH$ is assumed,
the sets $M^{\bot}$ and $^{\bot}N$ are nothing
but the orthogonal spaces of $M$ and $N$,
where $N^{\bot} \xsubset \hilH^*$ can be seen as a linear subspace of $\hilH$,
which is a consequence of \textsc{Riesz}'s representation theorem
(establishing a one-to-one relationship
between the spaces $\hilH$ and $\hilH^*$).

\subparagraph{}
Furthermore,
the adjoint $\linA^*$ of a densely defined linear operator $\linA$
can be regarded as a mapping acting in $\hilH$.
In fact,
if $\linA$ is linear and densely defined in $\hilH$,
then the domain of definition $\domain{\linA^*}$
of the adjoint $\linA^*$ consists of those elements $y \xeof \hilH$,
for which there is $y^* \xeof \hilH$ fulfilling
\[
\scalar{\linA \,x}{y} \eq \scalar{x}{y^*} \fa x \xeof \domain{\linA}.
\]
We then define $\linA^* y := y^*$ and denote $\linA^*$
the \textit{Hilbert adjoint} of $\linA$.

\paragraph{}
%------------------------------------------------------------------------------%
Recall that if a linear operator is closable,
then its adjoint operator is densely defined and the bi-adjoint is well-defined.
In the Hilbert space setting,
however,
we can say more about the bi-adjoint:

\begin{theorem}
If the linear operator $\linA$ is closable,
then $\linA^{**} \eq \overline{\linA}$,
which implies that
if $\linA$ is closed,
then $\linA^{**} \eq \linA$ holds.
\end{theorem}

\begin{comment}
\Proof:
An essential tool to prove properties of adjoint operators
in Hilbert space is the mapping
\[
\begin{cases}[1.2]
\quad \Psi \: \hilH \xtimes \hilH \xto \hilH \xtimes \hilH,\\
\quad (x,y) \mapsto (-y,x),\\
\end{cases}
\]
which obviously owns the property $\Psi^2 \eq \1_{\hilH}$.
But this mapping is also bilinear and unitary,
since
\[
\norm{\Psi(x,y)} = \norm{(x,y)} := (\norm{x} ^2 + \norm{y} ^2)^{1/2}
\fa (x,y) \in \hilH \xtimes \hilH.
\]
It also follows that for any subspace $U \xsubset \hilH$ we have
$\Psi(U^{\bot}) \eq \Psi(U)^{\bot}$.
Now,
if $\Gamma_A$ is the graph of a linear operator $\linA$,
then
\[
(\Gamma_A)^{\bot} = \{(y,z) \in \hilH \xtimes \hilH:
\scalar{x}{y} + \scalar{z}{\linA x} = 0 \fa x \in \domain{\linA} \}.
\]
On the other hand,
the image of the set $(\Gamma_A)^{\bot}$ under $\Psi$ is given by
\[
\Psi \, (\Gamma_A)^{\bot} =
\{(-z,y) \in \hilH \xtimes \hilH:
\scalar{x}{y} + \scalar{\linA x}{z} = 0
\fa x \in \domain{\linA}\}
\]
\[
=
\{(z,y) \in \hilH \xtimes \hilH:
z \in \domain{\linA^*} \and y = \linA^*z\}
= \Gamma_{\linA^*},
\]
so after all
$\Psi \,(\Gamma_{\linA})^{\bot} \eq \Gamma_{\linA^*}$ holds.
Since the set $\Gamma_A$ forms a linear subspace of the Hilbert space
$\hilH \xtimes \hilH$,
we even have the identity
\begin{equation}\label{GRAPH}
\overline{\Gamma_A} =
(\Gamma_A)^{\bot \bot} =
\Psi^2 \, (\Gamma_A)^{\bot \bot} =
(\Psi \, (\Psi (\Gamma_A)^{\bot}))^{\bot} =
(\Psi \, \Gamma_{\linA^*})^{\bot},
\end{equation}
that will especially be used in the proof of statement (d) in our next theorem.

To prove assertion (d),
recall that the operator $\linA^{**}$ exists by (c).
Then the identity
$\Psi \,(\Gamma_{\linA})^{\bot} \eq \Gamma_{\linA^*}$
together with \eqref{GRAPH} implies
$\Gamma_{\overline{A}} \eq \overline{\Gamma_A} \eq \Gamma_{\linA^{**}}$,
which implies $\overline{\linA} = \linA^{**}$.
\qed
\end{comment}

\subsection{} %-- normal operators in Hilbert space
%------------------------------------------------------------------------------%
The concept of Hilbert adjoint gives rise to introduce the following property
that a Hilbert space operator possibly owns:
the densely defined linear operator
$\linA \: \domain{\linA} \xsubset \hilH \xto \hilH$
is called \textit{normal},
if $\linA \, \linA^* \eq \linA^* \, \linA$ in the sense of section 2.

\subparagraph{} %- examples of normal operators
Let us consider some special cases of operators in Hilbert space,
wich are defined by means of the adjoint.
A densely defined linear operator $\linA$ acting in $\hilH$ is called

\begin{itemize}[leftmargin=5mm]
\item[1.]
\textit{symmetric} if $\linA \prec \linA^*$;
\item[2.]
\textit{self-adjoint} if $\linA \eq \linA^*$;
\item[3.]
\textit{skew-adjoint} if $\linA \eq - \linA^*$;
\item[4.]
\textit{unitary},
if $\domain{\linA} \eq \hilH$, $\linA$ is bijective
and $\linA^{-1} \eq \linA^*$.
\end{itemize}

\subparagraph{}
It should be clear that every self-adjoint or skew-adjoint or unitary operator
is normal.
Symmetric and self-adjoint operators and their spectral theory
will be discussed more in detail in sections 7 and 8.

%------------------------------------------------------------------------------%
\section{Resolvent and Spectrum}                                               %
%------------------------------------------------------------------------------%
%-- 
%-- 4. RESOLVENT AND SPECTRUM
%--    ----------------------
\textit{Spectral theory} is a subject within functional analysis
that proves to be helpful in the study of abstract linear equations
of type $\linA x \eq y$
as well as of eigen\-value equations of type $\linA x {-} \lambda x \eq y$.

\subparagraph{} %- complexification of a real Banach space
In spectral theory one is working with complex Banach spaces.
But notice that each real Banach space $\banX$ can be made
into a complex Banach space $\banX_{\C}$,
where the related procedure is denoted the \textit{complexification} of $\banX$:
let $\banX_{\C} := \banX + i \banX \eq \{x+iy: x,y \xeof \banX\}$
and furnished with the addition and multiplication
derived from those in $\banX$.
Then the norm in $\banX_{\C}$ is defined according to
\[
\norm{x+iy} \eqdef
\sup \, \{\norm{x} \cos \theta + \norm{y} \sin \theta: \theta \xeof [0,2\pi)\},
\]
in which the space $\banX_{\C}$ proves to be complete.

\subsection{} %-- resolvent set and spectrum
%------------------------------------------------------------------------------%
Let $\1_{\banX}$ be the identity mapping in a Banach space $\banX$.
Given a linear operator $\linA \: \domain{\linA} \xsubset \banX \xto \banX$,
the basic idea in spectral theory is to split the complex plane $\C$
into two disjoint sets $\rho(\linA)$ and $\sigma(\linA)$,
denoted the \textit{resolvent set} and the \textit{spectrum} of $\linA$,
respectively,
where by definition $\zeta \xeof \rho(\linA)$ \iff

\begin{itemize}[leftmargin=12mm]
\item[(R-1)]
$\linA {-} \zeta := \linA {-} \zeta \, \1_{\banX} \: \domain{\linA} \xto \banX$
is bjective,

\item[(R-2)]
the inverse mapping $(\linA {-} \zeta)^{-1} \: \banX \xto \domain{\linA}$
is continuous,
\end{itemize}

otherwise $\zeta$ is said to belong to $\sigma(\linA)$.
The numbers in the sets $\rho(\linA)$ and $\sigma(\linA)$
are called \textit{regular values} and \textit{spectral values},
respectively.
If $\linA$ is closed,
however,
property (R-2) is a consequence of the closed graph theorem
and does not have to be postulated by its own.

\subparagraph{} %- properties of the resolvent set
The resolvent set $\rho(\linA)$
may be the empty set, a proper subset of the complex plane
or the set $\C$ itself.
We summarize a couple of further properties of the resolvent set
in the following

\begin{theorem}
Let $\linA$ be a linear operator acting in $\banX$.
Then:
\begin{itemize}[leftmargin=9mm]
\item[(a)]
$\rho(\linA)$ is open.
\item[(b)]
If $\rho(\linA) \neq \emptyset$, then $\linA$ is closed
(and often denoted \textit{regular}).
\item[(c)]
If $\linA$ is continuous and everywhere defined in $\banX$,
then we actually have $\emptyset \neq \rho(\linA) \neq \C$ and
$\abs{\lambda} \leq \norm{\linA}$ for all $\lambda \xeof \sigma(\linA)$.
\end{itemize}
\end{theorem}

\begin{comment}
\Proof:
To prove assertion (b),
let $\zeta \xeof \rho(\linA)$ such that
$\linA {-} \zeta$ is bijective with domain of definition
$\domain{(\linA {-} \zeta)^{-1}} \eq \banY$ and range
$\range{(\linA {-} \zeta)^{-1}} \eq \domain{\linA}$.
Let $x_n \xto x$ with $x_n \xeof \domain{\linA}$ for $n \xeof \N$
and $\linA x_n \xto y$.
Then
\begin{equation}\label{XXX_TOOL}
x = \lim_{n \xto \infty} x_n =
\lim_{n \xto \infty} (\linA {-} \zeta)^{-1}(\linA {-} \zeta) x_n =
(\linA {-} \zeta)^{-1}(y-\zeta x).
\end{equation}
Thus $x \xeof \range{(\linA {-} \zeta)^{-1}} \eq \domain{\linA}$.
Applying $(\linA {-} \zeta)^{-1}$ from left to \eqref{XXX_TOOL},
we obtain $\linA x - \zeta x \eq y -\zeta x$,
hence $\linA x = y$,
which is nothing but closeness of $\linA$.
\qed
\end{comment}

\subsection{} %-- resolvent functions
%------------------------------------------------------------------------------%
If $\banX$ is a Banach space and $\linA$ is assumed
to be regular,
the operator-valued mapping
$R(\linA,\cdot) \: \rho(\linA) \xto \BO{\banX}$
defined by
\[
\zeta \mapsto R(\linA,\zeta) \eqdef (\linA {-} \zeta)^{-1}
\]
is called the \textit{resolvent function}
or briefly the \textit{resolvent} of $\linA$.

\subparagraph{} %- vector-valued holomorphic functions
Remember the concept of a vector-valued holomorphic function:
if $\G \xsubset \C$ is an open domain,
a function $f \: \G \xto \banX$ is called \textit{holomorphic}
if for every $\phi \xeof \banX^*$ the composition $\phi \circ f \: \G \xto \C$
is holomorphic in the sense of complex analysis.
The function $f$ is holomorphic \iff $f$ is \textit{analytic},
that is,
if $f(\zeta)$ for $\zeta$ close to $\zeta_0 \xeof \openO$ 
can locally be represented by a power series
\[
f(\zeta) = \sum_{k=0}^{\infty} a_k(\zeta_0) \, (\zeta - \zeta_0)^k
\quad \for \; \abs{\zeta-\zeta_0} < R(\zeta_0).
\]

\begin{theorem}
The resolvent is holomorphic in every $\zeta_0 \xeof \rho(\linA)$,
and $R(\linA,\zeta)$ can be represented by a power series
\[
\sum_{n \xeof \N} \, b_n(\zeta_0) \, (\zeta-\zeta_0)^n
\]
for $\abs{\zeta-\zeta_0} $ sufficiently small
and with coefficents $b_n(\zeta_0) \xeof \BO{\banX}$.
\end{theorem}

\subparagraph{} %- resolvent equation
All the mappings $R(\linA,\zeta)$ are coupled by the \textit{resolvent equation}
\begin{equation}\label{RESOLVENT_EQUATION}
R(\linA,\zeta_1) - R(\linA,\zeta_2)
\=
(\zeta_2-\zeta_1) \, R(\linA,\zeta_1) \circ R(\linA,\zeta_2)
\end{equation}
for $\zeta_1,\zeta_2 \xeof \rho(\linA)$,
which especially implies that
the linear operators $R(\linA,\zeta_1)$ and $R(\linA,\zeta_2)$
commute with each another.

\paragraph{} %- norm estimate of the resolvent
%------------------------------------------------------------------------------%
Estimates for the operator norm of resolvents play an important role
in later sections.
In fact,
the estimate
\begin{equation}\label{RESOLVENT_ESTIMATE_I}
\norm{R(\linA,\zeta)} \, \geq \,
r(R(\linA,\zeta))     \, \geq \,
\mbox{dist}(\zeta,\sigma(\linA))^{-1}
\end{equation}
holds for any $\zeta \xeof \rho(\linA)$,
where $r(R(\linA,\zeta)) \eq
\sup \{ \abs{\lambda} \: \lambda \xeof \sigma(R(\linA,\zeta)) \}$
is the \textit{spectral radius} of $R(\linA,\zeta)$ and
$\mbox{dist}(\zeta,\sigma(\linA))$ is the minimal distance
between the singleton $\{\zeta\}$ and the spectrum $\sigma(\linA)$.

\paragraph{} %- spectral mapping theorem
%------------------------------------------------------------------------------%
The estimate \eqref{RESOLVENT_ESTIMATE_I} is a consequence
of the so-called \textit{spectral mapping theorem},
establishing a connection between the spectra of the resolvents
$R(\linA,\zeta)$ and the operator $\linA$ itself and reads as follows:

\begin{theorem}
Let $\linA$ be a linear operator with spectrum $\sigma(\linA)$ and resolvent
function $R(\linA,\cdot)$.
Then for $\zeta \xeof \rho(\linA)$ we have
\[
\sigma(R(\linA,\zeta)) \setminus \{0\} \=
\{(\lambda-\zeta)^{-1} \: \lambda \xeof \sigma(\linA)\}.
\]
\end{theorem}

\subsection{} %-- various kinds of spectra
%------------------------------------------------------------------------------%
We next will have a look at the spectral values
of a closed linear operator $\linA$ acting in a Banach space $\banX$
and study for $\lambda \xeof \sigma(\linA)$ the possible reasons
that prevent the operator $\linA {-} \lambda$ to be bijective.

\begin{itemize}[leftmargin=4mm,itemsep=5pt]
\item[1.]
Recall from linear algebra
that a complex number $\lambda$ is called an \textit{eigen\-value}
of a linear operator $\linA$,
if there is an element $x \xeof \banX$ such that
$x \neq 0$ and $\linA \,x \eq \lambda x$.
In this case the element $x$ is called an \textit{eigenvector} of $\linA$
{\wrt} $\lambda$,
and all these eigenvectors form a linear subspace of $\banX$,
called the \textit{eigenspace} of $\linA$ for $\lambda$.

\subparagraph{} %- point spectrum
The \textit{point spectrum} of $\linA$ is defined to be
the set $\sigma_p(\linA)$ of all eigen\-values of $\linA$.
Clearly,
for $\lambda \xeof \sigma_p(\linA)$ the operator
$\linA {-} \lambda$ is not injective,
since its kernel is non-trivial.
\end{itemize}

If $\linA {-} \lambda$ is injective,
but not surjective,
then $(\linA {-} \lambda)^{-1}$ is defined on the range $\range{\linA {-} \lambda}$.
We distinguish between two cases:

\begin{itemize}[leftmargin=4mm]
\item[2.]
if $\range{\linA {-} \lambda}$ is dense in $\banX$,
then $\lambda$ is said to be an element of the set $\sigma_c(\linA)$,
called the \textit{continuous spectrum} of $\linA$,
otherwise

\item[3.]
the number $\lambda$ belongs to the \textit{residual spectrum}
$\sigma_r(\linA)$.
\end{itemize}

\subparagraph{}
In summary,
for each closed linear operator
$\linA \: \domain{\linA} \xsubset \banX \xto \banX$
the set of complex numbers forms a disjoint union
of the resolvent set and the three different kinds of spectra, namely
\[
\C \= \rho    (\linA) \, \sqcup
        \, \sigma_p(\linA) \, \sqcup
        \, \sigma_c(\linA) \, \sqcup
        \, \sigma_r(\linA).
\]
The decomposition of the spectrum into these three parts
is the classical approach.

\paragraph{} %- approximate point spectrum
%------------------------------------------------------------------------------%
The \textit{approximate point spectrum} $\sigma_{ap}(\linA)$
of a closed linear operator $\linA$ is the set of $\lambda \xeof \C$,
for which $\linA {-} \lambda$ is not injective
or $\range{\linA {-} \lambda}$ is not closed in $\banX$.
The elements of $\sigma_{ap}(\linA)$ are called
the \textit{approximate eigenvalues} of $\linA$.
Any number $\lambda \xeof \C$ belongs to $\sigma_{ap}(\linA)$
\iff
there is a sequence $\{x_n\}_{n \in \N} \xsubset \domain{\linA}$ such that
$\norm{x_n} \eq 1$ for all $n \xeof \N$ and
$\lim_{\,n \to \infty} \norm{\linA \,x_n - \lambda x_n} \eq 0$.
Since $\linA$ is closed,
the topological boundary $\partial \sigma(\linA)$ of the spectrum
forms a subset of $\sigma_{ap}(\linA)$.

\subparagraph{}
It is easy to see that
each eigenvalue of $\linA$ belongs to $\sigma_{ap}(\linA)$,
so this set generalizes the point spectrum $\sigma_p(\linA)$.
But $\sigma_{ap}(\linA)$ can be empty only in the case,
where $\sigma(\linA) \eq \emptyset$ or $\sigma(\linA) \eq \C$.

\paragraph{} %- resolvent set and spectrum of adjoint operators
%------------------------------------------------------------------------------%
Up to now we did not assume that
the linear operators under consideration are densely defined.
But if this is true for a linear operator $\linA$,
then the adjoint $\linA^*$ is uniquely defined,
and one may ask for relations between the spectra of $\linA$ and $\linA^*$.
In fact,
it turns out that
the resolvent set $\rho(\linA^*)$ and spectrum $\sigma(\linA^*)$
are mirrored sets of $\rho(\linA)$ and $\sigma(\linA)$,
and each one can be obtained by the other one by reflection at the real line.
There is also an interesting identity between the residual spectrum of $\linA$
and the point spectrum of $\linA^*$,
namely
$\sigma_r(\linA) \eq \sigma_p(\linA^*)$ (see \cite{EN}, p.\,161).

\subsection{} %-- spectrum and functional calculus for bounded operators
%------------------------------------------------------------------------------%
If a linear operator $\linA$ is bounded and defined in the entire space $\banX$,
we have $\sigma(\linA) \neq \emptyset$,
which is a consequence of \textsc{Liouville}'s theorem in complex analysis,
quoting that a bounded and holomorphic function $f \:\C \xto \C$
must be constant.

\subparagraph{} %- spectral radius
Furthermore,
the \textit{spectral radius} of $\linA$ is defined according to
\[
r(\linA) \eqdef \sup \, \{\abs{\lambda} \: \lambda \xeof \sigma(\linA)\}.
\]
It can be shown that $r(\linA) \leq \norm{\linA}$,
and since each resolvent $R(\linA,\zeta)$ is bounded,
we have the estimation
\begin{equation}\label{RESOLVENT_ESTIMATE_2}
\norm{R(\linA,\zeta)} \, \geq \,
r(R(\linA,\zeta))     \, \geq \,
\delta(\zeta,\sigma(\linA))^{-1}
\end{equation}
for any $\zeta \xeof \rho(\linA)$,
where $\delta(\zeta,\sigma(\linA))$ is the minimal distance between
the singleton $\{\zeta\}$ and the spectrum $\sigma(\linA)$.
Notice that inequality \eqref{RESOLVENT_ESTIMATE_2} is a consequence
of the spectral mapping theorem.

\paragraph{} %-  Riesz-Dunford calculus
%------------------------------------------------------------------------------%
A further main subject in this section
is the exposition of the \textit{holomorphic functional calculus},
sometimes also called the \textsc{Riesz-Dunford} \textit{calculus}
for bounded linear operators.
It allows to define functions $f(\linA)$ for a given linear operator $\linA$.

\subparagraph{}
To introduce this kind of functional calculus,
let $H(\G)$ be the space of analytic functions in $\G$.
We denote $H^{\infty}(\G)$ the subspace
of all \textit{bounded} functions in $H(\G)$,
which is a Banach space {\wrt} the norm
\[
\norm{f}_{\infty} \eqdef \sup_{\zeta \in \G} \, \abs{f(\zeta)} .
\]

\subparagraph{} %- contour integrals
Let $\G \xsubset \C$ be an open and connected subset of the complex plane,
$g \: [a,b] \xto \G$ be a rectificable path
and $f \: \G \xto \C$ be a complex-valued function.
Then
\[
\int_{\,\gamma} \, f(\zeta) \, d\zeta \eqdef \int_a^b f \, dg
\]
is called the \textit{contour integral} of $f$ {\wrt} the function $g$,
where $\gamma$ is the image of $[a,b]$ under $g$.

\paragraph{} %- Cauchy's integral theorem
%------------------------------------------------------------------------------%
\textsc{Cauchy}\textit{'s integral theorem} is a corner stone in
complex analysis and quotes that if $f$ is analytic in $\G$,
then
for \textit{every} closed curve $\gamma$ we have
\[
\int_{\,\gamma} \, f(\zeta) \, d\xi \= 0.
\]
Furthermore,
\textsc{Cauchy}\textit{'s integral formula} says
that for each $z_0 \xeof \G$ belonging to the interior
of a closed curve $\gamma$ the value $f(z_0)$ is nothing but
\begin{equation}\label{CAUCHY_INTEGRAL_FORMULA_1}
f(z_0) \= \frac{1}{2 \pi i} \int_{\,\gamma} \, \frac{f(z)}{z-z_0} \, dz\,,
\end{equation}
where the \textit{winding number} of $\gamma$ has value $1$.
This representation of $f(z_0)$ does not depend on the shape of $\gamma$.

\subparagraph{}
If we consider $z_0 \xeof \C$ as a linear operator
$Z_0 \: \C \xto \C \xeof \BO{\C}$ defined by
$Z_0 \,\zeta := z_0 \cdot \zeta$ for $\zeta \xeof \C$,
then the spectrum of $\Z_0$ consists of the single point $z_0$
and the resolvent adopts the form
$R(Z_0,\zeta) \eq (\zeta-z_0)^{-1}$.
Using \eqref{CAUCHY_INTEGRAL_FORMULA_1}, 
we obtain
\begin{equation}\label{CAUCHY_INTEGRAL_FORMULA_2}
f(z_0) \=
\frac{1}{2 \pi i} \, \int_{\,\gamma} \, f(\zeta) \, R(Z_0,\zeta) \, d\zeta.
\end{equation}

\subparagraph{} %- Cauchy's integral formula
The formula \eqref{CAUCHY_INTEGRAL_FORMULA_2} serves as a pattern
for the \textit{holomorphic} or \textit{analytic functional calculus}
in unital and complex Banach algebras like the operator algebra $\BO{\banX}$.
Let $\linA$ be a bounded linear operator acting in $\banX$
(so the spectrum $\sigma(\linA)$ is non-emtpy and bounded),
$\gamma$ be a closed and rectificable path with winding number $1$
containing $\sigma(\linA)$
and $f \: \G \xto \C$ be a bounded and analytic function
defined on an open set $\G \supset \sigma(\linA)$.
Then the linear operator $f(\linA)$ is defined
by the \textit{Riesz-Dunford integral}
\begin{equation}\label{RIESZ-DUNFORD_FORMULA}
f(\linA) \eqdef
\frac{1}{2 \pi i} \, \int_{\,\gamma} \, f(\zeta) R(\linA,\zeta) \, d\zeta.
\end{equation}
The right hand side is a vector-valued Riemann-Stieltjes integral,
and convergence of the intermediate sums to the element
$f(\linA)$ is understood in the sense of the norm topology in $\BO{\banX}$.
The value $f(\linA)$ is well-defined,
since $f$ is bounded and analytic.

\subparagraph{}
Recall
that the set $H^{\infty}(\G)$ of bounded and analytic functions
$f \: \openO \xto \C$ defined in an open neighbourhood $\openO$
of $\sigma(\linA)$ forms a Banach algebra.
Thus the mapping
\[
\begin{cases}[1.2]
\quad \Phi_{\linA} \: H^{\infty}(\openO) \xto \BO{\banB},\\
\quad
f \mapsto f(\linA) \; \mbox{as defined in \eqref{RIESZ-DUNFORD_FORMULA}} \\
\end{cases}
\]
forms a \textit{Banach algebra homomorphism}
from $H^{\infty}(\openO)$ to $\BO{\banB}$,
that is,
the mapping $\Phi_{\linA}$ is linear, norm-preserving
and respects the multiplicative operations in both algebras.
But $\Phi_{\linA}$ also maps polyomials $p \xeof H^{\infty}(\openO)$
to polynomials $p(\linA) \xeof \BO{\banB}$.
Due to these properties,
$\Phi_{\linA}$ is called a \textit{functional calculus},
which in addition deserves the appendix "analytic",
since the space $H^{\infty}(\openO)$ consists of analytic functions.

\subparagraph{}
It should be mentioned 
that for this kind of functional calculus the
\textit{spectral mapping theorem} is valid,
that is,
$\sigma(f(\linA)) \eq f(\sigma(\linA))$
for $\linA \xeof \BO{\banB}$ and $f \xeof H^{\infty}(\openO)$.
The resulting operator $f(\linA) \xeof \BO{\banX}$
defined by \eqref{RIESZ-DUNFORD_FORMULA}
is called the \textsc{Riesz-Dunford} \textit{operator}
associated with $\linA$ and $f \xeof H^{\infty}(\G)$.

\subsection{} %-- projections in a Banach space
%------------------------------------------------------------------------------%
In the following we suppose that
$\lambda$ is an \textit{isolated} point of the spectrum $\sigma(\linA)$,
so $\sigma(\linA) \cap \U_{\lambda} \eq \{\lambda\}$ holds
for every sufficiently small neighbourhood $\U_{\lambda}$ of $\lambda$.
Let $\gamma_{\lambda}$ be a closed and rectificable path,
which surrounds only $\lambda$,
but not any other element of $\sigma(\linA)$.
Then for $x \xeof \banX$ the contour integral
\[
\pi_{\lambda}(x) \eqdef
\frac{1}{2 \pi i} \, \int_{\,\gamma_{\lambda}} \, R(\linA,\zeta) x \, d\zeta
\]
gives rise to a further linear mapping $\pi_{\lambda}$,
called the \textsc{Riesz} \textit{projection} of $\linA$
{\wrt} the spectral value $\lambda$.

\paragraph{} %- properties of the Riesz projection
%------------------------------------------------------------------------------%
To the end,
we collect the properties of the mapping $\pi_{\lambda}$ to the following list:

\begin{itemize}[leftmargin=12pt]
\item[1.]
$\pi_{\lambda}$ is linear.

\item[2.]
$\pi_{\lambda}$ owns the \textit{projection property}
$\pi_{\lambda} \circ \pi_{\lambda} \eq \pi_{\lambda}$.

\item[3.]
The spectrum of this mapping only consists of the value $\lambda_0$ itself.

\item[4.]
$\pi_{\lambda}$ commutes with $\linA$,
so $\linA \circ \pi_{\lambda} \eq \pi_{\lambda} \circ \linA$ holds.

\item[5.]
The range $\Pi_{\lambda}$ of $\pi_{\lambda}$
is a closed linear subspace of $\banX$.

\item[6.]
There is a closed subspace $\Pi_{\lambda}^{\bot} \xsubset \banX$
such that
$\banX \eq \Pi_{\lambda} \, \oplus \, \Pi_{\lambda}^{\bot}$,
where $\Pi_{\lambda}^{\bot}$ is well-defined by the decomposition theorem.

\item[7.]
Both subspaces $\Pi_{\lambda}$ and $\Pi_{\lambda}^{\bot}$
are $\linA$-invariant, so
$\linA \, \Pi_{\lambda}        \xsubset \Pi_{\lambda}$ and
$\linA \, \Pi_{\lambda}^{\bot} \xsubset \Pi_{\lambda}^{\bot}$.
One also says that the subspaces 
$\Pi_{\lambda}$ and $\Pi_{\lambda}^{\bot}$
\textit{reduce} the entire space $\banX$.

\item[8.]
The mapping $\pi_{\lambda}^{\bot} := \1_{\banX} {-} \pi_{\lambda}$
is a projection again,
and both projections $\pi_{\lambda}$ and $\pi_{\lambda}^{\bot}$
are \textit{orthogonal} to each another,
that is,
$\pi_{\lambda}        \circ \pi_{\lambda}^{\bot} \eq 
 \pi_{\lambda}^{\bot} \circ \pi_{\lambda}        \eq \0$.
\end{itemize}

\subparagraph{}
The concept of {Riesz projection} can be extended into different directions.
First,
the linear operator $\linA$ may be assumed to be closed
instead of being continuous,
and second
the subset $\{\lambda\} \xsubset \sigma(\linA)$ may be replaced
by any connected component $\sigma_0(\linA)$ of the set $\sigma(\linA)$.

%------------------------------------------------------------------------------%
\section{Normally Solvable Operators}                                          %
%------------------------------------------------------------------------------%
%-- 
%-- 5. NORMALLY SOLVABLE OPERATORS
%--    ---------------------------
The general subject in this section is the review of linear operators
having closed range.
In fact,
given a densely defined linear operator $\linA$
acting between Banach spaces $\banX$ and $\banY$,
the orthogonality relations \eqref{ORTHOGONALITY_RELATIONS}
mentioned in section 3 can be improved
if the ranges of $\linA$ and $\linA^*$ are \textit{closed} subspaces
of $\banY$ and $\banX^*$,
so in these cases we actually would have
\[
{\range{\linA} } \= ^{\bot}\kernel{\linA^*}
\quad \and \quad
{\range{\linA^*} } \= \kernel{\linA}^{\bot}.
\]
One also says that
a linear operator $\linA$ is \textit{normally solvable}
if it is densely defined and
$\range{\linA} \eq \, ^{\bot}\kernel{\linA^*}$ holds,
which means by our next theorem that $\linA$ has closed range.

\subsection{} %-- closed range theorem
%------------------------------------------------------------------------------%
The \textit{closed range theorem}
is one of the classical results in functional analysis 
and provides equivalent conditions
for a densely defined and closed linear operator to own closed range:

\begin{theorem}\label{CLOSED_RANGE_THEOREM}
Let $\banX,\,\banY$ be Banach spaces
and $\linA \: \domain{\linA} \xsubset \banX \xto \banY$
be densely defined and closed.
Then it is equivalent to quote:

\begin{itemize}[leftmargin=9mm]
\item[(a)] $\range{\linA}$ is closed.
\item[(b)] $\range{\linA^*}$ is closed.
\item[(c)] $\range{\linA} $ is the pre-annihilator of $\kernel{\linA^*}$,
           {\ie} $\range{\linA} \eq ^{\bot}\kernel{\linA^*} $.
\item[(d)] $\range{\linA^*}$ is the annihilator of $\kernel{\linA}$,
           {\ie} $\range{\linA^*} \eq \kernel{\linA} ^{\bot}$.
\end{itemize}
\end{theorem}

\subparagraph{}
For a proof of this theorem see for example \cite{Yo}, pp. 205-207.
The statement was originally proved by \textsc{S.\,Banach}
for continuous linear operators and
have later been extended to closed linear operators by \textsc{F.\,Browder},
see \cite{Br1}, theorem\,1.1 and theorem\,1.2.

\subsection{} %- operators bounded below
%------------------------------------------------------------------------------%
It is natural to search for further conditions
that imply normal solvability of operators, and the first of them is
what is usually called the \textit{bounded-below theorem}.

\subparagraph{} %- bounded-below property
Let $\banX$ and $\banY$ be Banach spaces.
A linear operator $\linA$
with domain of definition $\domain{\linA} \xsubset \banX$
and range space\,$\banY$ is called \textit{bounded below}
if there is $\theta {\,>\,} 0$ such that
\begin{equation}\label{BOUNDED_BELOW_CONDITION}
\norm{\linA x}_Y \, \geq \, \theta \norm{x}_{\banX} \fa x \xeof \domain{\linA}.
\end{equation}

\subparagraph{} %- bounded below theorem
The \textit{bounded-below theorem} provides a fundamental assertion
in operator theory and is helpful to verify the closeness of $\range{\linA}$:

\begin{theorem}\label{BOUNDED_BELOW_THEOREM}
If $\linA$ fulfills \eqref{BOUNDED_BELOW_CONDITION},
then $\linA$ is injective and the operator
$\linA^{-1} \: \range{\linA} \xsubset \banY \xto \banX$ is bounded
with operator norm $\norm{\linA^{-1}} \leq \theta^{-1}$.
Furthermore,
if $\linA$ is closed,
then $\range{\linA}$ is a closed set in $\banY$.
\end{theorem}

Indeed,
injectivity of the operator $\linA$ can easily proved,
supposed \eqref{BOUNDED_BELOW_CONDITION} holds,
and closeness of $\range{\linA}$ is a consequence of sequentially closeness
of $\linA$.

\subsection{} %-- semi-Fredholm operators
%------------------------------------------------------------------------------%
We want to extend our considerations
to a more general class of linear operators.
Preliminary,
let $\alpha(\linA)$ (the \textit{nullity index} of $\linA$)
and $\beta(\linA)$ (the \textit{deficiency index} of $\linA$)
be the Hamel dimensions of the linear subspaces
$\kernel{\linA} \xsubset \banX$ and $\cokernel{\linA} \xsubset \banY$
of a given linear operator $\linA \: \domain{\linA} \xsubset \banX \xto \banY$.

\subparagraph{} %- Fredholm index
Both integers together can be used
to measure the deviation of $\linA$ from bijectivity.
Moreover,
if at least one of them is finite,
then the so-called \textit{Fredholm index}
\[
\chi(\linA) \eqdef
\begin{cases}[1.2]
\; + \infty, & \alpha(\linA) = \infty\\
\; \alpha(\linA) \, - \, \beta(\linA), &
\alpha(\linA) < \infty, \, \beta(\linA) < \infty\\
\; - \infty, & \beta(\linA) = \infty\\
\end{cases}
\]
is well-defined and takes values in the set
$\{{-} \infty\} \cup \N \cup \{{+} \infty\}$.
In the special case,
where $\alpha(\linA) {\,<\,} \infty$,
$\beta(\linA) {\,<\,} \infty$
and consequently $\chi(\linA) \xeof \Z$,
the operator $\linA$ is said to have \textit{finite index}.

\paragraph{} %- definition of semi-Fredholm operator
%------------------------------------------------------------------------------%
The linear operator $\linA \: \domain{\linA} \xsubset \banX \xto \banY$
is called \textit{semi-Fredholm},
if
\begin{itemize}[leftmargin=12mm]
\item[(F-1)] the range $\range{\linA}$ is closed,
\item[(F-2)] the kernel $\kernel{\linA} $ or the cokernel $\cokernel{\linA} $
             is finite-dimensional
             (rsp. $\alpha(\linA) {\,<\,} \infty$ or
                   $\beta(\linA)  {\,<\,} \infty$).
\end{itemize}

Notice the remarkable fact
that if $\linA$ is continuous with $\domain{\linA} \eq \banX$
and $\beta(\linA) {\,<\,} \infty$,
then the range $\range{\linA}$ is necessarily closed,
since it constitutes the topological complement
of the finite-dimensional linear subspace $\cokernel{\linA} \xsubset \banY$,
that is,
$\banX \eq \range{\linA} \oplus \cokernel{\linA}$.

\subparagraph{} %- lower and upper semi-Fredholm operators
We refine the concept of semi-Fredholm operator as follows:
a linear operator $\linA \: \domain{\linA} \xsubset \banX \xto \banY$
with closed range $\range{\linA}$ is called

\begin{itemize}[leftmargin=12mm]
\item[(F-3)]
\textit{upper semi-Fredholm},
if $\alpha(\linA) {\,<\,} \infty$ rsp. $\chi(\linA) {\,<\,} \infty$,

\item[(F-4)]
\textit{lower semi-Fredholm},
if $\beta(\linA) {\,<\,} \infty$ rsp. $\chi(\linA) {\,>\,} {-}\infty$.
\end{itemize}

\subparagraph{} %- example of upper semi-Fredholm operator
Examples of upper semi-Fredholm operators are the bounded-below ones.
In fact,
theorem \ref{BOUNDED_BELOW_THEOREM} implies that
these operators are upper semi-Fredholm with vanishing nullity index.

\paragraph{} %- Schauder estimates and Peetre's lemma
%------------------------------------------------------------------------------%
If a continuous linear operator $\linA$ satisfies
an abstract \textit{Schauder estimate},
one may quote a sufficient condition for $\linA$ to be upper semi-Fredholm,
which is known as \textit{Peetre}\textit{'s lemma}:

\begin{theorem}\label{PEETRE_THEOREM_II}
Let $\banX,\,\banY,\,\banZ$ be three Banach spaces,
where $\banX$ is compactly embedded into $\banZ$,
that is,
there is an injective and compact linear mapping $\iota \: \banX \xto \banZ$.
Then a continuous operator $\linA \: \banX \xto \banY$ is upper semi-Fredholm
iff there is $c {\,>\,} 0$ such that
\[
\norm{x}_{\banX} \, \leq \,
c \, (\norm{\iota x}_{\banZ} \+ \norm{\linA \,x}_{\banY})
\quad \quad (x \in \banX).
\]
\end{theorem}

\subparagraph{}
For a proof of this theorem, see \cite{Wl}, theorem 12.12, page 180.

\subsection{} %-- perturbation theory of semi-Fredholm operators
%------------------------------------------------------------------------------%
Nowadays there is a well-elaborated \textit{perturbation theory}
for semi-Fredholm operators available,
where conditions are studied,
under which a linear operator of the form $\linA {+}\linB$ is semi-Fredholm,
assumed $\linA$ is semi-Fredholm and $\linB$ fulfills certain properties.
In this case the operator $\linB$ is regarded as a perturbation of $\linA$.

\subparagraph{}
We next quote two theorems {\wrt} conditions on $\linB$,
which ensure the semi-Fredholm property of the operator $\linA {+}\linB$.
We basically assume that $\banX$ and $\banY$ are Banach spaces and
first quote \textit{Yood}\textit{'s theorem}:

\begin{theorem}\label{YOOD_THEOREM}
Let $\linA \: \banX \xto \banY$ be semi-Fredholm
and $\linB \: \banX \xto \banY$ be linear and bounded
with sufficiently small norm $\norm{\linB}$.
Then $\linA {+} \linB$ remains semi-Fredholm
with $\chi(\linA {+} \linB) \eq \chi(\linA)$.
\end{theorem}

\subparagraph{} %- punctured neighbourhood theorem
In other words,
this theorem asserts
that the set of semi-Fredholm operators is open in the set $\BO{\banX,\banY}$.
Especially linear operators of the form $\linA {+}\lambda$ for $\abs{\lambda} $ 
small enough are semi-Fredholm again.
But more can be said on the values $\alpha(\linA+\lambda)$,
$\beta(\linA+\lambda)$ and $\chi(\linA+\lambda)$
by \textsc{Gohberg}\textit{'s punctured neighbourhood theorem}:

\begin{theorem}\label{PUNCTURED_NEIGHBOURHOOD_THEOREM}
Let the continuous linear operator $\linA \: \banX \xto \banX$ be semi-Fredholm.
Then there is $\lambda_0 {\,>\,} 0$
such that for all $\lambda \xeof (0,\lambda_0)$ we have
$\alpha(\linA {+}\lambda) \leq \alpha(\linA)$,
$\beta (\linA {+}\lambda) \leq \beta (\linA)$
and
$\chi  (\linA {+}\lambda) \eq  \chi  (\linA)$.
\end{theorem}

\subsection{} %-- Fredholm operators
%------------------------------------------------------------------------------%
In the following we consider the case,
where both values are finite,
so we refine the concept of semi-Fredholm operator to a stronger one:
in addition to (F-1) to (F-4),
a continuous linear operator $\linA \: \banX \xto \banY$
acting between normed linear spaces $\banX$ and $\banY$ is called
\begin{itemize}[leftmargin=12mm,topsep=1mm]
\item[(F-5)]
\textit{Fredholm},
if $\chi(\linA)$ is finite,
{\ie} $\alpha(\linA) {\,<\,} \infty$ and $\beta(\linA) {\,<\,} \infty$.
\end{itemize}

\subparagraph{} %- Fredholm property of the adjoint
The relationships \eqref{ORTHOGONALITY_RELATIONS} imply
that if $\linA$ is densely defined and Fredholm,
then $\linA^*$ is also Fredholm and $\chi(\linA^*) \eq -\chi(\linA)$ holds.
Sure,
if $\linA$ is at once self-adjoint and Fredholm,
we have $\chi(\linA) \eq 0$.

\paragraph{} %- the Fredholm property of closed linear operators
%------------------------------------------------------------------------------%
We have introduced the Fredholm property of linear operators
under the assumption that
these are continuous and everywhere defined in their domain spaces.
If a linear operator $\linA$ is closed instead of continuous,
however,
we define the possible Fredholm property by regarding
$\linA$ as a continuous linear operator $\linA \: \domain{\linA} \xto \banY$,
where the domain of definition $\domain{\linA}$ is endowed
with the graph norm.
Indeed,
the pair $(\domain{\linA},\norm{\,\cdot\,}_{\linA})$
constitutes a Banach space due since $\linA$ is assumed to be closed.

\paragraph{} %- properties of Fredholm operators
%------------------------------------------------------------------------------%
We first collect a couple of properties of Fredholm operators:

\begin{itemize}[leftmargin=4mm]
\item[1.]
$\linA^*$ is Fredholm \iff $\linA$ is so.

\item[2.]
If one of $\linA$ or $\linA^*$ is Fredholm,
then the dimensions of the kernels and ranges of $\linA$ and $\linA^*$
are related to each other by
$\alpha(\linA^*) \eq \beta(\linA)$ and
$\beta(\linA^*)  \eq \alpha(\linA)$,
which implies $\chi(\linA^*) \eq - \chi(\linA)$.

\item[3.]
Fredholm operators acting in $\banX$ and
belonging to one and the same connected component
within the algebra $\BO{\banX}$ have identical index.
Thus,
if $\linA$ is Fredholm and the linear operator $\linB$ can be reached
by a continuous path in $\BO{\banX}$ with starting point $\linA$,
then $\linB$ is Fredholm,
too,
and $\chi(\linB) \eq \chi(\linA)$ holds.
\end{itemize}

\paragraph{} %- left and right parametrix
%------------------------------------------------------------------------------%
Our next subject is the representation of Fredholm operators
using compact linear operators.
More precisely,
a Fredholm operator can be represented by means of
a \textit{left} and \textit{right parametrix},
respectively.

\subparagraph{}
Let $\linA \: \banX \xto \banY$ be a continuous linear operator.
Then a linear mapping $P_l \: \banY \xto \banX$
is called a \textit{left parametrix} of $\linA$
if $\linA \circ P_l \eq \1_{\banX} {+} \compA_l$
for some compact operator $\compA_l \xeof \BO{\banX}$.
In a similar way,
a linear mapping $P_r \: \banY \xto \banY$
is said to be a \textit{right parametrix} of $\linA$,
if $P_r \circ \linA \eq \1_{\banY} {+} \compA_r$
for some compact operator $\compA_r \xeof \BO{\banY}$.

\subparagraph{} %- characterization of Fredholm operators
A linear operator $\linA$ is Fredholm if and only if
it owns a left and right parametrix,
where the involved compact operators $\compA_l$ and $\compA_r$
are even of \textit{finite-rank},
that is,
they have finite-dimensional ranges.

\paragraph{} %- Fredholm's alternative theorem
%------------------------------------------------------------------------------%
In the following we cite the \textit{Fredholm} \textit{alternative theorem}
as the outstanding result in Fredholm theory
{\wrt} the solvability of the linear equation $\linA x \eq y$.
Remember the \textit{rank-nullity theorem} known from linear algebra,
claiming that if a linear space $\vecV$ is finite-dimensional,
then a linear mapping $\linA \: \vecV \xto \vecV$ is surjective
\iff it is injective.

\subparagraph{}
The alternative theorem is a considerable generalization
of the rank-nullity theorem to arbitrary Banach spaces and Fredholm operators
with vanishing index:

\begin{theorem}\label{FREDHOLM_ALTERNATIVE_THEOREM}
Let $\banX,\,\banY$ be Banach spaces and
the continuous linear operator $\linA \: \banX \xto \banY$ be Fredholm
with index $\chi(\linA) \eq 0$.

\subparagraph{}
According to the solvability of the linear equation $\linA \, x \eq y$,
exactly one of the following cases occurs:

\begin{itemize}[leftmargin=9mm]
\item[(a)]
$\linA \, x \eq y$ is uniquely solvable for all $y \xeof \banX$,
so $\linA$ is bijective,

\item[(b)]
$\linA \, x \eq y$ is not for all $y \xeof \banX$ solvable,
and $\linA$ is neither injective nor surjective in this case.
However,
if the equation is solvable for an element $y_0 \xeof \banX$,
then the set of solutions forms a finite-dimensional subspace of $\banX$.
Moreover,
the equation $\linA \, x \eq y$ owns a solution $x \xeof \banX$
\iff
$y \xeof \,^{\bot}\kernel{\linA^*}$.
\end{itemize}
\end{theorem}

\subsection{} %-- essential spectrum
The various concepts of semi-Fredholm operators can be used to characterize
the \textit{essential spectrum} $\sigma_{ess,0}(\linA)$
of a closed operator $\linA$ acting in a complex Banach space $\banX$:
a number $\lambda \xeof \C$ belongs to $\sigma_{ess,0}(\linA)$
\iff $\linA {-} \lambda$ is \textit{not} semi-Fredholm.
But there are further approaches to define the notion of essential spectrum:

\begin{itemize}[leftmargin=6mm]
\item[\,] 
$\sigma_{ess,1}(\linA) := \{\lambda \xeof \C:
\linA {-} \lambda \mbox{ is not upper semi-Fredholm} \}$,

\item[\,] 
$\sigma_{ess,2}(\linA) := \{\lambda \xeof \C:
\linA {-} \lambda \mbox{ is not lower semi-Fredholm} \}$,

\item[\,] 
$\sigma_{ess,3}(\linA) := \{\lambda \xeof \C:
\linA {-} \lambda \mbox{ is not emi-Fredholm} \}$,

\item[\,] 
$\sigma_{ess,4}(\linA) := \{\lambda \xeof \C:
\linA {-} \lambda \mbox{ is Fredholm} \}$,

\item[\,] 
$\sigma_{ess,5}(\linA) := \{\lambda \xeof \C:
\linA {-} \lambda \mbox{ is not Fredholm with } \chi(\linA {-} \lambda) \eq 0\}$,

\item[\,] 
$\sigma_{ess,6}(\linA) := \{\lambda \xeof \sigma_{ess,5}(\linA) \:
\lambda \mbox{ is an isolated point in } \sigma(\linA)\}$.
\end{itemize}

\subparagraph{} %- properties of the essential spectrum
The sets $\sigma_{ess,4}(\linA)$ and $\sigma_{ess,5}(\linA)$
are called the \textit{Wolf spectrum} and the \textit{Schechter spectrum},
respectively,
whereas the set $\sigma_{ess,6}(\linA)$ is denoted
the \textit{Browder spectrum}.

\subparagraph{}
All sets $\sigma_{ess,k}(\linA)$ are closed subsets of $\sigma(\linA)$,
and by definition we have the chain of inclusions 
\[
\sigma_{ess,3}(\linA) \xsubset
\sigma_{ess,1}(\linA) \cup
\sigma_{ess,2}(\linA) \xsubset
\sigma_{ess,4}(\linA) \xsubset
\sigma_{ess,5}(\linA) \xsubset
\sigma_{ess,6}(\linA).
\]

\subparagraph{}
In the Hilbert space setting,
however
all sets $\sigma_{ess,k}(\linA)$ coincide and form the largest subset
$\sigma_{ess}(\linA) \xsubset \sigma(\linA)$ such that
for every compact linear operator $\compA \xeof \BO{\banX}$ 
\[
\sigma_{ess}(\linA + \compA) \= \sigma_{ess}(\linA).
\]

%------------------------------------------------------------------------------%
\section{Symmetry and Self-Adjointness}                                        %
%------------------------------------------------------------------------------%
%-- 
%-- 6. SYMMETRY AND SELF-ADJOINTNESS
%--    -----------------------------
This section serves as an overview on symmetric and self-adjoint linear
operators acting in a Hilbert space.
Let us begin with some remarks on the numerical range.

\subparagraph{} %-- numerical range and external field of regularity
Besides the spectrum $\sigma(\linA)$,
the \textit{numerical range} $\Theta(\linA)$ plays an important role
in the general theory of a Hilbert space operator $\linA$ 
and is defined by
\[
\Theta(\linA) :=
\{ \scalar{\linA}{x} \: x \xeof \domain{\linA}, \norm{x} \eq 1 \}.
\]
It is a well-known fact,
given by the \textit{Toeplitz-Hausdorff theorem},
that $\Theta(\linA)$ is a convex set within the complex plane.
There are relationships between the numerical range and the some kinds
of spectra of the operator $\linA$.
Indeed,
we have $\sigma_p(\linA) \xsubset \Theta(\linA)$,
$\sigma_{ap}(\linA) \xsubset \overline{\Theta(\linA)}$
and even $\sigma(\linA) \xsubset \Theta(\linA)$ for $\linA \xeof \BO{\hilH}$,
(see \cite{Ha}, theorem A.23 on p.\,266).

\paragraph{} %- external field of regularity
Let $\Delta(\linA) := \C \setminus \overline{\Theta(\linA)}$
be the \textit{external field of regularity} of $\linA$.
If $\Delta(\linA) \neq \emptyset$,
then it consists of one or two connected components.
More precisely,
if $\Delta(\linA)$ has two connected components,
then the numerical range is either a strip or a ray,
otherwise it is a half-strip or a half-ray.

\subparagraph{} %- field of regularity
For $\zeta \xeof \Delta(\linA)$ we of course have
$\scalar{\linA x}{x} \neq \zeta$
for $x \xeof \domain{\linA}$ and $\norm{x} \eq 1$,
and the following chain of inequalities
\[
0 \, < \, \mbox{dist}(\zeta,\linA) \, \leq \,
\abs{\scalar{\linA x}{x} - \zeta} \=
\abs{\scalar{(\linA {-} \zeta)x}{x}} \, \leq \,
\norm{(\linA {-}\zeta) x}
\]
is valid,
where $\mbox{dist}(\zeta,\linA)$ denotes the distance
between the singleton $\{\zeta\} \xsubset \Delta(\linA)$
and the set $\overline{\Theta(\linA)}$,
which in fact is positive,
since these sets are closed and have empty intersection.
This implies that for $\zeta \xeof \rho(\linA) \cap \Delta(\linA)$ we have
\[
\norm{R(\linA,\zeta)} \, \leq \, \frac{1}{\mbox{dist}(\zeta,\Theta(\linA))}.
\]
For $\zeta \xeof \Delta(\linA)$
the operator $\linA {-} \zeta$ is bounded below and
$\Delta(\linA)$ is a subset of the \textit{field of regularity}
$\Pi(\linA) := \{\zeta \xeof \C \: \linA {-} \zeta \mbox{ is bounded below}\}$.

\subsection{} %-- symmetric operators
%------------------------------------------------------------------------------%
A linear operator $\linA$ acting in $\hilH$
and with domain of definition $\domain{\linA}$ dense in $\hilH$ 
is called \textit{symmetric},
if the adjoint $\linA^*$ forms an extension of $\linA$,
and this property is equivalent to
\[
\scalar{\linA x}{y} \= \scalar{x}{\linA y} \qfa x,y \in \domain{\linA}.
\]
But if $\domain{\linA}$ is \textit{not} dense in $\hilH$,
it is still possible to define potential symmetry of $\linA$:
in this case, $\linA$ is called \textit{symmetric},
if $\Theta(\linA) \xsubset \R$.
It can be shown
that if $\linA$ is densely defined,
then the both mentioned concepts of symmetry are equivalent. 

\subparagraph{} %- half-bounded linear operators
Finally,
a symmetric linear operator $\linA$ acting in $\hilH$
is called \textit{half-bounded},
if $\Theta(\linA) \xsubset ({-}\infty,a)$
or $\Theta(\linA) \xsubset (b,\infty)$ holds.

\paragraph{} %- properties of symmetric operators
%------------------------------------------------------------------------------%
Densely defined and symmetric linear operators distinguish
by a couple of properties.
Indeed, if $\linA$ is symmetric, then

\begin{itemize}[leftmargin=9mm]
\item[(a)]
$\linA$ is closable,
since $\linA^*$ forms a closed extension of $\linA$;
\item[(b)]
the closure $\overline{\linA}$ is symmetric, too;
\item[(c)]
all eigenvalues of $\linA$ are real;
\item[(d)]
eigenvectors for different eigenvalues of $\linA$
are \textit{orthogonal} to each other.
\end{itemize}

\subparagraph{} %- A+i and A-i are bounded below
A useful property of symmetric closed operators is given by the fact
that the operators $\linA {-} \zeta$ are bounded below for $\im \zeta \neq 0$,
hence they are injective and own closed ranges,
so they are semi-Fredholm operators with $\alpha(\linA {-} \zeta) \eq 0$.
Furthermore,
the codimensions $\beta(\linA {-} \zeta)$ of the ranges $\range{\linA {-} \zeta}$
are constant in the upper and lower half-plane of $\C$,
respectively,
which is a consequence of the stability theorem for semi-Fredholm operators.
The values $d_{\pm} := d(\pm i) \xeof \N_0 \cup \{\infty\}$
are called the \textit{defect indices} of $\linA$.

\subsection{} %-- spectrum of symmetric operators
Now the decisional question is
whether the operators $\linA \pm \zeta$ are surjective and
the defect indices $d_{\pm}$ of $\linA$ are zero,
respectively.
To the end we can distinguish four possible cases:

\begin{itemize}[leftmargin=4mm]
\item[1.]
$d_+ \neq 0$ and $d_- \neq 0$.
The operators $\linA {+} i$ and $\linA {-} i$ are not surjective,
hence the upper and lower complex half-plane
belong to the spectrum $\sigma(\linA)$.
Because the spectrum is closed,
the real line also belongs to $\sigma(\linA)$,
and in conclusion we have $\sigma(\linA) \eq \C$.
In addition,
for $\im \zeta \neq 0$ we have $\zeta \xeof \sigma_r(\linA)$.

\item[2.]
$d_+ \eq 0$ and $d_- \neq 0$.
The operator $\linA {+} i$ is surjective,
but the operator $\linA {-} i$ is not so.
Hence the lower complex half-plane and the real line belong to $\sigma(\linA)$,
whereas the upper complex half-plane belongs to $\rho(\linA)$.

\item[3.]
$d_+ \neq 0$ and $d_- \eq 0$.
The operator $\linA {-} i$ is surjective,
but the operator $\linA {+} i$ is not so.
Hence the upper complex half-plane and the real line belong to $\sigma(\linA)$,
whereas the lower complex half-plane belongs to $\rho(\linA)$.

\item[4.]
$d_+ \eq 0$ and $d_- \eq 0$.
The operators $\linA {+} i$ and $\linA {-} i$ are surjective,
hence the upper and lower complex half-plane belong
to the resolvent set $\rho(\linA)$.
The spectrum $\sigma(\linA)$ forms a subset of the real line in this case.
\end{itemize}

\subsection{} %-- self-adjoint operators
%------------------------------------------------------------------------------%
It should be notified that the adjoint $\linA^*$ of a symmetric linear operator
$\linA$ is not necessarily symmetric again.
If $\linA$ is symmetric,
then every other linear operator $\linS$ having the property
$\linA \prec \linS$ is called a \textit{symmetric extension} of $\linA$.
Notice 
that $\linA \prec \linS$ always implies $\linS^* \prec \linA^*$,
which is also true for non-symmetric operators $\linA$ and $\linS$.

\subparagraph{}
Now,
let $\linS_1,\ldots,\linS_k$ be a finite sequence
of symmetric extensions of $\linA$,
where $\linS_{i+1}$ extends $\linS_i$ for $i \eq 1,\ldots,k-1$.
Since $\linS_k \prec \linS_k^*$ and
$\linS_{i+1}^* \prec \linS_i^*$ for $i \eq 1,\ldots,k-1$,
we obtain by concatenating all these extensions the chain
\[
\linA \, \prec \, \linS_1
\, \prec \,
\linS_2
\, \prec \,
\ldots
\, \prec \,
\linS_k
\, \prec \,
\linS_k^*
\, \prec \,
\ldots
\, \prec \,
\linS_2^*
\, \prec \,
\linS_1^*
\, \prec \,
\linA^*.
\]
Thus one may hope to find a \textit{maximal} symmetric extension
$\linS_{sa}$ of $\linA$ such that
\[
\linA \, \prec \, \linS_{sa} \= \linS_{sa}^* \, \prec \, \linA^*
\]
holds,
and in this case
the operator $\linS_{sa}$ cannot be extended to a larger symmetric operator.
This expectation suggests to introduce the notion of \textit{self-adjointness}:
a symmetric linear operator $\linA$ is called \textit{self-adjoint},
if $\linA^* \prec \linA$,
or equivalently,
if $\linA \eq \linA^*$.

\paragraph{} %- Fredholm property of self-adjoint linear operators
%------------------------------------------------------------------------------%
Let us mention
that if a self-adjoint linear operator $\linA$ is Fredholm,
then $\chi(\linA) \eq 0$,
which follows from the equality $\chi(\linA) \eq - \chi(\linA^*)$.
Remember from section 5
that the various definitions of the essential spectrum of $\linA$ 
are equivalent for any linear operator $\linA$ acting in a Hilbert space.
Hence we have
\[
\sigma_{ess,1}(\linA)  \eq
\sigma_{ess,2}(\linA)  \eq
\sigma_{ess,3}(\linA)  \eq
\sigma_{ess,4}(\linA)  \eq
\sigma_{ess,5}(\linA)  \eq
\sigma_{ess,6}(\linA),
\]
which gives rise to define $\sigma_{ess}(\linA) := \sigma_{ess,k}$
for any $1 \leq k \leq 6$.

\subparagraph{} %- Weyl sequence
Recall that the discrete spectrum
$\sigma_d(\linA) := \sigma(\linA) \setminus \sigma_{ess}(\linA)$
is nothing but the set of isolated eigenvalues of $\linA$
having finite algebraic multiplicity,
so $\kernel{\linA {-} \lambda}$ is finite-dimensional
for every $\lambda \xeof \sigma_d(\linA)$.
On the other hand,
any value $\lambda$ belongs to $\sigma_{ess}(\linA)$,
if there is a sequence $\{x_n\}_{n \xeof \N}$ of orthonormal elements in $\hilH$
such that $\norm{\linA x_n - \lambda x_n} \xto 0$ for $n \xto \infty$,
which is denoted a \textit{Weyl} \textit{sequence} {\wrt} $\lambda$.

\paragraph{} %- range condition
%------------------------------------------------------------------------------%
To ensure self-adjointess of a symmetric linear operator $\linA$,
the \textit{range condition} reflects the forth case in our considerations
on the defect indices of $\linA$
and serves as the main tool to verify this pleasant property:

\begin{theorem}\label{RANGE_CONDITION}
A densely defined, closed and symmetric linear operator $\linA$
acting in $\hilH$ is self-adjoint iff there is $\zeta \xeof \C$ such that
\[
\range{\linA + \zeta} \= \range{\linA + \overline{\zeta}} \= \hilH,
\]
or equivalently,
if the defect indices of $\linA \pm \zeta$ fulfill $d(\pm \zeta) \eq 0$.
\end{theorem}

\subparagraph{}
If it is even possible to find $\zeta \xeof \R$ fulfilling the range condition 
$\range{\linA {+}\zeta} \eq \hilH$,
then $\linA$ also proves to be self-adjoint.

\paragraph{} %- essentially self-adjointness
%------------------------------------------------------------------------------%
A symmetric linear operator is denoted \textit{essentially self-adjoint},
if its closure proves to be self-adjoint.
There is a comparable criterion for essentially self-adjointness available:

\begin{theorem}\label{RANGE_CONDITION_2}
A densely defined and symmetric linear operator $\linA$
acting in $\hilH$ is essentially self-adjoint
iff there is $\zeta \xeof \C$ such that
\[
\overline{\range{\linA + i}} \=
\overline{\range{\linA - i}} \= \hilH.
\]
\end{theorem}

\subparagraph{}
Essentially self-adjointness can also be characterized by means of the adjoint:
a densely defined and symmetric linear operator $\linA$ is essentially
self-adjoint iff the adjoint $\linA^*$ is symmetric, too.

\subsection{} %-- perturbation theory of self-adjoint operators
%------------------------------------------------------------------------------%
Let us now deal with operators of the form $\linA {+}\linB$,
where $\linA$ is self-adjoint and $\linB$ is symmetric
with domain of definition $\domain{\linB} \supset \domain{\linA}$,
so $\linA {+} \linB$ is well-defined
with $\domain{\linA {+} \linB} := \domain{\linA}$.
It is our aim to find
conditions on the operator $\linB$
to obtain self-adjointness of $\linA {+} \linB$.

\subparagraph{} %- A-boundedness
The following concept proves to be useful:
a symmetric linear operator $\linB \: \domain{\linB} \xsubset \hilH \xto \hilH$
with $\domain{\linA} \xsubset \domain{\linB}$
is called $\linA$\textit{-bounded},
if there are constants $\theta, c \geq 0$ such that
\begin{equation}\label{A_BOUNDEDNESS}
\norm{\linB \,x} \, \leq \, \theta \norm{\linA \,x} \, + \, c \norm{x}
\qfa x \xeof \domain{\linA}.
\end{equation}
In order to ensure $\linA$-boundedness of $\linB$,
it is sufficient to ensure \eqref{A_BOUNDEDNESS}
only for the elements in a core $\linA_0$ of $\linA$.

\subparagraph{} %- A-bound
The infimum of all numbers $\theta$,
for which there is $c$ such that \eqref{A_BOUNDEDNESS} holds,
is called the $\linA$\textit{-bound} of $\linB$.
Notice that if $\linB$ is a bounded operator,
then it is $\linA$-bounded with $\linA$-bound $\theta <1$ for any $\linA$.
In addition,
the validness of the estimate
\[
\norm{\linB \,x}^2 \, \leq \,
\theta^2 \norm{\linA \,x}^2 \, + \, c^2 \norm{x}^2
\qfa x \xeof \domain{\linA}.
\]
is a sufficient condition for \eqref{A_BOUNDEDNESS} to be true.

\paragraph{} %- Kato-Rellich theorem
%------------------------------------------------------------------------------%
The \textsc{Kato-Rellich} \textit{theorem}
provides a criterion on self-adjoint\-ness of the operator $\linA {+} \linB$:

\begin{theorem}\label{KATO-RELLICH_THEOREM_1}
Let the linear operator $\linA$ be self-adjoint and 
the linear operator $\linB$ with $\domain{\linB} \supset \domain{\linA}$ 
be symmetric and $\linA$-bounded with $\theta < 1$.
Then $\linA {+} \linB$ with domain of definition 
$\domain{\linA {+} \linB} \eq \domain{\linA}$ is self-adjoint.
\end{theorem}

\subparagraph{} %- Kato-Rellich theorem for essentially self-adjoint operators
A similar result is obtained for the case,
where \textit{essentially self-adjointness} of a linear operator $\linA$
should be preserved by adding
a symmetric and $\linA$-bounded perturbation $\linB$:

\begin{theorem}\label{KATO-RELLICH_THEOREM_2}
Let $\linA$ be essentially self-adjoint and assume that
$\linB$ is symmetric and $\linA$-bounded with bound $\theta < 1$.
Then the operator $\linA {+} \linB$ is essentially self-adjoint,
where
$\overline{\linA {+} \linB} \eq \overline{\linA} + \overline{\linB}$.
\end{theorem}

\subsection{} %-- multiplication operators
%------------------------------------------------------------------------------%
We finally want to mention an important class of symmetric and self-adjoint
linear operators.
Remember from integration theory
that for a given measure space $(X,\salgB,\mu)$
the linear space $L^2(X)$ of measurable and square-integrable functions
$f\:X \xto \C$ forms a Hilbert space with inner product
\[
\scalar{f}{g} \eqdef \int_X f \, \overline{g} \, d\mu.
\]
If $m \: X \xto \C$ is a measurable function,
then the linear mapping $\linM_m$ defined by
\[
\begin{cases}[1.2]
\quad \domain{\linM_m} & := \{f \xeof L^2(X) \: \; m \, f \in L^2(X)\}, \\
\quad (\linM_m f)(x)   & := m(x) \cdot f(x)   \;
                            \for x \xeof X \and f \xeof \domain{\linM_m} \\
\end{cases}
\]
is called a \textit{multiplication operator} with \textit{generator} $m$.
It turns out that the mapping $\linM_m$ is symmetric
iff
$m$ is real-valued,
and since the domains of definition
of the operators $\linM_m$ and $\linM_m^*$ coincide,
each multiplication operator with real-valued generator is self-adjoint.

%------------------------------------------------------------------------------%
\section{Spectral Theory of Self-Adjoint Operators}                            %
%------------------------------------------------------------------------------%
%-- 
%-- 7. SPECTRAL THEORY OF SELF-ADJOINT OPERATORS
%--    -----------------------------------------
The spectral theory of self-adjoint linear operators is one of the pearls
in the mathematical area of functional analysis
and has important applications in physics and engineering science
with a multiplicity of useful results.

\subparagraph{} %-- the spectrum of self-adjoint operators
First of all,
we may characterize the self-adjointness of a linear operator
acting in a Hilbert space by means of its spectrum,
which is an immediate consequence of the range condition \ref{RANGE_CONDITION}:

\begin{theorem}
A linear operator $\linA$ is self-adjoint
iff $\sigma(\linA) \xsubset \R$.
\end{theorem}

\subparagraph{} %- resolvent estimate for self-adjoint operators
A further result is given by the \textit{resolvent estimate}:
if $\linA$ is self-adjoint,
then for $\im \zeta \neq 0$ the resolvent estimation
\[
\norm{R(\linA,\zeta)} \, \leq \, \abs{\im \zeta} ^{-1}
\]
holds,
which follows from the inequality 
$\abs{\im \zeta} \leq \mbox{dist}(\zeta,\sigma(\linA))$
and the general estimate \eqref{RESOLVENT_ESTIMATE_I}.

\subparagraph{} %- residual spectrum of a self-adjoint operator is empty
Moreover,
if $\linA$ is self-adjoint and $\lambda \xeof \sigma(\linA)$,
then $\lambda$ is either an eigenvalue of $\linA$
or $\range{\linA {-} \lambda}$ is dense in $\hilH$,
but not the whole of $\hilH$.
Hence in the first case the value $\lambda$
belongs to the {point spectrum} $\sigma_p(\linA)$ of $\linA$,
whereas the second case means that
$\lambda$ belongs to the {continuous spectrum} $\sigma_c(\linA)$.
This means to say the {residual spectrum} $\sigma_r(\linA)$ is empty.

\subsection{} %-- spectral resolutions
%------------------------------------------------------------------------------%
The most important theorem {\wrt} self-adjoint linear operators,
however,
is the celebrated \textit{spectral theorem}.

\subparagraph{} %- spectral resolutions
Recall that a \textit{projection} in $\hilH$ is a symmetric linear operator
$\linP$ such that $\linP {\circ} \linP \eq \linP$.
Then a \textit{spectral resolution} is a one-parameter family
$\{\linP_{\lambda}\}_{\lambda \xeof \R}$
of linear operators $\linP_{\lambda} \: \hilH \xto \hilH$
such that

\begin{itemize}[leftmargin=12mm]
\item[(R-1)]
$\linP_{\lambda}$ is a projection for each $\lambda \xeof \R$;
\item[(R-2)]
for each $x \xeof \hilH$
the mapping $\lambda \mapsto \linP_{\lambda} x$ is right-continuous;
\item[(R-3)]
$\lambda {\,<\,} \mu$ implies $\linP_{\lambda} \leq \linP_{\mu}$,
so $\scalar{\linP_{\lambda} x}{x} \leq \scalar{\linP_{\mu} x}{x}$
for $x \xeof \hilH$;
\item[(R-4)]
$\lim_{\lambda \xto -\infty} \linP_{\lambda} \, x \eq 0$
for all $x \xeof \hilH$;
\item[(R-5)]
$\lim_{\lambda \xto  \infty} \linP_{\lambda} \, x \eq x$
for all $x \xeof \hilH$.
\end{itemize}

\subparagraph{}
Notice also that each projection $\linP_{\lambda}$ is related
to a closed linear subspace $\setA_{\lambda} \xsubset \hilH$.
Moreover,
the difference $\linP_{\mu} - \linP_{\lambda}$
for $\mu \geq \lambda$ also constitutes a projection 
related to the closed subspace $\setA_{\mu} \ominus \setA_{\lambda}$.

\subsection{} %-- spectral operators
%------------------------------------------------------------------------------%
A spectral resolution can be used to define an operator $\linA$ in $\hilH$
by taking the spectral integral
\begin{equation}\label{SPECTRAL_INTEGRAL}
\linA x \= \int_{\,\R} \lambda \, d\linP_{\lambda}x,
\end{equation}
where the domain of definition $\domain{\linA}$
consists of those elements $x \xeof \hilH$,
for which this improper Riemann-Stieltjes integral converges.
The mapping $\linA$ is said to be the \textit{spectral operator}
associated with the resolution $\{\linP_{\lambda}\}_{\lambda \xeof \R}$
and proves to be linear and even self-adjoint.

\paragraph{} %-- the spectral theorem
%------------------------------------------------------------------------------%
The \textit{spectral theorem} for self-adjoint operators claims
that the inverse statement also holds:

\begin{theorem}\label{SPECTRAL_THEOREM}
Each self-adjoint linear operator $\linA$
owns a spectral resolution $\{\linP_{\lambda}\}_{\lambda \xeof \R}$
such that
\[
\domain{\linA} \=
\{\, x \xeof \hilH:
    \int_{\,\R}
\lambda^2 \, d \scalar{\linP_{\lambda} \,x}{x} < \infty\,\}
\]
and
\[
\scalar{\linA \, x}{y} \=
\int_{\,\R} \lambda \, d\scalar{\linP_{\lambda} \,x}{y}
\quad \mbox{for } x \xeof \domain{\linA} \mbox{ and } y \xeof \hilH.
\]
\end{theorem}

\subparagraph{}
In summary,
self-adjoint linear operators are in a one-to-one relationship
to spectral operators defined by formula \eqref{SPECTRAL_INTEGRAL}.

\paragraph*{} %- relationship between the spectral resolution and the spectrum
%------------------------------------------------------------------------------%
We next want to point out the relationship between
the shape of the spectrum $\sigma(\linA)$ and the spectral resolution of $\linA$.
The following theorem shows
that the regular values and the different kinds of spectral values of $\linA$
are determined by the spectral resolution:

\begin{theorem}\label{SPECTRUM_RESOLUTION_THEOREM}
Let $\linA$ be self-adjoint
and $\setA_{\mu}$ be the closed subspace in $\hilH$
associated with the projection $\linP_{\mu}$ for some $\mu \xeof \R$.
Then

\begin{itemize}[leftmargin=9mm]
\item[1.]
$\mu \xeof \R$ is a \textit{regular value} of $\linA$
iff
there is an open interval $(\alpha,\beta) \xsubset \R$
such that $\mu \xeof (\alpha,\beta)$
and $\linP_{\nu} \eq \linP_{\mu}$ for each $\nu \xeof (\alpha,\beta)$.

\item[2.]
$\mu \xeof \R$ belongs to the \textit{point spectrum} $\sigma_p(\linA)$
iff
\[
0 < \dim \; (\setA_{\mu+0} \ominus \setA_{\mu}) < \infty.
\]
%Besides for some number $\epsilon > 0$
%the subspace
%$\setA_{\mu+\epsilon} \ominus \setA_{\mu-\epsilon}$ is finite-dimensional.
%Hence $\mu$ is not only an eigenvalue of $\linA$,
%but also belongs to the set $\sigma_d(\linA)$.

\item[3.]
$\mu \xeof \R$ belongs to the \textit{continuous spectrum} $\sigma_c(\linA)$
iff
for each $\epsilon > 0$ the subspace
$\setA_{\mu+\epsilon} \ominus \setA_{\mu-\epsilon}$ is infinite-dimensional.
\end{itemize}
\end{theorem}

\subsection{} %-- the functional calculus for self-adjoint operators
%------------------------------------------------------------------------------%
The spectral theorem forms the base
for the so-called \textit{functional calculus} for self-adjoint operators:
if a given function $f \: \R \xto \C$ is sufficiently regular
(for example,
if $f$ is piecewise continuous and the points of discontinuity
do not have any accumulation point in the real line),
then it is possible to define the operator $f(\linA)$
by means of the spectral resolution
$\{\linP_{\lambda}\}_{\lambda \xeof \R}$ of $\linA$:
\[
f(\linA) \, x \eqdef \int_{\,\R} f(\lambda) \, d\linP_{\lambda}x,
\]
where the domain of definiton of the operator $f(\linA)$ is defined by
\[
\domain{f(\linA)} \eqdef
\{x \xeof \hilH \:
\int_{\R} \abs{f(\lambda)} ^2 \, d\linP_{\lambda}x < \infty \}.
\]
Notice that the operator $f(\linA)$ is self-adjoint iff $f$ is real-valued.

\subparagraph{} %- functional calculus for unbounded self-adjoint operators
The functional calculus for unbounded self-adjoint operators owns
the following properties:

\begin{itemize}[leftmargin=9mm]
\item[1.] $f(\alpha \linA) \eq \alpha f(\linA)$ for $\alpha \xeof \C$.
\item[2.] $f(\linA) + g(\linA) \prec (f+g)\linA$.
\item[3.] $f(\linA) g(\linA) \prec f(g(\linA))$.
\item[4.] $f(\linA)^* \eq f(\linA^*)$.
\item[5.] $\sigma(f(\linA)) \eq f(\sigma(\linA))$
          (\textit{spectral mapping theorem}).
\end{itemize}

\paragraph{} %- functional calculus for self-adjoint positive operators
%------------------------------------------------------------------------------%
We are especially interested in the case,
where the self-adjoint operator $\linA$ is positive and
$f(\lambda) := 0$ for $t<0$ and $e^{-\lambda t}$ for $t \geq 0$,
in which the family $\{e^{-t\linA}\}_{t \geq 0}$ of bounded linear operators
is shown to constitute an analytic semigroup, see section 11).
But also other operators $f(\linA)$ are meaningful in applications,
which for example are given by the functions
$x \mapsto \sqrt{x}$,
$x \mapsto \abs{x} $,
$x \mapsto \cos{x}$,
$x \mapsto \sin{x}$
and are denoted as
$\sqrt{\linA}$,
$\abs{\linA} $,
$\cos{\linA}$,
$\sin{\linA}$.

\subparagraph{}
Especially the resolvents $R(\linA,\zeta) := (\linA-\zeta)^{-1}$
for $\zeta \xeof \rho(\linA)$ of a self-adjoint linear operator $\linA$
can be represented by means of the spectral resolution,
that is
\[
R(\linA,\zeta) \, x \= \int_{-\infty}^{\infty}
\frac{1}{\lambda-\zeta} \, dP_{\lambda} x,
\quad (x \in \hilH, \, \im \zeta \neq 0).
\]

\subparagraph{} %- Stone's formula
On the other hand,
\textsc{Stone}\textit{'s formula} provides a constructive method
to describe the spectral resolution of $\linA$ in an interval $[a,b]$
by means of the resolvents $R(\linA,\zeta)$.
Let $x \xeof \hilH$ be arbitrary, then
\[
\frac{1}{2} \{P[a,b] + P(a,b)\} \, x
=
\lim_{\epsilon \downarrow 0} \, \frac{1}{2 \pi i}
\int_a^b
\{R(\linA,t + i \epsilon)\,x -
  R(\linA,t - i \epsilon)\,x \} \,
  dt.
\]

\subsection{} %-- evolution equations governed by self-adjoint operators
%------------------------------------------------------------------------------%
The presented functional calculus provides applications
to initial-value problems for evolution equations
governed by self-adjoint and positive linear operators.
In the rest of this section we shall elaborate the usage of this method.

\subparagraph{} %- differentiable vector-valued functions
First of all let us clarify
what is meant by a differentiable vector-valued function.
If $\banX$ is an arbitrary Banach space and $[a,b] \xsubset \R$,
then a vector-valued function $x \: [a,b] \xto \banX$
is said to be \textit{classical differentiable} in $t_0 \xeof (a,b)$,
if there is an element $\dot{x}(t_0) \xeof \banX$ such that
\[
\lim_{\,t \xto t_0} \,
\norm{ \, \frac{x(t)-x(t_0)}{t-t_0}\, - \dot{x}(t_0) \, } \= 0.
\]
Moreover,
the function $x$ is called {classical differentiable in} $(a,b)$,
if $x$ is so in each $t \xeof (a,b)$.
The linear space consisting of functions having \textit{continuous}
classical derivatives in $(a,b)$ is denoted $C^1(a,b;\banX)$.

\subparagraph{}
The spaces $C^k(a,b;\banX)$ ($k \xeof \N$) are defined recursively as usual,
and especially the space $C^{\infty}(a,b;\banX)$ is defined
by taking the intersection of all spaces $C^k(a,b;\banX)$ for $k \xeof \N_0$.

\paragraph{} % initial-value problems governed by self-adjoint linear operators
%------------------------------------------------------------------------------%
The initial-value problems we want to consider in this section
are the homogeneous Cauchy problems of first and second order
and the abstract {\SCHROEDINGER} equation.
Let us notice in advance that they can be solved as well by semigroup methods,
see also the sections 9,\,10 and 11.

\begin{itemize}[leftmargin=12pt,itemsep=2mm,topsep=1mm]
\item[1.]
The \textit{homogeneous Cauchy problem of first order}
is the initial-value problem given by
\begin{equation}\label{HOMOGENEOUS_CAUCHY_PROBLEM_I}
\begin{cases}[1.2]
\quad \dot{x}(t) \+ \linA \, x(t) \= 0,\\ 
\quad x(0) \= x_0 \xeof \hilH,\\
\end{cases}
\end{equation}
where $\linA$ is self-adjoint and positive,
that is,
$\scalar{\linA x}{x} \geq m \norm{x}^2$
for some $m \geq 0$ and all $x \xeof \domain{\linA}$.

\subparagraph{} %- functional calculus approach
The task is to find \textit{vector-valued} solutions $x \: [0,\T) \xto \hilH$
of the equation $\dot{x} + \linA x \eq \0$ 
satisfying $x \xeof C^1((0,\T);\hilH) \cap C([0,\T);\hilH)$ and
the initial value condition $x(0) \eq x_0$.
If $\linA$ is positive,
then any solution $x\:[0,\T) \xto \hilH$
of problem \eqref{HOMOGENEOUS_CAUCHY_PROBLEM_I} satisfies the estimate
\begin{equation}\label{ESTIMATION_FOR_NONNEGATIVE_OPERATOR}
\norm{x(t)} \, \leq \, \norm{x(0)}
\quad (t > 0),
\end{equation}
which is a consequence of the inequality
$\norm{x(t)}^2 \eq - 2 \, \scalar{\linA x}{x} \leq 0$
and therefore causes the solution to be unique.

\subparagraph{}
To verify the existence of a solution for \eqref{HOMOGENEOUS_CAUCHY_PROBLEM_I},
let $(P_{\lambda})_{\lambda \in \R}$ be the spectral resolution
of the self-adjoint operator $\linA$ given by the spectral theorem.
We notice $P_{\lambda} \eq \0$ for $\lambda < 0$,
since $\linA$ is positive.
Then the functional calculus applied to the continuous function
$t \mapsto e^{-t}$ gives rise to the spectral operator
\[
\linL_t(x)
\eqdef
e^{-t \linA} x
\=
\int_{\,0}^{\,\infty} e^{-\lambda t} \, dP_{\lambda} \, x
\quad \quad (t \geq 0, \, x \in \hilH).
\]
We apparently have $\linL_0 \eq \1_{\banX}$,
and \eqref{ESTIMATION_FOR_NONNEGATIVE_OPERATOR} implies
$\norm{\linL_t} \leq 1$ for $t > 0$.

\subparagraph{} %- existence of a solution for the homogeneous Cauchy problem
In conclusion,
the existence of the solution of \eqref{HOMOGENEOUS_CAUCHY_PROBLEM_I}
is ensured by the following

\begin{theorem}\label{EXISTENCE_THEOREM_FOR_ABSTRACT_HEAT_EQUATION}
The function $x\:t \mapsto \linL_t(x_0)$ is continuous in $[0,\infty)$,
continuously differentiable in $(0,\infty)$
and forms a solution of \eqref{HOMOGENEOUS_CAUCHY_PROBLEM_I}.
\end{theorem}

\item[2.] %- second order equations
The \textit{homogeneous Cauchy problem of second order} consists of solving
the initial value problem
\begin{equation}\label{HOMOGENEOUS_CAUCHY_PROBLEM_II}
\begin{cases}[1.2]
\quad \ddot{x}(t) \+ \linA \, x(t) \= \0,\\
\quad x(0)  \= x_0,
\quad \dot{x}(0) \= x_1,\\
\end{cases}
\end{equation}
where the operator $\linA$ again is self-adjoint and positive.
We also suppose $x_0 \xeof \domain{\linA}$,
but the element $x_1$ may belong to any other linear subspace of $\banX$.
A function $x$ is called a classical solution of this problem,
if $x \xeof C([0,\T],\hilH) \cap C^2((0,\T),\hilH)$
and the initial conditions are satisfied.

\paragraph{} %- functional calculus approach for the second-order Cauchy problem
%------------------------------------------------------------------------------%
The functional calculus for self-adjoint linear operators
may also be applied to this problem.
Preliminary,
each solution $x$ of \eqref{HOMOGENEOUS_CAUCHY_PROBLEM_II} satisfies
the \textit{energy conservation law}
\[
\norm{\dot{x}(t)}^2 \+ \snorm{\sqrt{\linA} \, x(t)}^2
\=
\norm{\dot{x}(0)}^2 \+ \snorm{\sqrt{\linA} \, x(0)}^2
\quad (0 < t < \T),
\]
which implies that the problem owns mostly one solution.
Let $\{\linP_{\lambda}\}_{\lambda \xeof \R}$ be
the spectral resolution of $\linA$.
Then we define a function $x \: [0,\T) \xto \hilH$
by means of the spectral operator $\linL_t$,
which itself is induced by the functions
$t \mapsto \cos(t)$ and $t \mapsto \sin(t)$.

\subparagraph{}
Indeed,
for $0 \leq t {\,<\,} \T$
the operator $x \mapsto \linL_t x$ is built up
as the sum of two single spectral operators,
namely
\[
x(t)
:=
\linL_t(x_0,x_1)
:=
\int_{\,0}^{\,\infty} \cos{\sqrt{\lambda}t} \, dP_{\lambda} \, x_0
\+
\int_{\,0}^{\,\infty} \frac{\sin{\sqrt{\lambda}t}}{\sqrt{\lambda}}
\, dP_{\lambda} \, x_1.
\]

\subparagraph{} %- well-posedness of the second-order Cauchy problem
Having defined the family of linear operators $\{\linL_t(x_0,x_1)\}$ this way,
we obtain the following existence and uniqueness theorem
for the second-order homogeneous Cauchy problem:

\begin{theorem}
The mapping $t \mapsto \linL_t(x_0,x_1)$ is continuous in $[0,\T)$,
twice continuously differentiable in $(0,\T)$ and
constitutes the unique solution
of problem \eqref{HOMOGENEOUS_CAUCHY_PROBLEM_II}.
\end{theorem}

\item[3.]
By the \textit{abstract {\SCHROEDINGER} equation} is meant
the linear evolution equation
\begin{equation}\label{ABSTRACT_SCHROEDINGER_EQUATION}
\begin{cases}[1.2]
\quad \dot{x}(t) \, + \, i \linA \, x(t) \= \0,\\
\quad x(0) \= x_0 \xeof \hilH,\\
\end{cases}
\end{equation}
where $\linA$ is a self-adjoint linear operator.
Each solution $x$ of \eqref{ABSTRACT_SCHROEDINGER_EQUATION}
satisfies the \textit{energy equation}
\begin{equation}\label{ENERGY_CONSERVATION_LAW}
\norm{x(t)} = \norm{x(0)}
\end{equation}
for $t \xeof \R$,
which immediately follows from
\[
\frac{d}{dt} \, \norm{x(t)}^2
\=
\scalar{-i\linA x(t)}{x(t)} + \scalar{x(t)}{-i \linA x(t)} \= 0.
\]
If an initial value $x(0) \eq x_0$ is given,
then \eqref{ENERGY_CONSERVATION_LAW} implies
that the solution of the initial-value problem 
for \eqref{ABSTRACT_SCHROEDINGER_EQUATION} is unique.

\subparagraph{}
To construct a solution
by the functional calculus method for self-adjoint operators,
we define the spectral operator
\[
L_t(x) 
\eqdef
e^{-it \linA} x
\=
\int_{\,-\infty}^{\,\infty} e^{-it \lambda} dP_{\lambda} \, x
\quad \quad (t \xeof \R, x \xeof \hilH),
\]
where $P_{\lambda}$ again is the spectral resolution of $\linA$
having support in the whole real line.
Then the family of operators $L_t$ yields the solution
of the abstract {\SCHROEDINGER} problem:

\begin{theorem}\label{EXISTENCE_THEOREM_FOR_ABSTRACT_SCHROEDINGER_EQUATION}
The mapping $t \xeof \R \mapsto \linL_t(x_0)$ is the unique solution
of the Cauchy problem for \eqref{ABSTRACT_SCHROEDINGER_EQUATION}.
\end{theorem}

\subparagraph{}
Moreover,
for each $t \xeof \R$ the operator $\linL_t$ is unitary,
and we have $\linL_t^{-1} \eq \linL_{-t}$ for all $t \xeof \R$.
\end{itemize}

%------------------------------------------------------------------------------%
\section{Operators with Pure Point Spectrum}                                   %
%------------------------------------------------------------------------------%
%-- 
%-- 8. OPERATORS WITH PURE POINT SPECTRUM
%--    ----------------------------------
%-- discrete spectrum
%------------------------------------------------------------------------------%
Recall from section 5 that we have introduced various notations of essential
spectrum and especially the Browder spectrum $\sigma_{ess,6}(\linA)$.
On the other hand,
the \textit{discrete spectrum} $\sigma_d(\linA)$ of a linear operator $\linA$
acting in a complex Banach space is defined to be
the relative complement of $\sigma_{ess,6}(\linA)$ {\wrt} $\sigma(\linA)$.
Hence the set $\sigma_d(\linA)$ contains
those points $\lambda \xeof \sigma(\linA)$,
which are isolated in $\sigma(\linA)$
and for which the rank of the \textit{Riesz projection} $P_{\lambda}$ is finite.
In this case we have a disjoint decomposition of $\sigma(\linA)$
according to
\[
\sigma(\linA) \eq \sigma_d(\linA) \, \sqcup \, \sigma_{ess,6}(\linA).
\]
where $P_{\lambda}$ is defined by the \textit{Dunford integral}
\begin{equation}\label{DUNFORD_FORMULA}
P_{\lambda}(x) \eqdef
\frac{1}{2 \pi i} \, \int_{\,\gamma_{\lambda}} \, R(\linA,\zeta) x \, d\zeta
\quad (x \xeof \banX)
\end{equation}
and $\gamma_{\lambda}$ is a closed continuous path enclosing $\lambda$
but no other values in $\sigma(\linA)$.
One may also introduce the set $\sigma_d(\linA)$ as the collection
of isolated points $\lambda \xeof \sigma(\linA)$
having finite algebraical multiplicity
(this magnitude is the Hamel dimension of the linear space
\[
E_{gen}(\linA,\lambda) \eqdef \cup_{n \in \N} E(\linA^n,\lambda)
\]
of \textit{generalized} eigenvalues of $\linA$).
Elements of the set $\sigma_d(\linA)$ are called \textit{normal eigenvalues}.
It is meaningful to say
that $\linA$ has \textit{pure point spectrum}
if $\sigma(\linA) \eq \sigma_d(\linA)$ and 
$\sigma_{ess,6}(\linA) \eq \emptyset$,
respectively.

\subsection{} %-- Riesz-Schauder theory for compact linear operators
%------------------------------------------------------------------------------%
The theory of linear operators having pure point spectrum
mainly relies on two fundamental theorems:

\begin{itemize}[leftmargin=12pt]
\item[1.]
the \textit{spectral mapping theorem for the discrete spectrum},
quoting that
the discrete spectrum of any resolvent $R(\linA,\zeta)$ depends
on that of $\linA$ according to
\[
\sigma_d(R(\linA,\zeta)) \setminus \{0\} \=
\{(\lambda-\zeta)^{-1} \: \lambda \xeof \sigma_d(\linA)\};
\]

\item[2.]
the \textit{Riesz-Schauder} \textit{spectral theorem},
which illustrates
that the spectral theory of compact linear operators does in principle
not differ from that of linear transformations
in the finite-dimensional linear spaces $\C^d$.
Before quoting this theorem,
let us have a short look to this case,
which can completely be covered by the theory of square matrices.

\subparagraph{} %- geometric and algebraic multiplicity
In finite-dimensional linear algebra,
the spectral theory begins with the study of eigenvalues and eigenvectors
of a square matrix $\A \xeof \C^{d \times d}$.
Remember that an eigenvalue $\lambda$ of $\A$ is a complex number
satisfying $\A x \eq \lambda x$ for some nonzero vector $x \xeof \C^d$,
called an eigenvector corresponding to $\lambda$.
The set of all eigenvalues of $\A$, counting multiplicities,
forms the spectrum $\sigma(\A)$.
Moreover,
the eigenvalues of $\A$ are precisely the roots
of its characteristic polynomial $p_{\A}(\lambda) \eq det(\lambda {-} \A)$,
a monic polynomial of degree $d$.
The algebraic multiplicity of an eigenvalue $\lambda$ is
the multiplicity of $\lambda$ as a root of $p_{\A}(\lambda)$.
The geometric multiplicity of $\lambda$,
however,
is the dimension of the eigenspace $\kernel{\A {-} \lambda}$,
which is at most the algebraic multiplicity.

\subparagraph{} %- Jordan canonical form
A matrix $\A$ is diagonalizable \iff
for every eigenvalue geometric multiplicity equals the algebraic multiplicity.
When $\A$ is not diagonalizable,
its structure is captured by the \textit{Jordan canonical form},
which decomposes $\C^d$ into generalized eigenspaces.
For an eigenvalue $\lambda$,
the generalized eigenspace is nothing but $\kernel{\A{-}\lambda}^m$,
where $m$ is the \textit{index} of $\lambda$,
defined as the smallest positive integer
such that $\kernel{\A{-}\lambda}^m \eq \kernel{\A{-}\lambda}^{m+1}$.
This index equals the size of the largest Jordan block
associated with $\lambda$.
The space $\C^n$ then decomposes as a direct sum
of these generalized eigenspaces over all distinct eigenvalues.
The \textit{Jordan canonical form theorem} states that
every square matrix $\A$ over an algebraically closed field such as $\C$
is similar to a unique (up to permutation of blocks) Jordan matrix $J$,
meaning there exists an invertible matrix $P$ such that $\A \eq PJP^{-1}$.
The matrix $J$ is block-diagonal,
consisting of Jordan blocks $J_k(\lambda)$ for each eigenvalue $\lambda$,
where a $k \xtimes k$ Jordan block has the value $\lambda$ on the main diagonal
and the value $1$ on the superdiagonal.
This form was established by \textsc{C.\,Jordan} in his 1870 treatise.
A simple example of a non-diagonalizable matrix is
\[
\A \eq \begin{pmatrix} 0 & 1 \\ 0 & 0 \end{pmatrix},
\]
which has the eigenvalue $0$ with algebraic multiplicity $2$
but geometric multiplicity $1$.

\paragraph{} %- Riesz-Schauder spectral theorem
A linear operator $\compA$ acting from a normed linear space $\banX$
into a normed linear space $\banY$ is called \textit{compact},
if it maps the unit ball $B_1(\0) \xsubset \banX$
into a pre-compact subset of\,$\banY$,
that is,
the closure of the image $\compA(B_1(\0))$ is compact.
Then the announced theorem
elaborated by \textsc{F.\,Riesz} and \textsc{J.\,Schauder}
reads as follows:

\begin{theorem}\label{RIESZ-SCHAUDER_SPECTRAL_THEOREM}
Let $\banX$ be a complex Banach space
and $\compA \: \banX \xto \banX$ be compact.
Then the following assertions hold:

\begin{itemize}[leftmargin=15pt]
\item[(a)]
The spectrum $\sigma(\compA) \subset \C$ is mostly countable.

\item[(b)]
$\lambda \eq 0$ is the only possible limit point of $\sigma(\compA)$.

\item[(c)]
If $\banX$ is infinite-dimensional, then $0 \xeof \sigma(\compA)$.

\item[(d)]
Each $0 \neq \lambda \xeof \sigma(\compA)$ is an eigenvalue of $\compA$
with $\alpha(\compA-\lambda) < \infty$. %of finite-dimensional eigenspace.

\item[(e)]
Every spectral value $\lambda \neq 0$ is a pole of finite order
of the resolvent function $\zeta \xeof \rho(\compA) \mapsto R(\compA,\zeta)$.
\end{itemize}
\end{theorem}
\end{itemize}

\subparagraph{} %-- operators with compact resolvents
Both theorems together imply
that a linear operator $\linA$ has pure point spectrum
\iff
at least one of the resolvents is compact,
hence $\linA$ is said to \textit{have compact resolvents} in this case.
But notice that if one resolvent of $\linA$ is compact,
then all other resolvents are so,
which is a consequence of the resolvent equation \eqref{RESOLVENT_EQUATION}.

\subparagraph{} %- characterization of the property to have compact resolvents
Furthermore,
a closed linear operator $\linA$ has compact resolvents
\iff the embedding of the Banach space $(\domain{\linA},\norm{\,\cdot\,}_A)$
into the space $(\banX,\norm{\,\cdot\,})$ is compact,
where the domain of definition $\domain{\linA}$ is furnished with
the \textit{operator norm} $\norm{x}_A := \norm{x} + \norm{\linA x}$.

\paragraph{} %- order of the eigenvalues
%------------------------------------------------------------------------------%
The eigenvalues of a linear operator $\linA$ with pure point spectrum
are mostly countable and can be ordered by their moduli,
that is,
there is a sequence $\{ \lambda_n \}_{n \xeof \N}$ consisting
of eigenvalues of $\linA$
such that
\[
\abs{\lambda_1} \,\leq\, \abs{\lambda_2} \,\leq\, \abs{\lambda_3} \,\leq\,
\ldots \, .
\]

\subparagraph{} %- the symmetric case
But more can be said in the case,
where $\banX \eq \hilH$ is a Hilbert space
and the operator $\linA$ is symmetric.
Then all eigenvalues $\lambda_n$ are real-valued
and consequently $\linA$ proves to be self-adjoint.
If in addition $\linA$ is half-bounded with lower bound $m$,
we may improve the order of the sequence of eigenvalues to
\[
m \,\leq\, \lambda_1 \,\leq\, \lambda_2 \,\leq\, \lambda_3 \,\leq\, \ldots \, .
\]

\subsection{} %-- the discrete spectral theorem
%------------------------------------------------------------------------------%
To the end,
if the underlying Hilbert space is separable,
then the collection of eigenvectors $\{x_n\}_{n \xeof \N}$
of the operator $\linA$ forms an orthonormal base.
This statement is the content of
the celebrated \textit{Hilbert-Schmidt} \textit{spectral theorem}:

\begin{theorem}\label{DISCRETE_SPECTRAL_THEOREM}
If the linear operator $\compA$
acting in a separable Hibert space $\hilH$
is compact and symmetric,
then the collection $\orthO \xsubset \hilH$ of normed eigenvectors of $\compA$
is complete and constitutes an orthonormal system in $\hilH$,
respectively.
\end{theorem}

\subparagraph{} %- spectral decomposition of operators with pure point spectrum
We conclude from this theorem,
often also denoted the \textit{discrete spectral theorem},
that a symmetric linear operator $\linA$ with pure point spectrum
admits the \textit{spectral decomposition}
\[
\linA \,x \= \sum_{n=1}^{\infty} \lambda_n \scalar{x}{x_n} \, x_n
\]
for all $x \xeof \domain{\linA}$,
hence the domain of definition of $\linA$ can be described according to
\[
\domain{\linA} \=
\{x \xeof \hilH:
\sum_{n \eq 1}^{\infty} \, \lambda_n^2 \, \abs{\scalar{x}{x_n}} ^2 < \infty \}.
\]

\subparagraph{} %- spectral resolution of a self-adjoint operator wpps.
The discrete spectral theorem can be seen as a special case
of the general spectral theorem \ref{SPECTRAL_THEOREM}
for unbounded self-adjoint operators.
Indeed,
the spectral resolution of a self-adjoint operator with pure point spectrum
consists of the projections
\[
P_{\lambda} \eqdef \sum_{\lambda_k < \lambda} \pi_{\lambda_k}
\quad \quad (\lambda \in \R),
\]
where $\pi_{\lambda_k}$ is the projection onto the eigenspace $E_{\lambda_k}$
of the eigenvalue $\lambda_k$,
hence the mappings $P_{\lambda}$ are constant for
$\lambda \xeof [\lambda_k,\lambda_{k+1})$.

\paragraph{} %- equivalence of linear equations
%------------------------------------------------------------------------------%
Let us in the second part of this section
turn back to the class of Fredholm operators
and illustrate the interplay of Fredholm theory
with the theory of linear operators having pure point spectrum.

\subparagraph{} %- the operator 1+k is Fredholm
If $k \: \banX \xto \banX$ is compact,
then it follows from theorem \ref{RIESZ-SCHAUDER_SPECTRAL_THEOREM} (d),
that the operator $\1_{\banX}{+}k$ has finite nullity index and closed range,
hence it is upper semi-Fredholm.
But this is also true for the operator $\1_{\banX^*}{+}k^*$,
since $k^*$ is compact, too,
and both properties together imply that $\1_{\banX}{+}k$ is Fredholm with
index $\chi(\1_{\banX}{+}k) \eq 0$.

\subparagraph{} %- equivalence of equations
The following computation is of central meaning 
and for example useful in solving the \textit{Dirichlet problem}
for linear-elliptic differential operators.
Consider the linear equation $\linA x \eq y$,
where $\banX$ is a Banach space and
$\linA \: \domain{\linA} \xsubset \banX \xto \banX$
is a linear operator with pure point spectrum.
Then for $\zeta \xeof \rho(\linA)$ we surely have
\[
\linA \,x \= (\linA {-}\zeta)\,x + \zeta \,x \= y,
\]
and applying the mapping $(\linA {-}\zeta)^{-1}$ from the left yields
\[
x \+ \zeta (\linA {-} \zeta)^{-1}\,x
\=
(\linA {-}\zeta)^{-1}\,y \, =: \, \tilde{y}
\]
or equivalently $x + \tilde{\compA} \,x \eq \tilde{y}$,
where $\tilde{\compA} := \zeta (\linA {-} \zeta)^{-1}$ is compact.
The equations $\linA \,x \eq y$ and
$(\1_{\banX} + \tilde{\compA})\,x \eq \tilde{y}$
are called \textit{equivalent to each other} in the sense
that the element $x$ is a solution of the first one
\iff
it is a solution of the second one.

\paragraph{}
%------------------------------------------------------------------------------%
We summarize the result of the previous computation to the following

\begin{theorem}\label{FREDHOLM_ALTERNATIVE_THEOREM_II}
Let $\linA \: \domain{\linA} \xsubset \banX \xto \banX$
be an operator with pure point spectrum.
Then the equation $\linA x \eq y$ is equivalent to an equation
of the form $(\1_{\banX} + \tilde{\compA})x \eq \tilde{y}$,
where the operator $\tilde{\compA}$ is compact.
\end{theorem}

\subparagraph{}
Since $\1_{\banX} + \tilde{\compA}$ is Fredholm with vanishing index,
\textsc{Fredholm}'s alternative theorem \ref{FREDHOLM_ALTERNATIVE_THEOREM}
is applicable and therefore provides also assertions
on the solution of the equation $\linA x \eq y$.

\subsection{} %-- concepts of solvability for linear equations
%------------------------------------------------------------------------------%
Let us finalize the elaborations on linear equations
with a complete collection of concepts
regarding the solvability of $\linA \, x \eq y$,
where $\linA$ is assumed to be closed
and acting between some Banach spaces $\banX$ and\,$\banY$.
The equation $\linA \, x \eq y$ is called

\begin{itemize}[leftmargin=9mm]
\item[(a)]
\textit{densely solvable} (\cite{Kr}),
if $\overline{\range{\linA}} \eq \banY$;

\item[(b)]
\textit{normally solvable} (\cite{Kr}),
if $\overline{\range{\linA}} \eq \range{\linA}$;

\item[(c)]
\textit{correctly solvable} (\cite{Kr}),
if $\linA$ is bounded below;

\item[(d)]
\textit{semi-solvable} (\cite{Br1}), if $\linA$ is Fredholm;

\item[(e)]
\textit{almost solvable} or \textit{regularly solvable} (\cite{Br1}),
if $\linA$ is semi-solvable and $\chi(\linA) \eq 0$;

\item[(f)]
\textit{uniquely solvable} (\cite{Kr}),
if $\linA$ is injective;

\item[(g)]
\textit{everywhere solvable} (\cite{Kr}),
if $\linA$ is surjective;

\item[(h)]
\textit{solvable} (\cite{Br1}),
if it is uniquely and everywhere solvable;

\item[(i)]
\textit{well-posed},
if it is solvable and $\linA^{-1}$ is continuous,
or equivalently,
if $0 \xeof \rho(\linA)$;

\item[(j)]
\textit{completely solvable} (\cite{Br1}) or \textit{compactly solvable}
if it is solvable and $\linA^{-1}$ is compact.
\end{itemize}

%------------------------------------------------------------------------------%
\section{Infinitesimal Generators of Semigroups}                               %
%------------------------------------------------------------------------------%
%-- 
%-- 9. INFINITESIMAL GENERATORS OF SEMIGROUPS
%--    --------------------------------------
This section contains a short outlook to \textit{semigroup theory},
the subject that deals with families of bounded linear operators
fulfilling two fundamental features.
The theory of semigroups proves to be one of the most suitable tools
in the study of linear evolution equations.

\subparagraph{} %-- strongly continuous semigroups of bounded linear operators
Given a complex Banach space $\banX$,
a \textit{strongly continuous semigroup of bounded linear operators}
(shortly denoted a \textit{C$_0$-semigroup}) is an operator-valued mapping
$
\linS \: [0,\infty) \xto \BO{\banX},
$
also denoted {\SG},
which distinguishes by the conditions

\begin{itemize}[leftmargin=15mm]
\item[(G-0)]
\;$\linS(0) = \1_{\banX}$;

\item[(G-1)]
\;$\linS(t+s) \eq \linS(t) \circ \linS(s)$ for $t,s \geq 0$
(\textit{semigroup property});

\item[(G-2)]
\;$\lim_{\,t \downarrow 0} \, \linS(t)\,x \eq x$ for all $x \xeof X$.
\end{itemize}

\subparagraph{} %- continuity from the right
It follows from (G-2),
which is nothing but strong continuity from the right in $t \eq 0$,
that $\linS(0)$ is the identity map $\1_{\banX}$.
The conditions (G-1) and (G-2) together
imply strong continuity from the right of the semigroup
even in every $t {\,>\,} 0$.
In case of $\banX \eq \R$,
the family $\{\cal{E}(t)\}_{\, t \geq 0}$ of exponential functions
$\cal{E}(t) \: x \xeof \R \mapsto e^{\lambda tx}$
fulfills both conditions,
so it serves as an elementary example of a C$_0$-semigroup.

\subparagraph{} %- classification of strongly continuous semigroups
The totality of strongly continuous semigroups
may be classified as follows:
for each such semigroup {\SG} there are constants $M \geq 0$
and $\omega \xeof \R$
such that
\begin{equation}\label{X_SEMIGROUP_GROWTH}
\norm{\linS(t)} \, \leq \, M e^{\omega t}
\quad \for t > 0.
\end{equation}

\subparagraph{} %- growth bound
We denote $\omega$ the \textit{growth bound} of the semigroup.
In the special case,
where $(M,\omega) \eq (0,1)$ and $\norm{S(t)} \leq 1$ for all $t \geq 0$,
respectively,
the semigroup is said to be \textit{contractive}.
If $\omega \eq 0$,
however,
the semigroup is called \textit{uniformly bounded}.

\subsection{} %-- infinitesimal generators
%------------------------------------------------------------------------------%
By the \textit{orbit map} or \textit{trajectory} $\eta_x$
of a C$_0$-semigroup {\SG}
and associated with the initial value $x \xeof \banX$
we mean the vector-valued mapping $\eta_x \: [\,0,\infty) \xto \banX$
defined by $\eta_x(t) := S(t) \, x$.
The conditions (G-1) and (G-2) together
imply continuity of the map $\eta_x$.
But $\eta_x$ is even differentiable under a mild assumption:
if {\SG} is strongly continuous,
then for each $x \xeof \banX$
the mapping $\eta_x$ is differentiable in $(0,\infty)$
iff $\eta_x$ is right-differentiable in $t \eq 0$.

\subparagraph{} %- definition of generating operator
The last statement serves as motivation for the following notion:
if {\SG} is strongly continuous,
its \textit{infinitesimal generator} of the semigroup
(briefly denoted the \textit{generator} or the \textit{generating operator})
is defined as the mapping $\linA \: \domain{\linA} \xsubset \banX \xto \banX$
given by the graph
\[
\Gamma_{\linA} \eqdef
\{\,(x,y) \xeof X \xtimes X: \;
y \= \lim_{\,h \downarrow 0} \, \frac{1}{h}(\linS(h)x-x) \mbox{ exists} \}.
\]
Hence the domain of definition $\domain{\linA}$ consists
of those elements $x \xeof \banX$,
for which the orbit map $\eta_x$ is differentiable from the right in $t \eq 0$.
But this means that
$\eta_x$ is also (left and right) differentiable in every $t>0$,
which is a consequence of property (G-1).

\subparagraph{} %- properties of infinitesimal generators
It is easy to show that
the infinitesimal generator of a strongly continuous semigroup is linear,
so it forms a linear operator in $\banX$.

\paragraph{} % infinitesimal generators are densely defined and closed
%------------------------------------------------------------------------------%
Our first theorem in this section describes two basic properties
of infinitesimal generators:

\begin{theorem}
The infinitesimal generator of a C$_0$-semigroup is densely defined
and closed.
\end{theorem}

Hence the generators of strongly continuous semigroups fit
into our base class of linear operators considered in section 2.

\subsection{} %-- well-posedness of first order Cauchy problems
%------------------------------------------------------------------------------%
Functional analysis provides a rich theory
on \textit{homogeneous} linear evolution equations of the form
\begin{equation}\label{HOMOGENEOUS_CAUCHY_PROBLEM_IB}
\begin{cases}[1.2]
\quad \dot{x}(t) \+ \linA \, x(t) \= 0,\\ 
\quad x(0) \= x_0 \xeof \banX,\\
\end{cases}
\end{equation}
where $\linA$ is a densely defined and closed linear operator
acting in a Banach space $\banX$
and with domain of definition $\domain{\linA} \xsubset \banX$.

\subparagraph{}
Like in section 7,
it is to find a solution $x \: [0,\T) \xto \banX$
of the equation $\dot{x} + \linA x \eq 0$ 
satisfying the initial value condition $x(0) \eq x_0$.
This problem has already been considered in the Hilbert space setting,
where $\linA$ was assumed to be self-adjoint and positive,
and the solution of the problem \eqref{HOMOGENEOUS_CAUCHY_PROBLEM_IB}
was constructed by means of the spectral resolution of $\linA$
and the functional calculus for unbounded self-adjoint operators.

\paragraph{} %- classical and mild solutions
%------------------------------------------------------------------------------%
We first have to define concepts of solution
for problem \eqref{HOMOGENEOUS_CAUCHY_PROBLEM_IB}:

\begin{itemize}[leftmargin=12pt]
\item[1.]
A vector-valued function
$x \xeof C([0,\T),\domain{\linA}) \, \cap \, C^1((0,\T),\banX)$
is called a \textit{classical solution}
of \eqref{HOMOGENEOUS_CAUCHY_PROBLEM_IB},
if it satisfies the initial condition $x(0) \eq x_0$
and the differential equation in each point $t \xeof (0,\T)$.

\item[2.]
A continuous function $x \: [0,\T) \xto \banX$ is said to be
a \textit{mild solution} of \eqref{HOMOGENEOUS_CAUCHY_PROBLEM_IB} if
\[
\mbox{(1) } \, \tilde{x}(t) \eqdef
\int_{\,0}^t x(s) \, ds \, \xeof \domain{\linA}
\quad \and \quad
\mbox{(2) } \, \linA \, \tilde{x}(t) \eq x(t)-x_0
\]
holds for all $0 {\,<\,} t {\,<\,} \T$,
where the integral on the right-hand side of (1) 
is meant in the Riemannian sense.
One is led to the notion of mild solution in a natural way
by integrating both sides of the equation $\dot{x} + \linA \, x \eq \0$,
which then becomes into
\[
\int_{\,0}^t \dot{x} ds \=
x(t)-x(0) \= -
\int_{\,0}^t \linA \, x(s) \, ds.
\]
If the operator $\linA$ is closed,
then the action of $\linA$ and the execution of the integral commutes
for continuous functions $x \: [0,t] \xto \domain{\linA}$,
that is,
\[
\int_{\,0}^t \linA \, x(s) \, ds  \= \linA \, \int_{\,0}^t x(s) \, ds,
\]
and to the end we obtain the equation
\[
x(t) \= x(0) - \linA \, \int_{\,0}^t x(s) \,ds
\quad \quad (t \xeof (0,\T)).
\]
\end{itemize}

\subparagraph{} %- comparing classical and mild solutions
Notice that each classical solution is a mild solution,
and reversely a mild solution is classical
iff it is continuously differentiable in $(0,\T)$.
Thus classical and mild solutions of homogeneous Cauchy problems
differ {\wrt} this regularity condition.
It should also be mentioned
that if a solution $x$ is classical,
then $x(0) \eq x_0$ must belong to the domain of definition $\domain{\linA}$,
whereas for a mild solution $x$
the initial value $x_0$ may be an arbitrary element in $\banX$.

\paragraph{} %- correctly posed and well-posed problems
%------------------------------------------------------------------------------%
The homogeneous Cauchy problem \eqref{HOMOGENEOUS_CAUCHY_PROBLEM_IB}
is called \textit{correctly posed},
if there is a unique classical solution $x$
for any initial value $x_0 \xeof \domain{\linA}$,
that is,
$x \xeof C^1((0,\T);\banX) \cap C([0,\T);\banX)$.

\subparagraph{}
Moreover,
the problem is called \textit{well-posed},
if it is correctly posed and
the solution $x$ depends continuously on the initial value $x_0$.
In other words,
the problem is \textit{well-posed},
if a classical solution exists and
for each $\T {\,>\,} 0$ there is $c_{\T} {\,>\,} 0$ such that
$\norm{x(t)} \leq c_{\T} \norm{x_0}$
for all $t \xeof (0,\T)$ and all $x_0 \xeof \banX$.

\paragraph{} %- well-posedness theorem for the Cauchy problem
%------------------------------------------------------------------------------%
The following fundamental theorem describes
well-posedness of \eqref{HOMOGENEOUS_CAUCHY_PROBLEM_IB}
in terms of the property of the linear operator ${-}\linA$
to be an infinitesimal generator:

\begin{theorem}\label{WELL_POSEDNESS_THEOREM_I}
Let $\linA$ be a densely defined and closed linear operator
acting in a Banach space $\banX$.
Then the problem \eqref{HOMOGENEOUS_CAUCHY_PROBLEM_IB} is well-posed
\iff
${-}\linA$ is the infinitesimal generator of a C$_0$-semigroup.
\end{theorem}

\subparagraph{}
Indeed,
the family {\SG} of linear operators
having ${-}\linA$ as its infinitesimal generator
yields the classical solutions $x$ with initial values
$x_0 \xeof \domain{\linA}$ of \eqref{HOMOGENEOUS_CAUCHY_PROBLEM_IB},
that is,
$x(t,x_0) \eq S(t)\,x_0$ for $0 \leq t {\,<\,} \T$.
The proof of this fundamental well-posedness theorem can be found in \cite{EN},
p.\,112 as that of theorem II.6.6.

\subparagraph{}
Let us finally mention that problem \eqref{HOMOGENEOUS_CAUCHY_PROBLEM_IB}
governed by a densely defined and closed linear operator $\linA$ is well-posed
iff
it is correctly posed and $\rho(\linA) \ne \emptyset$ holds.

\paragraph{} %- necessary conditions to be a generator
%------------------------------------------------------------------------------%
Besides the fact that
the operator $\linA$ has to be densely defined and closed,
the further conditions on ${-}\linA$ to be the infinitesimal generator
of a C$_0$-semigroup can be expressed 
\textit{(i)}  by the location of the resolvent set $\rho(\linA)$
              within the complex plane and
\textit{(ii)} by a collection of estimates for the resolvents $R(\linA,\zeta)$
              in a subset of $\rho(\linA)$.

\subparagraph{}
Let the semigroup {\SG} be strongly continuous.
Then for fixed $x \xeof \banX$ and $\zeta \xeof \C$,
the value $L_n(\zeta,x)$ defined by
\begin{equation}\label{LAPLACE}
L_n(\zeta,x) \eqdef
\frac{1}{(n-1)!} \int_{\,0}^{\infty} t^{n-1}e^{\zeta t}S(t)x \,dt
\end{equation}
is called the (right side) \textit{Laplace transform} of order $n$
of the semigroup {\wrt} $x$ and $\zeta$.
This expression has to be understood as
improper vector-valued Riemann-Stieltjes integral.

\subparagraph{}
We cannot expect that $L_n(\zeta,x)$ exists without assuming further conditions.
But $L_n(\zeta,x)$ is surely well-defined for
values $\zeta$ with $\re \zeta < \omega$ and all $x \xeof \banX$.
Hence the mapping $L_n(\zeta)x := L_n(\zeta,x)$ constitutes a linear operator
in $\banX$ with norm estimate
\[
\norm{L_n(\zeta,\cdot)}
\, \leq \,
\frac{1}{(n-1)!} \int_0^{\infty} t^{n-1} e^{-\re \zeta t}e^{\omega t}\,dt
\= \frac{M}{(\re \zeta - \omega)^n}.
\]
A short computation shows that
\[
\begin{cases}[1.2]
\quad L_1(\zeta)(\zeta-\linA) \= (\zeta-\linA)L_1(\zeta) \= \1_{\banX},     \\
\quad L_n(\zeta)(\zeta-\linA) \= (\zeta-\linA)L_n(\zeta) \= L_{n-1}(\zeta), \\
\end{cases}
\]
which implies that $(-\infty,\omega) \xsubset \rho(\linA)$
and $R(\linA,\zeta) \eq L_n(\zeta)$ holds for all $n \xeof \N$.

\subsection{} %-- generator theorems for strongly continuous semigroups
%------------------------------------------------------------------------------%
Reversely,
the question arises,
under which conditions a linear operator ${-}\linA$ acting in $\banX$ forms
the infinitesimal generator of any C$_0$-semigroup.
For bounded operators the answer can be given quickly:
if $\norm{\linA}$ is the operator norm of $\linA$,
then the mappings
\[
\exp({-}\linA)(t)\,x \eqdef \sum_{n=0}^{\infty} \frac{(-t \linA)^n}{n!} \,x
\]
for $x \xeof \banX$ and $t>0$ form a strongly continuous semigroup.

\subparagraph{} %- Hille-Yosida theorem
The general answer,
however,
is given by the \textit{Feller-Miyadera-Hille-Yosida} \textit{theorem},
which characterizes properties of the operator ${-}\linA$
to be the infinitesimal generator of a C$_0$-semigroup: 

\begin{theorem}\label{MAIN_THEOREM_OF_SEMIGROUP_THEORY}
Let $\linA$ be a densely defined and closed linear operator
acting in an arbitrary Banach space $\banX$.
Then the following assertions are equivalent to each other:

\begin{itemize}[leftmargin=9mm]
\item[(a)]
The operator ${-}\linA$ generates a strongly continuous semigroup
of type $(\omega,M)$.
\item[(b)]
The resolvent set of $\linA$ fulfills $(-\infty,\omega) \xsubset \rho(\linA)$
and the resolvents $R(\linA,\zeta)$ satisfy the \textit{Hille-Yosida estimates}
\begin{equation}\label{HILLE-YOSIDA_ESTIMATES}
\norm{R(\linA,\zeta)^n} \, \leq \,
\frac{M}{(\omega - \zeta)^n} \quad \quad (n \xeof \N, \, \zeta < \omega).
\end{equation}
\end{itemize}
\end{theorem}

\subparagraph{}
The theorem shows that the already known conditions on ${-}\linA$
to be a generating operator are also sufficient.

\subparagraph{} %- Hille-Yosida theorem
However,
since the number of conditions on the operator $\linA$ is infinite,
this theorem has limited worth {\wrt} applications
and even in the simplest cases.
But the original theorem for generators of \textit{contraction semigroups}
($\omega \eq 0$ and $M \eq 1$)
due to \textsc{Hille} and \textsc{Yosida} quotes
that in this case only the first inequality
($n \eq 1$) in \eqref{HILLE-YOSIDA_ESTIMATES} has to be verified:

\begin{theorem}\label{HILLE-YOSIDA_THEOREM}
The operator ${-}\linA$ acting in a Banach space $\banX$
is the infinitesimal generator of a C$_0$-semigroup of \textit{contractions}
iff

\begin{itemize}[leftmargin=9mm]
\item[(a)]
$\linA$ is densely defined and closed,
\item[(b)]
the interval $(-\infty,0)$ belongs to the resolvent set $\rho(\linA)$,
\item[(c)] the resolvent estimation in the contractional case holds:
\begin{equation}\label{RESOLVENT_ESTIMATION_HY}
\norm{\zeta \, R(\linA,\zeta) } \leq 1 \quad (\zeta {\,<\,} 0).
\end{equation}
\end{itemize}
\end{theorem}

\subparagraph{} %- remarks on the proof of the Hille-Yosida theorem
We want to sketch how the assertions (b) and (c) together imply
that $\-\linA$ generates a contractional C$_0$-semigroup.
First of all,
it can be verified that \eqref{RESOLVENT_ESTIMATION_HY} implies
$
\lim_{\,\zeta \to -\infty} \, \zeta R(\linA,\zeta) \eq x
$
for all $x \xeof \domain{\linA}$.
As a consequence,
the so-called \textit{Yosida operators}
\[
\linA_n \eqdef n \linA (n-\linA)^{-1} \= n^2 - (\linA-n)^{-1} - n
\]
approximate the operator $\linA$ as $n \xto {-} \infty$ in the sense that
\[
\lim_{\,n \to -\infty} \, \linA_n x \= \linA x
\]
for all $x \xeof \domain{\linA}$.
Since the operators $\linA_n$ are bounded,
one may define semigroups for each of them by letting
\[
e^{-t \linA_n}\,x \eqdef
\sum_{k=0}^{\infty} \frac{(-t \linA_n)^k}{k!}\;x
\quad \quad (x \xeof \banX).
\]
It follows that $\snorm{e^{-t \linA_{\zeta}}} \leq 1$ for $t > 0$,
so we obtain even semigroups of contractions.
Now the main part of the proof is to show that
this sequence of semigroups is Cauchy and
the semigroup generated by the operator $-\linA$ is given by
\begin{equation}\label{YOSIDA_SEMIGROUP}
S(t)x \eqdef \lim_{\,n \to -\infty} e^{-t \linA_n}\,x
\quad \quad (x \xeof \banX).
\end{equation}

To be precise,
one first has to verify property (S-1) and (S-2)
and second the operator ${-}\linA$ to be
the generator of the semigroup \eqref{YOSIDA_SEMIGROUP}.

%------------------------------------------------------------------------------%
\section{Accretive Operators}                                                  %
%------------------------------------------------------------------------------%
%-- 
%-- 10. ACCRETIVE OPERATORS
%--     -------------------
%-- concepts of accretivity
The shapes and locations 
of the spectrum and the numerical range of
a densely defined and closed linear operator
acting in a complex Hilbert space
give rise to introduce a couple of subclasses of such operators,
and one of them is the class of \textit{accretive operators}.

\subparagraph{} %- accretive operators
A linear operator $\linA$ acting in $\hilH$ is called \textit{accretive},
if the numerical range $\Theta(\linA)$ fulfills
\[
\Theta(\linA) \subset \C^+ := \{\zeta \xeof \C \: \re \zeta \geq 0\},
\]
or equivalently,
if $\re \scalar{\linA x}{x} \geq 0$ for all $x \xeof \domain{\linA}$.
Furthermore,
the operator $\linA$ is said to be \textit{quasi-accretive},
if there is $\lambda_0 \xeof \R$ such that $\linA {+}\lambda_0$ is accretive.

\subparagraph{}
One verifies that the condition for accretivity is equivalent to
\begin{equation}\label{ACCRETIVE_PROPERTY}
\norm{(\linA {-} \zeta)x} \, \geq \, \norm{\zeta x}
\end{equation}
for $\zeta < 0$ and $x \xeof \domain{\linA}$ and also equivalent to
\[
\norm{(\linA {-} \re \zeta)x} \, \geq \, \norm{\re \zeta \, x}
\]
for $\zeta < 0$ and $x \xeof \domain{\linA}$.
But this means that the operators $\linA {-} \zeta$ are bounded from below,
that is, 
they are injective and own closed ranges in case of $\linA$ being closed.

\paragraph{} %- m-accretive operators
%------------------------------------------------------------------------------%
Finally,
the operator $\linA$ is said to be \textit{m-accretive},
if it is accretive and does not have any proper accretive extension.
Equivalently,
an accretive operator $\linA$ is m-accretive,
if there is $\zeta_0 \xeof (-\infty,0)$
such that $\range{\linA {-} \zeta_0} \eq \hilH$.
In this case
the operator $\linA {-} \zeta_0$ is actually bijective.
Hence we have the estimate
\[
\norm{\, \zeta_0 \, R(\linA,\zeta_0)\,} \, \leq \, 1,
\]  
where $R(\linA,\zeta_0)$ again denotes the resolvent
for $\zeta_0 \xeof \rho(\linA)$.

\subparagraph{} %- range condition for m-dissipative operators
Like for self-adjoint operators,
the necessarity of the existence of a number $\zeta_0 \xeof \C$
such that $\linA {-} \zeta$ be surjective,
is called the \textit{range condition} for m-dissipativity.
It follows from the stability theory of semi-Fredholm operators
that if such a number $\zeta_0$ exists,
then $\linA {-} \zeta$ is surjective even for every $\zeta \xeof (-\infty,0)$,
so $\zeta_0$ may be replaced by any other number $\zeta < 0$.

\paragraph{} %- dense definiteness and closeness of m-accretive operators
%------------------------------------------------------------------------------%
A m-accretive operator fits into the main class of densely defined and
closed linear operators,
as the following theorem shows:

\begin{theorem}\label{M_ACCRETIVITY_IMPLIES_CLOSENESS}
Each m-accretive operator $\linA$ is densely defined and closed,
and the Hilbert adjoint $\linA^*$ is m-accretive,
too.
\end{theorem}

\begin{comment}
\Proof:
To verify the first assertion,
let $f \xeof \domain{\linA} ^{\bot}$ and $x \xeof \domain{\linA} $ be a vector
such that $\linA x + x  \eq f$.
Sure, such a vector exists since $\linA+1$ is surjective.
Then we compute
\[
0 \= \scalar{x}{f} \=
\scalar{x}{\linA x +x} \= \scalar{\linA x}{x} + \norm{x}^2
\, \geq \, \norm{x}^2 \, \geq \, 0,
\]
which implies $x \eq 0$ and $f \eq \linA x+x  \eq  0$.
Hence $\domain{\linA} ^{\bot}  \eq \{\0\}$
and consequently $\closure{ \domain{\linA} } \eq \hilH$.

Two show the second assertion,
we have to verify that if $x_n \xto x$ and $\linA x_n \xto y$,
then $x \xeof \domain{\linA} $ and $\linA x \eq y$.
Clearly, $(\linA {+}1)x_n \xto y+x$, hence
\[
x \= \lim_{n \to \infty} x_n \=
\lim_{n \to \infty} (\linA {+}1)^{-1}(\linA {+}1) \, x_n
\=
(\linA {+} \1_{\banX})^{-1}(y+x),
\]
so
$x \= (\linA {+}\1_{\banX})^{-1}(y+x)
\, \Longleftrightarrow \, x + \linA x \eq y + x
\, \Longleftrightarrow \, \linA x \eq y$.
\qed
\end{comment}

\subparagraph{}
Some further assertions on m-accretive operators are collected
to the following theorem:

\begin{theorem}
The following statements are equivalent:
\begin{itemize}[leftmargin=9mm,itemsep=2pt]
\item[(a)] $\linA$ is m-accretive;
\item[(b)] $\linA + i\delta$ is m-accretive for all $\delta \xeof \R$;
\item[(c)] $\linA + \epsilon$ is m-accretive for all $\epsilon > 0$;
\item[(d)] $-\zeta \xeof \rho(\linA)$ and
           $\norm{(\linA-\zeta)(\linA+\zeta)^{-1}} \leq 1$ for $\zeta > 0$;
\item[(e)] $\C^- := \{z \xeof \C \: \re z < 0\} \xsubset \rho(\linA)$ and
           \[ \norm{R(\linA,\zeta)} \leq \frac{1}{\abs{\re \zeta} }
             \quad (\re \zeta < 0);
           \] 
\item[(f)] $\linA$ is densely defined and closed, and $\linA^*$ is m-accretive
           (this is the inverse statement of theorem
           \ref{M_ACCRETIVITY_IMPLIES_CLOSENESS}).
\end{itemize}
\end{theorem}

\subsection{} %-- Lumer-Phillips theorem
%------------------------------------------------------------------------------%
The usage of m-accretivity is illustrated by the next theorem.
By means of this concept,
the celebrated \textit{Lumer-Phillips} \textit{theorem}
classifies those linear operators,
which are infinitesimal generators of contractive C$_0$-semigroups:

\begin{theorem}\label{LUMER-PHILLIPS_THEOREM}
Let $\linA$ be a linear operator in Hilbert space.
Then ${-}\linA$ generates a contractive C$_0$-semigroup
iff $\linA$ is m-accretive.
\end{theorem}

\subparagraph{}
The proof of this theorem relies on the classical Hille-Yosida theorem
for C$_0$-semigroups of type $(0,1)$.
Indeed, $\linA$ is m-accretive,
then $\linA$ is densely defined, closed
and $-\linA$ fullfills the estimate \eqref{ACCRETIVE_PROPERTY},
hence $-\linA$ is the generator of a contractive C$_0$-semigroup.

\subparagraph{}
Reversely,
if $-\linA$ is the infinitesimal generator of a C$_0$-semigroup of contractions,
then we have $(0,\infty) \xsubset \rho(\linA)$
and therefore $\range{\linA {-} \zeta} \eq \hilH$ for all $\zeta > 0$.
In addition, if $x \xeof \domain{\linA}$, then
\[
\abs{ \scalar{S(t) x}{x}} \,\leq\, \norm{S(t)x} \norm{x} \,\leq\, \norm{x}^2
\]
and
\[
\re \scalar{S(t)x-x}{x} \= \re \scalar{S(t)x}{x} \,-\,\norm{x}^2 \,\leq\, 0.
\]
Dividing this line by $t > 0$ and letting $t \downarrow 0$ yields
$\re \scalar{\linA x}{x} \leq 0$,
which shows that $-\linA$ is accretive,
hence in summary the operator $\linA$ even proves to be m-accretive.

\subsection{} %-- well-posedness of second-order Cauchy problems
%------------------------------------------------------------------------------%
Similar to the first-order Cauchy problem,
a function $x$ is called
a {classical solution} of the \textit{second-order} homogeneous Cauchy problem
\begin{equation}\label{HOMOGENEOUS_CAUCHY_PROBLEM_2}
\begin{cases}[1.2]
\quad \ddot{x}(t) + \linA x(t) \= \0,\\
\quad x(0) \= x_0 \xeof \banX,\\
\quad \dot{x}(0) \= x_1 \xeof \banX,\\
\end{cases}
\end{equation}
if $x \xeof C([0,\T),\domain{\linA}) \, \cap \, C^2((0,\T),\banX)$
and satisfies both initial conditions as well as
the differential equation in each point $t \xeof (0,\T)$.

\paragraph{} %- the semigroup approach to second order Cauchy problems
%------------------------------------------------------------------------------%
To apply results from semigroup theory to the second order Cauchy problem,
let $U := (x,\dot{x})$.
We suppose that $\linA$ is a \textit{positive} and
\textit{self-adjoint} linear operator in a Hilbert space $\hilH$.
Then the differential equation $\ddot{x}(t) + \linA \, x(t) \eq \0$
in \eqref{HOMOGENEOUS_CAUCHY_PROBLEM_2}
can be written as first order problem
\begin{equation}\label{TRANSFORMED_SECOND_ORDER_EQUATION}
\begin{cases}[1.2]
\quad U'(t) \= \linB \, U(t) \\
\quad U(0) = (x_0,x_1),\\
\end{cases}
\end{equation}
where the operator $\linB$ is defined by
\begin{equation}\label{BLOCK_MATRIX}
\linB \eqdef \begin{pmatrix} \0 & \1 \\ \sqrt{\linA} & \0 \end{pmatrix}.
\end{equation}

\subparagraph{}
The main idea in proving existence and uniqueness of a solution is to show
that
the operator $\linB$ defined on a suitable chosen domain $\domain{\linB}$
is the infinitesimal generator of a semigroup of contractions.
Its domain of definition $\domain{\linB}$ is given by
\[
\domain{\linB} \eqdef
\domain{\linA} \times \domain{\sqrt{\linA} }
\, \subset \,
\domain{\sqrt{\linA} }_{\norm{\cdot}_{\linA}} \times \hilH,
\]
where $\domain{\sqrt{\linA} }_{\norm{\cdot}_{\linA}}$
is the Hilbert space $\domain{\sqrt{\linA} }$
furnished with the graph norm.

\paragraph{} %- Lumer-Phillips theorem for second order problems
%------------------------------------------------------------------------------%
The following theorem is for example proved in \cite{Go}, p.\,111:

\begin{theorem}\label{LUMER-PHILLIPS_THEOREM_2}
Let the linear operator $\linA$ be positive self-adjoint
and the operator $\linB$ be defined by \eqref{BLOCK_MATRIX}.
Then $\linB$ is m-accretive,
hence the Cauchy problem \eqref{TRANSFORMED_SECOND_ORDER_EQUATION}
is well-posed.
\end{theorem}

\subsection{} %-- regularly accretive operators
%------------------------------------------------------------------------------%
In the Hilbert space setting an accretive operator often appears as
the form operator associated
with a \textit{quasi-coercive} and continuous sesquilinear form
$\sql\, \vecV \xtimes \vecV \xto \C$,
defined by means of a Hilbert space triple $(\vecV,\hilH,\vecV_a^*)$.

\subparagraph{} %- pairs of Hilbert spaces and Gelfand triples
We want to make this assertion precise.
Each Hilbert pair $(\vecV,\hilH)$ with continuous embedding
$\iota \: \vecV \xto \hilH$
gives rise to introduce the \textit{transposed pair} $(\hilH_a^*,\vecV_a^*)$,
where the conjugate linear space $\hilH_a^*$
is continuously and densely embedded into
the conjugate linear space $\vecV_a^*$
by a further continuous embedding $\iota_2 \: \hilH_a^* \xto \vecV_a^*$.
Then,
by identification of the spaces $\hilH$ and $\hilH_a^*$
via the {Riesz} map,
combining both Hilbert pairs leads to the triple
\begin{equation}\label{GELFAND_TRIPLE}
\vecV \, \overset{\iota_1}{\hookrightarrow} \,
\hilH \, \overset{\iota_2}{\hookrightarrow} \, \vecV_a^*,
\end{equation}
sometimes called a \textit{Gelfand triple},
an \textit{evolution triple} or a \textit{rigged Hilbert space}.
The space $\hilH$ is often denoted the \textit{pivot space}
of \eqref{GELFAND_TRIPLE}.

\subparagraph{} %- the main property of a Gelfand triple
Gelfand triples own a remarkable property.
Let $(\,\cdot\,,\,\cdot\,)$ and $\scalar{\,\cdot\,}{\,\cdot\,}$
be the inner products in $\vecV$ and $\hilH$,
respectively,
then for fixed chosen $x \xeof \hilH$
the action of the antilinear linear functional
$f_x := \iota_2 \, x \xeof \vecV_a^*$ to any element $v \xeof \vecV$
is represented by the inner product of $x$ and $v$ in $\hilH$,
{\ie} $f_x(v) \eq \scalar{x}{\iota_1 \, v}$ holds for all $v \xeof \vecV$.

\paragraph{} %- coercive forms and the Lax-Milgram theorem
%------------------------------------------------------------------------------%
Furthermore,
a continuous sesquilinear form $\sql \: \vecV \xtimes \vecV \xto \C$
is called \textit{coercive}
if there is $\theta {\,>\,} 0$ such that
\[
\re \, \sql(x,x) \, \geq \theta \norm{x}^2
\qfa x \xeof \vecV.
\]
The celebrated \textit{Lax-Milgram theorem} claims,
that if $\sql$ is coercive,
then the \textit{variational equation}
\[
\sql(x,y) = f(y) \quad \quad (y \in \vecV)
\]
owns for each $f \xeof \vecV_a^*$ exactly one solution $x \xeof \vecV$.
In other words,
the representation operator $\A_{\sql} \: \vecV \xto \vecV_a^*$,
defined by 
\[
(\A_{\sql}x)(y) \eqdef \sql(x,y) \quad \quad (x,y \in \vecV)
\]
is bijective.
Moreover,
the mapping $\A_{\sql}$ has a bounded inverse
with norm $\norm{\A_{\sql}^{-1}} \leq \theta^{-1}$,
hence the equation $\linA_{\sql} x = f$ is well-posed.

\subparagraph{} %- quasi-coercive forms
In applications it is often the case
that sesquilinear forms are merely continuous but not coercive.
However,
if a Hilbert pair $(\vecV,\hilH)$ is given,
it is possible to make a form $\sql$ coercive by adding an additional term,
which is just the multiplicity of the inner product in $\hilH$
with a real number.
More precisely,
if $\sql \: \vecV \xtimes \vecV \xto \C$ is continuous
and the slightly modified form
\[
\sql_{\lambda}(u,v) \eqdef
\sql(u,v) \+ \lambda \scalar{\iota_1 u}{\iota_1 v}\, \quad
(u,v \xeof \vecV)
\]
is coercive for some $\lambda \xeof \R$,
that is,
there is $\theta {\,>\,} 0$ such that
\[
\re \sql(u,u)  \+
\lambda \norm{\iota_1u}_{\hilH}^2 \, \geq \, \theta \norm{u}_{\vecV}^2
\quad (u \xeof \vecV),
\]
then $\sql$ is called
\textit{quasi-coercive},
\textit{H-elliptic}
or simply a \textit{\GARDING} \textit{form}.
It should be clear that
if $\sql_{\lambda_0}$ is coercive for some $\lambda_0 \xeof \R$,
then $\sql_{\lambda}$ is so for $\lambda {\,>\,} \lambda_0$.

\paragraph{} %- operators associated with H-elliptic forms
%------------------------------------------------------------------------------%
Quasi-coercive sesquilinear forms can be used
to define densely defined and closed linear operators
acting in the pivot space $\hilH$ of a Gelfand triple.
To elaborate this assertion,
let $\sql$ be a sesquilinear form defined on a Hilbert space $\vecV$,
which is the left-hand side of a Gelfand triple $(\vecV,\hilH,\vecV_a^*)$
with continuous embeddings
$\iota_1 \: \vecV \xto \hilH$ and
$\iota_2 \: \hilH \xto \vecV_a^*$.
It is our aim to introduce a linear operator $\linA_{\sql}$
with domain space $\hilH$ and range space $\hilH$.
This will be done by restricting the form operator
$\repA \: \vecV \xto \vecV_a^*$ associated with $\sql$
to those elements $u \xeof \vecV$,
for which $\repA \, u \xeof \hilH$.
Indeed,
the domain of definition $\domain{\linA_{\sql}}$ is defined
to be the subspace of those elements $u \xeof \vecV$,
for which $\repA \, u \xeof \iota_2(\hilH) \xsubset \vecV_a^*$ holds,
briefly
\[
\domain{\linA_{\sql}} \eqdef \{u \xeof \vecV \: \repA \, u \xeof \hilH \}.
\]

\subparagraph{} %- variational operator
The set $\domain{\linA_{\sql}}$ can be regarded as a linear subspace
of $\hilH$ by means of the embedding $\iota_1$.
Then for $u \xeof \domain{\linA_{\sql}}$ there is $y \xeof \hilH$
such that $\iota_2 y \eq \repA \, u$.
We set $x := \iota_1 u$ and $\linA_{\sql} x := y$
for $x \xeof \domain{\linA_{\sql}}$
and denote $\linA_{\sql}$
the \textit{operator in $\hilH$ associated with $(\vecV,\sql)$},
or shortly,
the \textit{variational operator} derived from $\sql$.
Another phrase is to say that
$\linA_{\sql}$ is the \textit{part of $\repA$ in $\hilH$}.
Clearly,
we have $\sql{x,y} \= \scalar{\linA x}{y}$
for all $x \xeof \domain{\linA}$ and $y \xeof \vecV$.

\paragraph{} %- main theorem on H-elliptic forms and associated operators
%------------------------------------------------------------------------------%
Our last theorem in this section
collects a couple of properties of the operator $\linA_{\sql}$
associated with a quasi-coercive form $\sql$:

\begin{theorem}\label{QUASI-COERCIVE_FORMS_AND_ASSOCIATED_OPERATORS}
Let $(\vecV,\hilH,\vecV_a^*)$ be a triple of Hilbert spaces
and the sesquilinear form $\sql$ be continuous and quasi-coercive
on $\vecV$ with parameter $\lambda_0 \xeof \R$.
Let $\linA_{\sql}$ be the operator associated with $\sql$.
Then the following assertions are true:

\begin{itemize}[leftmargin=9mm]
\item[(a)]
$\linA_{\sql}$ is closed.

\item[(b)]
$\linA_{\sql}$ is densely defined in $\hilH$.

\item[(c)]
If the embedding of $\vecV$ into $\hilH$ is \textit{compact},
then the operator $\linA_{\sql} + \lambda$
owns a compact inverse $(\linA_{\sql} +\lambda)^{-1}$
for each $\re \lambda \geq \lambda_0$.

\item[(d)]
The operator $\linA_{\sql} + \lambda$ is m-accretive
for each $\re \lambda \geq \lambda_0$.

\end{itemize}
\end{theorem}

\begin{comment}
\subparagraph{}
To prove that $\linA_{\sql}$ is closed,
let $\{x_n\}_{n \xeof \N}$ be a sequence in $\domain{\linA_{\sql}}$
with limit element $x \xeof \hilH$
and $(\linA_{\sql} x_n$ be the sequence mapped by $\linA_{\sql}$
with limit element $y \xeof \hilH$.
It is to show that $x \xeof \domain{\linA_{\sql}}$ and $\linA_{\sql} x \eq y$.
Let $\abs{x} $ and $\norm{x}$ be the norms in $\vecV$ and $\hilH$,
respectively.
Clearly,
we have
\[
\norm{(\linA_{\sql}+\lambda) x_n - y - \lambda x} \xto 0 
\quad \for n \xto \infty.
\]
In the following we shall use the fact
that for $\lambda {\,>\,} \lambda_0$ the operators
$\linA_{\sql}+\lambda$ and $(\linA_{\sql}+\lambda)^{-1}$
forms linear homeomorphism between the normed spaces
$(\hilH,\norm{\,\cdot\,})$
and $(\domain{\linA_{\sql}},\abs{\,\cdot\,} )$,
so the latter one forms a closed subspace
within the normed space $(\vecV,\abs{\,\cdot\,} )$.
By this reason,
\[
x_n = (\linA_{\sql}+\lambda)^{-1} (\linA_{\sql}+\lambda) x_n
\xto 
(\linA_{\sql}+\lambda)^{-1} (y + \lambda x)
\quad \for n \xto \infty
\]
{\wrt} the norm $\abs{\,\cdot\,} $ in $\domain{\linA_{\sql}}$.
Since $x$ belongs to $\hilH$ and $\domain{\linA_{\sql}}$ is closed,
the element $(\linA_{\sql}+\lambda)^{-1} x$
belongs to $\domain{\linA_{\sql}}$ with $\linA_{\sql} x \eq y$.

\subparagraph{}
To prove that $\domain{\linA_{\sql}}$ is dense in $\hilH$,
suppose that $y \xeof \hilH$ fulfills $\scalar{x}{y} \eq 0$
for all $x \xeof \domain{\linA_{\sql}}$.
It is to show that $y \eq 0$ in this case.
Since $\linA_{\sql}+\lambda$ is surjective,
there must exist $x_0 \xeof \domain{\linA_{\sql}}$ such that
$(\linA_{\sql}+\lambda) x_0 \eq y$.
Then
\[
0 \=
\scalar{\linA_{\sql} x_0}{x} \+ \scalar{\lambda x_0}{x} \=
\sql(x_0,x) \+ \lambda \scalar{x_0}{x}.
\]
Especially for $x \eq x_0$ the last equation leads to
\[
0 \=
\scalar{\linA_{\sql} x_0}{x_0} \+ \scalar{\lambda x_0}{x_0} \,\geq\,
\theta \norm{x_0}^2
\]
for some $\theta {\,>\,} 0$,
due to to quasi-coercivity of $\sql$,
hence $x_0 \eq 0$ and finally $y \eq 0$.

\subparagraph{}
Finally,
if $\iota_1 \: \vecV \xto \hilH$ is compact,
then $\iota_1 \circ (\linA_{\sql}+ \lambda)^{-1} \: \hilH \xto \hilH$
is compact,
too,
since the composition of a compact linear operator
and a continuous linear operator proves to be compact.
\qed.
\end{comment}

\subparagraph{}
It is usual to call a linear operator \textit{regularly accretive},
if it is the operator associated with a quasi-coercive sesquilinear form.

%------------------------------------------------------------------------------%
\section{Sectorial Operators and Analytic Semigroups}                          %
%------------------------------------------------------------------------------%
%-- 
%-- 11. SECTORIAL OPERATORS AND ANALYTIC SEMIGROUPS
%--     -------------------------------------------
The concept of sectorialness is helpful to study semigroups,
which are \textit{analytic} in some sense and form the solution of so-called
\textit{parabolic} evolution equations,
whose solutions are infinitely differentiable in every $t > 0$.

\subparagraph{} %-- definition of sectorial operator
A linear operator $\linA$ acting in a complex Banach space $\banX$
is called \textit{sectorial},
if there is $\psi \xeof [0,\pi)$ and $\omega \xeof \R$ such that

\begin{itemize}[leftmargin=9mm]
\item[(S-1)]
the spectrum $\sigma(\linA)$ is contained in the closure
$\overline{\Sigma_{\,\psi,\omega} }$ of a right-open \textit{sector}
\[
\Sigma_{\,\psi,\omega} \eqdef
\{\zeta \xeof \C: \zeta \neq \omega, arg(\zeta-\omega) < \psi\}
\]
with vertex $\omega$ and semi-angle $\psi$;
this implies that 
$\rho(\linA)$,
the resolvent set of $\linA$,
contains the complementary and left-open sector
$\Lambda_{\,\pi-\psi,\omega} :=
\C \setminus \overline{ \Sigma_{\,\psi,\omega} }$;

\item[(S-2)]
given the angle $\phi < \pi$,
for every $\psi {\,<\,} \phi$
the following resolvent estimation holds:
\begin{equation}\label{RESOLVENT_ESTIMATE_II}
\sup \;
\{\, \norm{\,\zeta R(\linA,\zeta)\,} \:
  \zeta \xeof \Lambda_{\,\pi-\psi,\omega}
\}
\, < \, \infty\,;
\end{equation}
\end{itemize}

\subparagraph{}
Moreover,
the operator $\linA$ is called \textit{sectorial},
if (S-1) and (S-2) hold for some pair $(\psi,\omega) \xeof [0,\pi) \xtimes \R$.
Notice that
if $\linA$ is sectorial {\wrt} the angle $\psi$,
then there is $0 {\,<\,} \epsilon {\,<\,} \psi$
such that $\linA$ is also sectorial {\wrt} the smaller angle $\psi - \epsilon$.
The infimum $\psi_0$ of all these angles is called
the \textit{angle of sectoriality} of $\linA$.

\subparagraph{}
It should be clear that each sectorial operator is quasi-m-accretive
if the resolvent estimate \eqref{RESOLVENT_ESTIMATE_II} holds with $M \leq 1$.

\paragraph{} %- criterion for sectorialness
%------------------------------------------------------------------------------%
The following theorem is useful to ensure sectorialness:

\begin{theorem}\label{CRITERION_FOR_SECTORIALNESS}
Let $\linA$ be a regular linear operator in $\banX$ with domain of definition
$\domain{\linA} \xsubset \banX$ having the following properties:

\begin{itemize}[leftmargin=9mm]
\item[(a)]
$\rho(\linA)$ contains
$\C_{\omega} := \{\zeta \xeof \C \: \re \zeta \leq \omega\}$
for some $\omega \xeof \R$;

\item[(b)]
there is $M {\,>\,} 0$ such that
$\norm{\zeta \, R(\linA,\zeta)} \, \leq \, M$ for all $\zeta \xeof \C_{\omega}$.
\end{itemize}
Then $\linA$ is sectorial with vertex $\omega$.
\end{theorem}

\subparagraph{} %- proof of the criterion for sectorialness.
The proof of this theorem can be found in \cite{Lu}, p.\,43.

\paragraph{} %- properties of sectorial operators in Hilbert spaces
%------------------------------------------------------------------------------%
We next quote two properties of a sectorial linear operator $\linA$,
when acting in a Hilbert space $\hilH$:

\begin{itemize}[leftmargin=7mm,itemsep=3pt]
\item[1.]
$\linA$ is densely defined and closed.
\item[2.]
$\hilH \eq \kernel{\linA} \oplus \overline{\range{\linA}}$.
\end{itemize}

The second property also holds
if $\linA$ is acting in a reflexive Banach space.

\subsection{} %-- the functional calculus for sectorial operators
%------------------------------------------------------------------------------%
One of the goals in this section is
to sketch a functional calculus for a sectorial linear operator $\linA$,
which especially enables us to derive from $\linA$
a family of bounded linear operators $\{e^{-t \linA}\}_{t \geq 0}$.

\subparagraph{}
More precisely,
if $\linA$ is sectorial,
then the mapping $t \mapsto e^{-t \linA}$ is defined
by means of %the contour integrals
\begin{equation}\label{CONTOUR_INTEGRAL}
e^{-t \linA} x \eqdef
\frac{1}{2 \pi i}
\int_{\gamma} \, e^{-t \zeta} \, R(\linA,\zeta)x \, d\zeta
\quad \quad (x \xeof \hilH),
\end{equation}
which generalizes the \textit{Riesz-Dunford} \textit{formula}
\eqref{RIESZ-DUNFORD_FORMULA}.
Here by the symbol $\gamma$ is meant a rectificable path in the complex plane
that drifts within the sector $S_{\phi,\omega}$
and includes the spectrum of $\linA$.

\subparagraph{}
In fact,
the family of mappings $\{e^{-t \linA} \}_{t \geq 0}$
given by \eqref{CONTOUR_INTEGRAL}
constitutes a semigroup of bounded linear operators,
which is strongly continuous
if $\linA$ is assumed to be densely defined.

\subsection{} %-- analytic semigroups
Recall that a semigroup of bounded linear operators is usually parametrized
by the half line $[0,\infty)$.
But it is also customary to study possible extensions of a semigroup to a
symmetric sector
\[
\Sigma_{\,\theta} \eqdef
\{z=re^{i\alpha} \in \C: r > 0\,, \abs{\alpha} < \theta\}
\]
of angle $2 \theta$ for some $\theta \xeof (0,\frac{\pi}{2}]$.
A semigroup {\SG} is called \textit{analytic} if

\begin{itemize}[leftmargin=12pt]
\item[1.]
there is $\theta \xeof (0,\frac{\pi}{2})$
and an analytic extension $\linS':\Sigma_{\,\theta} \xto \BO{X}$ of {\SG},
which is \textit{locally bounded},
that is,
\[
\sup \,
\{ \, \norm{\linS'(z)} \: z \in \Sigma_{\,\theta}, \, \norm{z} \leq 1 \}
\, < \, \infty,
\]
\item[2.]
the semigroup property $S(t{+}s) \eq S(t) {\circ} S(s)$ holds
for $s,t \xeof \Sigma_{\,\theta}$.
\end{itemize}

\subparagraph{} %- properties of analytic semigroups
The advantage of working with this class of semigroups
instead of the strongly continuous one is given by the following features
of an analytic semigroup.

\begin{itemize}[leftmargin=4mm]
\item[1.]
$\linA x \xeof \domain{\linA}$ for each $x \xeof \domain{\linA}$.

\item[2.]
The orbits $\eta_x$ $(x \xeof \banX)$ of an analytic semigroup
are indefinitely differentiable.
\end{itemize}

\subparagraph{} %- properties of analytic semigroups
Hence,
if $\linA$ is the generator of an analytic semigroup,
then the solutions $x$ of the the differential equation
$\dot{x} + \linA x \eq y$ are smooth in all intervals $(0,\T)$,
and by this reason the differential equation is called \textit{parabolic}.

\subsection{} %-- generator theorems for analytic semigroups
%------------------------------------------------------------------------------%
Remember that the functional calculus for a sectorial operator $\linA$
enables to define the semigroup $\{e^{-tA}\}_{t \geq 0}$
by a contour integral around the spectrum $\sigma(\linA)$.
But more is true in this case:

\begin{theorem}\label{SEMIGROUP_CREATED_BY_SECTORIAL_OPERATOR_IS_ANALYTIC}
$\linA$ is sectorial,
iff $\{e^{-tA}\}_{t \geq 0}$ is {analytic}.
\end{theorem}

\subparagraph{} %- sectorially contractive analytic semigroups
Let us now consider the case,
where the analytic semigroup {\SG} is contractive
in a sector $\Sigma_{\,\theta}$,
that is,
$\norm{\linS(\zeta)} \leq 1$ holds for all $\zeta \xeof \Sigma_{\,\theta}$,
and the semigroup is called \textit{sectorially contractive analytic}
of angle $\theta$,

\paragraph{}
%------------------------------------------------------------------------------%
The generation theorem for the class of sectorially contractive 
analytic semigroups reads as follows:

\begin{theorem}\label{CRITERION_FOR_SEMIGROUP_TO_BE_SECTORIALLY_CONTRACTIVE}
Let $\linA$ be densely defined and let $\theta \xeof (0,\frac{\pi}{2}]$.
Then the following assertions are equivalent:

\begin{itemize}[leftmargin=9mm]
\item[(a)]
${-}\linA$ generates a sectorially contractive analytic $C_0$-semigroup
(of angle $\theta$).

\item[(b)]
$-\Sigma_{\,\theta} \xsubset \rho(\linA)$
and $\norm{R(\linA,\zeta)} \leq \abs{\zeta} ^{-1}$
for all $\zeta \xeof -\Sigma(\,\theta)$.
\end{itemize}
\end{theorem}

\subsection{} %-- analytic semigroups created by quasi-coercive forms
%------------------------------------------------------------------------------%
Recall from section 10
that linear operators associated with quasi-coercive forms are m-accretive,
and in this case these operators are called regularly accretive.
Rather similar,
one of the main results in this section consists of the fact
that quasi-coercive sesquilinear forms give rise to analytic semigroups
and their (sectorial) infinitesimal generators are the operators
associated with these forms:

\begin{theorem}
The operator $\linA_{\sql}$ associated with a quasi-coercive form $\sql$
in a Hilbert space triple $(\vecV,\hilH,\vecV_a^*)$ is sectorial.
\end{theorem}

\subparagraph{}
In fact,
let $\sql$ be a coercive form with lower bound $\theta$,
that is,
$\re \sql(u,u) \geq \theta \norm{u}_V^2$,
and let $\lambda \xeof \C$ such that $\re \lambda \geq 0$.
Then the form $\sql(u,v) + \lambda \scalar{u}{v}$ is also coercive,
and the associated operator $\linA {+} \lambda \: \domain{\linA} \xto \hilH$
proves to be bijective.
This implies that for $f \xeof \hilH$ there is $u \xeof \domain{\linA}$
such that $\linA u + \lambda u \eq f$ and
\[
\sql_{\lambda}(u,v) \eqdef
\sql(u,v) + \lambda \scalar{u}{v}_H \= \scalar{f}{v}_H.
\]
Letting $v=u$, we especially have
\[
   \sql_{\lambda}(u,u)
\= \sql(u,u) + \lambda \norm{u}_{\hilH}^2
\= \scalar{f}{u}_{\hilH}.
\]
A first estimation gives
\[
\abs{\lambda} \norm{u}_H^2 \leq \Theta \norm{u}_V^2 + \norm{f}_H \norm{u}_H,
\]
while a second one, using the coercivity estimation, yields
\[
\theta \norm{u}_V^2 \leq \sql(u,u) \= \scalar{f}{u}_H - \lambda \norm{u}_H^2
\leq \norm{f}_H \norm{u}_H.
\]
Combining both estimates, we obtain
\[
\abs{\lambda} \norm{u}_H^2 \leq
(1 + \frac{\Theta}{\theta}) \norm{f}_H \norm{u}_H,
\]
hence
\[
\snorm {\lambda (\linA+\lambda)^{-1}f} \leq
(1 + \frac{\Theta}{\theta}) \norm{f}
\]
and $\norm{\lambda(\linA+\lambda)^{-1}} \leq M := 1 + \Theta/\theta$
for all $\lambda \xeof \C$ with $\re \lambda \geq 0$.
This shows that the suffcient criterion for sectorialness of $\linA$
as quoted in theorem \ref{CRITERION_FOR_SECTORIALNESS} is fulfilled.

%------------------------------------------------------------------------------%
\section{Applications to Linear Differential Operators}                        %
%------------------------------------------------------------------------------%
%-- 
%-- 12. APPLICATIONS TO LINEAR DIFFERENTIAL OPERATORS
%--     ---------------------------------------------
This final section deals with applications of abstract functional analysis
to partial differential operators and the problems associated with them.
We want to start with the concept
of \textit{linear partial differential operator}
and continue with the consideration of specific problems
associated with this kind of operators.

\subsection{} %-- linear partial differential operators
%------------------------------------------------------------------------------%
Remember first the notation of \textit{multi-index},
which is nothing but a $d$-tuple $\alpha \eq (\alpha_1,\ldots,\alpha_d)$
with values $\alpha_i \xeof \N_0$
and with modulus $\abs{\alpha} := \alpha_1 + \ldots + \alpha_d$.
Multi-indices are used to shorten multiple partial derivatives
$
\PA f \eqdef \partial_1^{\alpha_1}
             \partial_2^{\alpha_2} \ldots
             \partial_d^{\alpha_d} f
$
of a sufficiently often differentiable and complex-valued function $f$,
but also to denote the multiple powers
$\xi^{\alpha} := \xi_1^{\alpha_1}
\cdot \xi_2^{\alpha_2} \cdot \ldots \cdot \xi_d^{\alpha_d}$
of a vector $\,\xi \eq (\xi_1,\xi_2,\ldots,\xi_d) \xeof \C^d$.

\subparagraph{} %- linear differential operators
By the usage of multi-indices,
a \textit{linear partial differential operator} $\A$ of even order $2m$
is defined by
\begin{equation}\label{EVEN_ORDER_LPDO}
\A(x,\,\cdot\,) \eqdef
\sum_{\abs{\alpha} \leq 2m} c_{\alpha}(x) \, \PA
\end{equation}
for $x \xeof \clG$,
the closure of a \textit{domain} $\G$,
that is,
an open and connected subset of an Euclidean space $\Rd$.
This domain may be bounded or unbounded and even the whole of $\Rd$.

\subparagraph{} %- essentially real differential operators
We shall call $\A$ \textit{essentially real},
if the coefficients $c_{\alpha}$ of highest order $\abs{\alpha} \eq 2m$
are real-valued.
Moreover,
the expression $\A$ is said to be
an \textit{operator with constant coefficients},
if all functions $c_{\alpha}$ are constant in $\clG$.

\subsection{} %-- free space and boundary value problems
%------------------------------------------------------------------------------%
In case of $\G \eq \Rd$ respectively $\bndG \eq \emptyset$,
where $\bndG := \overline{\G} \setminus \interior{\G}$ is the boundary of $\G$,
we denote the linear partial differential equation $\A \, u \eq f$
with $\A$ given by \eqref{EVEN_ORDER_LPDO} a \textit{free space problem}.
On the other hand,
partial differential equations considered in domains
with non-empty boundary are usually furnished with \textit{boundary conditions},
which fix possible solutions and their normal derivatives up to a certain order
in the boundary $\bndG$.
Thus,
when having in mind linear equations of the form $\A\,u \eq f$
with imposed boundary conditions,
we shall speak of \textit{linear boundary value problems}.

\subparagraph{} %- general form of boundary value problems
The most general form of a boundary value problem
posed in a domain $\G \subsetneq \Rd$ is of the form
\begin{equation}\label{LEBVP}
\begin{cases}[1.2]
\quad \A(x,u)     \= f(x)   & (x \xeof \G),\\
\quad \BDO_k(x,u) \= g_k(x) & (x \xeof \bndG \and k \eq 0,\ldots,m{-}1),\\
\end{cases}
\end{equation}
where $\A$ is a {linear} differential operator of order $2m$
and $\BDO_0,\ldots,\BDO_{m-1}$ are \textit{boundary operators} of order $m_k$,
each of them defined by
\[
\BDO_k(x,u) \=
\sum_{\abs{\alpha} \leq m_k} b_{k,\alpha}(x) \, \PA u(x).
\]

\subparagraph{}
To introduce the boundary operators $\BDO_k$,
it is sufficent
that the real or complex-valued functions $b_{k,\alpha}$
are defined in a neighbourhood of the boundary $\bndG$ only.
Then each operator $\BDO_k$ is a linear differential operator in its own right
of order $m_k {\,\leq\,} 2m{-}1$,
which maps the function $u \: \clG \xto \C$
to a function $\BDO_k u \: \bndG \xto \C$
belonging to a so-called \textit{boundary space} $B_k(\bndG)$.

\paragraph{} %- homogeneous and normal boundary operators
%------------------------------------------------------------------------------%
The collection $\{B_k\}_{k=0}^{m-1}$ of boundary operators in \eqref{LEBVP}
is said to be \textit{homogeneous}
if $\BDO_k(x,\cdot) \eq 0$ for $x \xeof \bndG$ and $k \eq 0,\ldots,m{-}1$.
Moreover,
it is denoted \textit{normal} if

\begin{itemize}[leftmargin=15mm,itemsep=2pt]
\item[(B-1)]
$m_i \neq m_j$ for $i \neq j$,

\item[(B-2)]
for all $\xi \neq 0$ being normally positioned in $x \xeof \bndG$
we have
\[
\sum_{\abs{\alpha} = m_k} b_{k,\alpha}(x) \, \xi^{\alpha} \neq 0
\quad \quad (k=0,\ldots,m{-}1).
\]
\end{itemize}

\subparagraph{} %- homogeneous Dirichlet systems
Finally,
our collection of boundary operators is called a \textit{Dirichlet system},
if it is normal and the values $m_k$ form a permutation
of the set $\{0,\ldots,m{-}1\}$,
so one may reorder the operators $\BDO_k$ to obtain $m_k \eq k$
($0 \leq k \leq m{-}1$).
The most regarded set of boundary conditions in applications
is the \textit{homogeneous Dirichlet system} defined by
\begin{equation}\label{HOMOGENEOUS_DIRICHLET_SYSTEM}
\BDO_k(x,u) \eqdef (\partial_{\nu}^k u)(x) \= 0
\end{equation}
for $x \xeof \bndG$ and $k \eq 0,\ldots,m{-}1$,
where $(\partial_{\nu}^k u)(x)$ is the $k$-th derivative of a function $u$
along the outer normal vector $\nu$ in $x \xeof \bndG$.

\subparagraph{} %- classical solutions of BVP endowed with a Dirichlet system
If the set of boundary conditions in \eqref{LEBVP} is a Dirichlet system,
then we denote a solution $u \: \clG \xto \C$
of this problem \textit{classical},
if
\begin{itemize}[leftmargin=9mm]
\item[(a)] $u \xeof C^{2m}(\G) \, \cap \, C^{m-1}(\clG)$,
\item[(b)] $u$ fullfills $\A(x,u) \eq f(x)$ in each $x \xeof G$ and
\item[(c)] the $m{-}1$ equations $\BDO_k(x,u) \eq g_k(x)$ hold
           in every $x \xeof \bndG$.
\end{itemize}

\paragraph{} %- boundary condition for second-order problems
%------------------------------------------------------------------------------%
Especially to each \textit{second-order} boundary value problem
there is assigned exactly one boundary operator $\BDO_0$,
hence problem \eqref{LEBVP} becomes into
\begin{equation}\label{SECOND_ORDER_BOUNDARY_VALUE_PROBLEM}
\begin{cases}[1.2]
\quad \A(x,u)     = f(x)   & (x \xeof \G),\\
\quad \BDO_0(x,u) = g_0(x) & (x \xeof \bndG).\\
\end{cases}
\end{equation}

\subparagraph{} %- enumeration of boundary conditions for second-order problems
Boundary value problems of the form \eqref{SECOND_ORDER_BOUNDARY_VALUE_PROBLEM}
and governed by a linear differential operator $\A$ of second order
are obtained by various boundary operators $\BDO_0$.
We denote

\begin{itemize}[leftmargin=15mm,itemsep=2pt]
\item[(BC-1)]
\;$\BDO_0(x,u) \eq u(x)$ for $x \xeof \bndG$ a \textit{Dirichlet condition},

\item[(BC-2)]
\;$\BDO_0(x,u) \eq u_{\nu}(x)$ for $x \xeof \bndG$ a \textit{Neumann condition},

\item[(BC-3)]
\;$\BDO_0(x,u) \eq \alpha(x)u(x)+\beta(x) u_{\nu}(x) \eq 1$ for $x \xeof \bndG$ 
a \textit{Robin condition},
where the functions $\alpha, \beta$ must fulfill 
\[
\alpha(x)^2 + \beta(x)^2 \neq 0 \quad \quad (x \in \bndG).
\]
\end{itemize}

\subparagraph{}
Due to physical reasons,
one often has to split the boundary set $\bndG$
into disjoint subsets $\Gamma_1,\Gamma_2,\Gamma_3$,
on which the conditions (BC-1), (BC-2) and (BC-3) are prescribed.
We then shall speak of \textit{mixed} boundary value problems
for the second-order operator $\A$.

\subsection{} %-- operators of elliptic, parabolic and hyperbolic type
%------------------------------------------------------------------------------%
The concept of \textit{ellipticity} has been shown to be helpful in order
to classify linear partial differential expressions
by a simple algebraic criterion.
Let $\A$ be a linear differential operator of order $2m$ be given
by \eqref{EVEN_ORDER_LPDO} and let
\[
p(x,\xi) \eqdef \sum_{\abs{\alpha} = 2m} c_{\alpha}(x) \, \xi^{\alpha}
\quad \for \xi = (\xi_1,\ldots,\xi_d) \in \Rd
\]
be the \textit{characteristic function} or the \textit{main symbol} of $\A$,
which is built up by the summands of highest order in \eqref{EVEN_ORDER_LPDO}.

\begin{itemize}[leftmargin=12mm]
\item[(E-1)]
The expression $\A$ is called \textit{elliptic} in $x \xeof \clG$,
if
\[
\abs{p(x,\xi)} \, \neq \, 0 \qfa \xi \neq 0,
\]
and $\A$ is said to be \textit{uniformly elliptic} in $\G$ (in $\clG$),
if it is elliptic in \textit{each} point $x \xeof \G$ (in $x \xeof \clG$).
We then call a boundary value problem $\A \,u \eq f$ \textit{elliptic},
if $\A$ is so.

\item[(E-2)]
$\A$ is called \textit{properly elliptic} in $\clG$,
if it is elliptic and 
if for each fixed $x \xeof \bndG$ and all linear independent vectors
$\xi, \eta \xeof \Rd$ the main symbol $P_{\A}(x, \xi + \tau \eta)$ has,
{\wrt} the variable $\tau$,
exactly $m$ roots with positive imaginary part and
exactly $m$ roots with negative imaginary part.
It can be shown that
if $\A$ is \textit{essentially real} or $\G \xsubset \Rd$ with $d \geq 3$,
then any elliptic differential operator is also {properly elliptic}.

\item[(E-3)]
The concept of ellipticity is often not good enough
to obtain satisfying results.
Hence the expression $\A$ of order $2m$
is called \textit{strongly elliptic} in $x \xeof \clG$,
if there is $\theta_x> 0$ such that
\[
\re \, (-1)^m \, p(x,\xi)
\, \geq \,
\theta_x \cdot \abs{\xi} ^{2m}
\quad
\fa \xi \neq 0.
\]

\item[(E-4)]
$\A$ is called \textit{uniformly strongly elliptic} in $\clG$,
if there is $\theta {\,>\,} 0$ such that 
\[
\re \, (-1)^m \, p(x,\xi)
\, \geq \,
\theta \cdot \abs{\xi} ^{2m}
\quad
\fa \xi \neq 0 \and x \xeof \clG.
\]
\end{itemize}

\subparagraph{} %- Laplace operator
The most well-known linear partial differential operator of second order
is the \textit{Laplacian}
\[
\Delta = \partial^2 / \partial {x_1}^2 + \ldots
       + \partial^2 / \partial {x_d}^2,
\]
whose negative version $-\Delta$ constitutes the prototype
of what is meant by an {uniformly strongly elliptic operator}.
One easily computes
$\re -p_{-\Delta}(x,\xi) \eq \abs{\xi} ^2$ for $\xi \xeof \Rd$,
hence the expression $-\Delta$ often serves as a first example
when studying linear-elliptic operators.

\paragraph{} %- parabolic and hyperbolic differential operators
%------------------------------------------------------------------------------%
Besides the class of {linear-elliptic} operators
there are two further important classes
of linear partial differential operators.
They are called to be of \textit{parabolic} and \textit{hyperbolic} type.
In order to introduce these concepts,
we consider differential expressions on cylindrical subsets
$\J \xtimes \G$,
where $\J \eq (0,\T) \xsubset \R$ and $\G \xsubset \Rd$ is a domain.
Then a \textit{parabolic} differential operator $\A$ is one of the form
\[
\A(t,x,u) \= \frac{\partial u}{\partial t}(x) \+ \A_1(x,u),
\]
where $\A_1$ is \textit{uniformly elliptic} in $\G$,
and $\A$ is of \textit{hyperbolic} type,
if it is of the form
\[
\A(t,x,u) \= \frac{\partial^2 u}{\partial t^2}(x) \+ \A_1(x,u),
\]
where again $\A_1$ is uniformly elliptic in $\G$.
Simple examples of para\-bolic and hyper\-bolic operators are
the \textit{heat operator} $\partial_t    - \Delta$ and
the \textit{wave operator} $\partial_{tt} - \Delta$,
respectively.
With these operators we associate the linear evolution equations
\[
\frac{\partial u}{\partial t} + \A(\cdot,u) = f
\quad \mbox{ and } \quad
\frac{\partial^2 u}{\partial t^2} + \A(\cdot,u) = f,
\]
both of them called \textit{initial value problems},
when considered with the prescribed start value $u(0) \eq u_0$
(and in addition $\dot{u}(0) \eq u_1$ in the hyperbolic case).

\subparagraph{} %- initial value problems
When considering parabolic or hyperbolic equations,
we speak of so-called \textit{initial-boundary value problems},
if for every time $t \xeof (0,T]$ the solution fulfills
a certain set of boundary conditions in every $x \xeof \bndG$,
where we especially have in mind the homogeneous Dirichlet boundary conditions
\eqref{HOMOGENEOUS_DIRICHLET_SYSTEM}.
However,
in case of $\G \eq \Rd$,
we obtain free space problems or pure initial value problems
for the differential operators of parabolic or hyperbolic type.

\paragraph{} %- overview on the section
%------------------------------------------------------------------------------%
It is the aim of this section to apply
all the presented functional analytic tools to the just enumerated problems.
Let $\A$ be a linear partial differential operator of order $k \xeof \N$
(not necessarily even) and given by
\begin{equation}\label{ANY_ORDER_LPDO}
\A(x,u) \eqdef
\sum_{\abs{\alpha} \leq k} c_{\alpha}(x) \, \PA u.
\end{equation}
Then by a \textit{realization} of $\A$ in a Banach space $\mathbb{E}(\G)$
of functions $u \: \G \xto \C$ is meant any linear operator
\[
\linA: \domain{\linA} \subset \mathbb{E}(G) \xto \mathbb{E}(G)
\]
with domain of definition $\domain{\linA}$.
The Banach space $\mathbb{E}(\G)$ should contain the space $\TF$ 
of test functions\footnote{
\,The space $\TF$ consists of indefinitely differentiable functions
having \textit{compact} support
\[
\supp{\,\phi} \eqdef \overline{\set{x \in \G}{\phi(x) \neq 0}}\,,
\]
usually denoted the \textit{space of test functions}.
This space plays an out\-standing role in analysis,
since it lies dense in a multiplicity of other functional spaces.}
as a dense linear subspace.
We also require for each $u \xeof \domain{\linA}$ that
first $\A \, u$ is well-defined by
$\linA u \eq \A \, u$ and second
$\A \, u \xeof \mathbb{E}(\G)$ holds.
Clearly,
our first condition means that the set $\domain{\linA}$
may contain only functions $u \: \G \xto \C$ being sufficiently often
differentiable in some sense.
But $\domain{\linA}$ should also include the space $\TF$ of test functions,
so $\domain{\linA}$ lies dense in $\mathbb{E}(\G)$ as well 
and $\linA$ appears as a densely defined operator in $\mathbb{E}$.

\paragraph{} %- content of the section
%------------------------------------------------------------------------------%
In the second part of this section we shall point out that

\begin{itemize}[leftmargin=2mm,itemsep=4pt]
\item[-] 
for a linear partial differential expression $\A$
there can be defined a couple of densely defined and closable linear operators
$\linA$ acting either in the {\HOELDER} spaces $C^{k,\gamma}(\G)$
($k \xeof \N_0$, $0 < \gamma \leq 1$)
or in the Lebesgue spaces $\Lp$ ($1 \leq p < \infty$);

\item[-] 
using elliptic estimates,
the realization of a strongly elliptic differentiable operator
acting in Sobolev spaces
under homogeneous Dirichlet boundary conditions is closed;

\item[-] 
the study on realizations of differential operators can be supported
by that on their adjoints,
where those are determined by so-called \textit{adjointness theorems};

\item[-] 
by starting with quasi-coercive sesquilinear forms,
created by strongly elliptic differential expressions,
we can obtain realizations in $\Lh$,
which prove to be densely defined and closed;

\item[-] 
weak solutions of initial-boundary value problems exist and are unique,
supposed the governing sesquilinear form is the Dirichlet form
of a strongly elliptic differential operator;

\item[-] 
the realizations of a regularly elliptic expression $\A$ in $\Lh$ are Fredholm,
so the alternative theorem \eqref{FREDHOLM_ALTERNATIVE_THEOREM}
may be applied to the equation $\A \, u \eq f$;

\item[-] 
strongly elliptic expressions give rise to m-accretive and 
sectorial realizations in $\Lh$,
so the related parabolic and hyperbolic initial value problem
prove to be well-posed;

\item[-] 
formally self-adjoint and half-bounded differential expressions
lead to self-adjoint realizations in $\Lh$;
\end{itemize}

\subparagraph{}
This enumeration is of course only a small selection 
within the totality of applications of functional analytic methods
to linear partial differential operators and related equations,
but may give a good impression of their capabilities.

\def\TITLEA{Realizations in the spaces $C^{\gamma}(\G)$ and $\Lp$.}
\def\TITLEB{The closeness of strongly elliptic Dirichlet operators.}
\def\TITLEC{Formally adjoint expressions and adjointness theorems.}
\def\TITLED{{\GARDING}'s inequality and quasi-coercive Dirichlet forms.}
\def\TITLEE{Weak solutions of boundary and initial-boundary value problems.}
\def\TITLEF{The Fredholm property of elliptic operators.}
\def\TITLEG{Spectral theory and sectorialness of strongly elliptic operators.}
\def\TITLEH{Self-adjoint extensions of half-bounded operators.}

%-- realizations of differential equations in C^{k,\gamma}(G) and L^2(G)
%------------------------------------------------------------------------------%
\subsubsection{I. \TITLEA}
We basically want to distinguish between two classes of functional spaces,
in which realizations of a linear differential operator $\A$ are studied:

\begin{itemize}[leftmargin=4mm]
\item[1.] %- realizations in Hölder spaces
The notation of classical derivative $\PA u$,
applied to a function $u \: \G \xto \C$,
leads us to the usage
of the spaces $C^{\gamma}(\G)$ of \textit{{\HOELDER}-constinuous} functions,
which constitute Banach spaces when endowed with a suitable norm.

\subparagraph{} %- Hölder continuity
Preliminary,
a function $u \: \G \xto \F$ is called \textit{{\HOELDER}-continuous}
{\wrt} the parameter $\gamma \xeof (0,1]$,
if its \textit{{\HOELDER}-norm} is finite,
that is,
\[
\abs{u} _{\,\gamma} \eqdef
\sup_{x,y \, \in \, \G, \, x \neq y} \,
\frac{\abs{u(x) - u(y)} } {\abs{x-y} ^{\gamma}} \, < \, \infty.
\]

\subparagraph{} %- Hölder spaces
Let $C^{k,\gamma}(\clG)$ be the linear space
of functions $u \xeof C^k(\clG)$,
whose partial derivatives $\PA u$
are {\HOELDER}-continuous for all $\abs{\alpha} \eq k$.
Then the \textit{{\HOELDER} norm} $\abs{u} _{k,\gamma}$
of a function $u \xeof C^{k,\gamma}(\clG)$ is defined
by the sum of the norm
$\norm{u}_k := \sum_{\abs{\alpha} \leq k} \,
\sup_{x \xeof \clG } \; \abs{\partial^{\alpha} u(x)} $
and the {\HOELDER} norms of the partial derivatives $\PA u$
for $\abs{\alpha} \eq k$:
\[
\abs{u} _{k,\gamma}
\eqdef
\norm{u}_k \+ \sum_{\abs{\alpha} \eq k} \abs{\,\PA u} _\gamma.
\]
The spaces $C^{k,\gamma}(\clG)$ prove to be complete {\wrt} this norm,
but they are not reflexive.

\subparagraph{}
A realization of a differential operator $\A$ of order $2m$
acting in {\HOELDER} space $C^{\gamma}(\G)$ is nothing but a mapping
\[ 
\linA \: \domain{\linA} \subset C^{\gamma}(\G) \to C^{\gamma}(\G),
\]
where $\domain{\linA}$ forms a linear subspace of $C^{2m,\gamma}(\G)$
restricted by boundary conditions,
for example by a Dirichlet system.

\subparagraph{}
There is an extensive theory on the realizations
of linear partial differential operators acting in {\HOELDER} spaces,
which provide nice results especially for operators of second order.
We should mention \textsc{J.\,Schauder} as the outstanding pioneer
in the investigation of this subject.
But \textsc{Schauder} uses classical methods for the setup of his theory,
coming from potential theory and complex analysis,
so we do not want to go into the details
(a good reference for this subject is \cite{GT}).

\item[2.] % realizations in Lebesgue spaces
Instead of classical derivatives,
the modern theory of partial differential equations
is based on the concept of \textit{weak derivative},
since this approach leads to larger spaces of differentiable functions
than the already mentioned {\HOELDER} spaces,
usually denoted as \textit{Lebesgue-Sobolev spaces}.

\subparagraph{} %- Hilbert space L^2
First recall that
for any domain $\G \xsubset \Rd$
the Hilbert space $\Lh$ is a complete normed space
consisting of measurable functions $f \: \G \xto \C$ such that
\[
\int_{\,\G} \, \abs{f} ^2 \, dx \, < \, \infty,
\]
where the inner product in $\Lh$ is defined by
\[
\scalar{f}{g} \eqdef \int_{\,\G} \, f {\cdot} \overline{g} \, dx
\quad \for f,g \in \Lh.
\]
This space is a specific member in the scale of the so-called
\textit{Lebesgue spaces} $\Lp$ ($1 \leq p < \infty$),
whose norms are defined by
\[
\abs{f} _p \eqdef ( \int_{\,\G} \, \abs{f} ^p \, dx)^{\frac{1}{p}} \,<\,\infty
\]
for equivalence classes of measurable functions $f \xeof \Lp$.

\subparagraph{}
The space $L^{\infty}(\G)$,
however,
consists of those equivalence classes of 
measurable functions $f \: \G \xto \C$,
for which 
\[
\norm{f}_{\infty} \eqdef \inf
\{\,a \geq 0: \abs{f} \leq a \,\, \mu \mbox{-almost everywhere}\,\}
\,<\,\infty.
\]

\paragraph{} %- weak derivatives
%------------------------------------------------------------------------------%
We next turn to the notation of \textit{weak derivative}.
Preliminary,
any \textit{locally integrable function}\footnote{
\,A function $u\:\G \xto \C$ is said to be \textit{locally integrable}
and to belong to the space $\Lc$,
if $u$ is Lebesgue-integrable on each compact subset of $\G$.}
$u:\G \xto \C$ is said to be \textit{weakly differentiable} 
{\wrt} a given multi-index $\alpha$,
if there is a further locally integrable function $v$ such that
\[
\int_{\,\G} v(x) \, \phi(x) \, dx
=
(-1)^{ \abs{\alpha} } \, \int_{\,\G} u(x) \, \PA \phi(x) \, dx
\fa \phi \in \TF.
\]
We denote $\PA u := v$ as the $\alpha$-\textit{weak derivative} of $u$
in this case.
The function $u$ is said to be \textit{$k$-times weakly differentiable},
if $u$ has $\alpha$-weak derivatives
for every index $\alpha$ with $\abs{\alpha} \,{\leq}\, k$.
Taking weak derivatives of locally integrable functions {\wrt} the multi-index
$\alpha$ arises as a linear mapping
\[
\PA \: D_{\alpha} \subset \Lc \to \Lc,
\]
where $D_{\alpha}$ forms a linear subspace of $\Lc$.
Let us also notice that
the concept of weak derivative is compliant
with that of the classical derivative:
if a function $u$ is classical differentiable in $\G$
with classical derivative $\PA u$,
then $u$ is also weakly differentiable,
and its weak derivative is equal to $\PA u$.

\subparagraph{} %- Lebesgue-Sobolev spaces
The space $W^{k,p}(\G)$ ($1 \leq p < \infty)$
is the \textit{Lebesgue-Sobolev space}
\[
W^{k,p}(\G) \eqdef
\{\,f \in \Lp \: \PA f \xeof \Lp, \; \abs{\alpha} \leq k \, \}.
\]
If $p \eq 2$,
the space $H^k(\G) := W^{k,2}(\G)$ can be furnished with inner product
\[
\scalar{f}{g}_k \eqdef
\sum_{\abs{\alpha} \leq k} \int_{\,\G} \;
\PA f \cdot \PA \overline{g} \, dx\,,
\quad \quad f,g \xeof H^k(\G)
\]
and norm
$\abs{f} _k \eq \scalar{f}{f}^{1/2}$,
so $H^k(\G)$ forms a Hilbert space in its own right.
In other words,
the space $H^k(\G)$ consists of square-integrable functions,
which are at least $k$-times {weakly differentiable}
and all these weak derivatives belong to $\Lh$ again.

\subparagraph{} %- Hilbert-Sobolev spaces with homogeneous boundary conditions
We also want to consider the linear subspace $\WN{k,p}{\G}$ of $W^{k,p}(\G)$
fulfilling homogeneous Dirichlet boundary conditions.
This space consists of those functions $u \xeof W^{k,p}(\G)$,
which are approximated by a sequence $\{\phi_n\}_{n \xeof \N}$
of \textit{test functions} {\wrt} the norm $\abs{\,\cdot\,} _p$,
that is,
each $\phi_n$ belongs to the space $\TF$ and $\abs{u - \phi_n} _p \xto 0$. 
Clearly,
since $\WN{k,p}{\G}$ is a closed subspace of $W^{k,p}(\G)$,
it forms a Banach space again.
Furthermore,
functions belonging to $\WN{k,p}{\G}$ fulfill the $k$ boundary conditions
of a homogeneous Dirichlet system.
For $p \eq 2$,
however,
the space $\HN{k}{\G}$ owns the inner product and norm
as defined in the larger space $H^k(\G)$,
so it is a Hilbert space as well.

\subparagraph{}
Finally,
a realization of a linear differential operator $\A$ of order $k$
in Lebesgue space $\Lp$ is any linear mapping
\[
\linA \: \domain{\linA} \subset \Lp \to \Lp,
\]
for which $\domain{\linA}$ forms a linear subspace of $W^{k,p}(\G)$,
for example restricted by homogeneous Dirichlet boundary conditions.
\end{itemize}

%-- the closeness of strongly elliptic Dirichlet operators
%------------------------------------------------------------------------------%
\subsubsection{II. \TITLEB}
One often has to deal with the realization $\linA_D$
of a strongly elliptic differential operator $\A$ in $\Lp$
having order $2m$ and domain of definition
\[
\domain{\linA_D} \eqdef \{ u \xeof \WN{m,p}{\G} \: \A\,u \xeof \Lp \},
\]
so the functions in $\domain{\linA_D}$ satisfy all the $m$ boundary conditions
of the homogeneous Dirichlet system \eqref{HOMOGENEOUS_DIRICHLET_SYSTEM}.
By this reason it is reasonable to denote the realization $\linA_D$
a \textit{Dirichlet operator}.

\paragraph{} %- closeness of strongly elliptic differential operators
%------------------------------------------------------------------------------%
We turn to the question of closeness of the realization $\linA_D$.
Let the domain $\G$ and the expression $\A$ fulfill the following conditions:

\begin{itemize}[leftmargin=9mm]
\item[(I)\;\;]
$\G$ is bounded in $\Rd$,

\item[(I')\;\,]
$\G$ satisfies the $m$-th \textit{extension property},
that is,
the linear space $R_0(\G)$ consisting
of those functions $\phi \xeof W^{m,p}(\G)$,
which are restrictions $\phi |_{\G}$
of functions $\phi \xeof C_c^{\infty}(\Rd)$,
is dense in $W^{m,p}(\G)$;

\item[(II)\;]
$\A$ is uniformly strongly elliptic in $\G$;

\item[(III)\,]
the coefficient functions $a_{\alpha \beta}$ are essentially bounded in $\G$,
that is,
$a_{\alpha \beta} \xeof L^{\infty}(\G)$;

\item[(III')]
the coefficient functions $a_{\alpha \beta}$ of highest order
$\abs{\alpha} + \abs{\beta} \eq 2m$ are uniformly continuous in the closure
$\clG$ of $\G$.
\end{itemize}

Then,
by showing the validness of the \textit{elliptic estimate}
\[
\abs{u} _{2m} \, \leq \, c \cdot \{ \,\abs{u} _p + \abs{\linA_D u} _p \, \}
\]
for $u \xeof \domain{\linA_D}$,
where $\abs{u} _p \eq \norm{u}$ is the norm of $u$ in $\Lp$,
theorem \ref{BROWDER_THEOREM} implies that the operator $\linA_D$ is closed.
Moreover,
$\linA_D$ is injective and the range of $\linA_D$ is a closed subspace in $\Lp$.

%-- formally adjoint expressions and adjointness theorems
%------------------------------------------------------------------------------%
\subsubsection{III. \TITLEC}
As we have already mentioned,
it is often helpful to combine the study of a linear operator
with that of its formally adjoint.
Supposed the linear partial differential operator $\A$ of order $k$
has coefficients $c_{\alpha} \xeof C^{\abs{\alpha} }(\G)$.
We then find for $u \xeof \TF$ and $\psi \xeof \TF$
by repeated partial integration
\[
\int_{\,\G} \sum_{\abs{\alpha} \leq k}
(c_{\alpha} \, \PA u) \cdot \overline{\psi} \, dx
\=
\int_{\,\G} u \cdot \sum_{\abs{\alpha} \leq k}
(-1)^{\abs{\alpha} } \, \overline{\PA(\overline{c_{\alpha}} \, \psi)} \, dx.
\]
Hence there is a \textit{formally adjoint} differential expression
\begin{equation}\label{FAO}
\A^{\dag}(x,u) \eqdef
\sum_{\abs{\alpha} \leq k} (-1)^{\abs{\alpha} }
\PA(\overline{c_{\alpha}(x)} \, u(x))
\end{equation}
such that the identity
\[
\int_{\,\G} \A \, \phi \, \overline{\psi} \, dx 
\=
\int_{\,\G} \phi \, \overline{\A^{\dag} \, \psi} \, dx 
\]
holds for $\phi,\psi \xeof \TF$,
briefly
$
\scalar{\A\,\phi}{\psi} \eq \scalar{\phi}{\A^{\dag}\,\psi}
$
by means of the inner product in $\TF$.
Using \textsc{Leibniz}'s rule to compute the terms
$\PA(\overline{c_{\alpha}} \, u)$ in \eqref{FAO},
there are functions $c_{\alpha}'\: \G \xto \C$
such that
\[
\A^{\dag}(x,u) \=
\sum_{\abs{\alpha} \leq k} c_{\alpha}'(x) \; \PA u(x).
\]
Thus $\A^{\dag}$ is uniquely determined by $\A$
and constitutes a differential expression of type \eqref{ANY_ORDER_LPDO},
too.
It may happen that $\A^{\dag} \eq \A$ holds,
and in this case
we shall call the expression $\A$ \textit{formally self-adjoint}.

\subparagraph{} %- a criterion for closability of differential operators
A useful criterion to ensure closability of a linear operator $\linA$
acting in Hilbert space $\Lh$
and induced by a differential expression $\A$ of order $k$
is given by the following theorem
(see also \cite{Yo}, p.\,78),
whose proof uses the presence of the formally adjoint expression $\A^{\dag}$:

\begin{theorem}\label{L2_CLOSABILITY_THEOREM}
Let the expression $\A$ of order $n$ be given by \eqref{ANY_ORDER_LPDO}
with coefficients $c_{\alpha} \xeof C^k(\G)$ for $\abs{\alpha} \leq k$ and
the linear operator
\[
\linA: \;
\begin{cases}[1.2]
\quad \domain{\linA} \eqdef \{u \xeof C^k(\G) \cap \Lh \:\A \, u \xeof \Lh \},\\
\quad \linA u \eqdef \A\,u \; \for u \xeof \domain{\linA}\\
\end{cases}
\]
be a realization of $\A$ in $\Lh$.
Then $\linA$ is closable.
\end{theorem}

\subparagraph{} %- description of the domain of definition of the closure
Note that the domain of definition $\domain{\linA}$
of the operator $\linA$ considered in theorem \ref{L2_CLOSABILITY_THEOREM}
is not restricted to any boundary conditions.
It should also be mentioned that
rather the closability of the differential operator $\linA$
with domain of definition $\domain{\linA}$
like in theorem \ref{L2_CLOSABILITY_THEOREM} is ensured,
it is not an easy task to give an exact description
for the domain of definition of its closure.
But if the coefficients of $\A$ are constant,
then the set $\domain{\overline{\linA} }$ can be determined
and coincides with the Hilbert-Sobolev space $H^k(\G)$.

\subparagraph{} %- minimal and maximal operators
The closure $\linA_{min} := \overline{\linA_0}$
of a closable linear differential operator $\linA_0$
with domain of definition $\domain{\linA_0} \eq \TF$
is denoted the \textit{minimal realization}
or the \textit{minimal operator} of $\A$.
The \textit{maximal operator} $\linA_{max}$ of $\A$,
however,
is by definition the adjoint of the minimal operator of $\A^{\dag}$
with domain of definition $\domain{{\linA^{\dag}}_0} \eq \TF$,
that is,
$\linA_{max} := ({\linA^{\dag}}_{min})^*$.
One always has $\linA_{min} \, {\prec} \, \linA_{max}$,
and often any operator $\linA$ fulfilling
\[
\linA_{min} \prec \linA \prec \linA_{max}
\]
is called a \textit{realization} of $\A$,
in opposite to our selected definition of this concept.
Using this definition,
we also want to consider the linear operators $\linA_{min}$ and $\linA_{max}$
as realizations of $\A$,
too.

\paragraph{} %- adjointness theorems
%------------------------------------------------------------------------------%
To the end,
the question arises whether any realization
of the formally adjoint expression $\A^{\dag}$ is in relationship
with the adjoint of some realization of $\A$.
This problem is examined in \cite{Br2} (section\,3: Adjoints)
even in the context of $L^p$-spaces.
\textsc{Browder}'s result for Hilbert spaces $\Lh$
and the Dirichlet realization $\linA_D$
of a regularly elliptic differential expression $\A$ sounds as follows
(corresponding to theorem 5, see p.\,57):

\begin{theorem}
Let the $2m$-order expression $\A$ be strongly elliptic
with sufficiently smooth coefficients
and defined on a smooth domain $\G \xsubset \Rd$
such that the formal adjoint $\A^{\dag}$ exists and its coefficients
are uniformly bounded in $\G$.

\subparagraph{} %- adjointness theorem in L^2
If $\linA$ and $\linA^{\dag}$ are the realizations of $\A$ and $\A^{\dag}$
in $\Lh$ with homo\-geneous Dirichlet boundary conditions
and domains of definition
\[
\domain{\linA} \= \domain{{\linA}^{\dag}} \eqdef H^{2m}(\G) \cap \HN{m}{\G},
\]
then $\linA^{\dag}$ is the Hilbert adjoint of ${\linA}$,
that is,
${\linA}^* \eq \linA^{\dag}$.
\end{theorem}

\subparagraph{} %- adjointness theorem in L^p
A similar result holds in the setting of the spaces $\Lp$ for $1 < p < \infty$.
As it is quoted in \cite{Pa}, lemma 7.3.4,
we have 
${\linA_p}^* \eq {{\linA}^{\dag}}_q$,
where $q \eq p/(p-1)$ respectively $\frac{1}{p} + \frac{1}{q} \eq 1$
and both operators
\[
\begin{cases}[1.2]
\quad \domain{\linA_p} \eqdef W^{2m,p}(\G) \cap \WN{m,p}{\G},\\
\quad \linA_p u \eqdef \A\,u \for u \in \domain{\linA_p},\\
\end{cases}
\]
\[
\begin{cases}[1.2]
\quad \domain{{\linA^{\dag}}_q} \eqdef W^{2m,q}(\G) \cap \WN{m,q}{\G},\\
\quad {\linA^{\dag}}_q u \eqdef \A^{\dag} u
      \for u \in \domain{{\linA^{\dag}}_q}.\\
\end{cases}
\]
are densely defined and closed.

%-- Gardings's inequality and quasi-coercive Dirichlet forms
%------------------------------------------------------------------------------%
\subsubsection{IV. \TITLED}
We next want to illustrate the usage of abstract sesquilinear
forms for the analysis of linear equations of the form $\A \, u \eq f$,
where $\A$ is a linear partial differential operator of order $2m$
and $u,f$ are real or complex-valued mappings
defined on a domain $\G \xsubset \Rd$.

\subparagraph{} %- differential expressions in divergence form
Preliminary,
a linear differential expression $\A$ of order $2m$
is said to be in \textit{divergence form},
if there are functions $a_{\alpha \beta} \xeof C^{\abs{\alpha} }(\G)$
for $\abs{\alpha} , \abs{\beta} \leq m$
such that $\A$ can be written as
\begin{equation}\label{LPDO_DIV}
\A \eqdef
\sum_{\abs{\alpha} \leq m, \abs{\beta} \leq m} (-1)^{\abs{\alpha} } \,
\PA(\,a_{\alpha \beta}(x) \, D^{\beta}).
\end{equation}

In opposite to this definition let us call an arbitary differential expression
given by \eqref{ANY_ORDER_LPDO} to be \textit{in non-divergence form}.

\subparagraph{} %- relationship between divergence and non-divergence forms
It is a matter of fact
that every expression $\A$ in divergence form
can be transformed into an expression $\A'$ in \textit{non-divergence form}.
Reversely,
a non-divergence form expression $\A'$ given by \eqref{ANY_ORDER_LPDO}
can be transformed into an expression $\A$ in divergence form \eqref{LPDO_DIV},
if $c_{\alpha} \xeof C^{2m - \abs{\alpha} }(\G)$ for $\abs{\alpha} \leq 2m$,
that is,
if the coefficients are sufficiently many differentiable in $\G$.
Comparing the expressions \eqref{ANY_ORDER_LPDO} and \eqref{LPDO_DIV},
it turns out that
$\A'$ is strongly elliptic in $\G$ \iff $\A$ is so.

\subparagraph{}
To show the purpose of a linear differential operator being in divergence form,
let the space $\TF$ of {test functions} be furnished with the inner product
\[
\scalar{\phi}{\psi} \eqdef
\int_{\,\G} \phi \, \overline{\psi} \,dx \qfa \phi,\psi \xeof \TF.
\]
Now,
applying $\A$ given in divergence form to $\phi \xeof \TF$,
then multiplying $\A \, \phi$ with $\overline{\psi} \xeof \TF$ and
finally performing $m$-times partial integration leads to
\[
\DF_{\A}(\phi,\psi) \eqdef \scalar{\A \, \phi}{\psi} \=
\sum_{\abs{\alpha} \leq m, \abs{\beta} \leq m} \,
\int_{\,\G}
a_{\alpha \beta} \; \PA \phi \; D^{\beta} \overline{\psi}
\, dx,
\]
since all the boundary terms vanish.
The mapping
\[
\DF_{\A} \eqdef \TF \xtimes \TF \xto \C
\]
is a sesquilinear form,
called the \textit{Dirichlet form} of $\A$ in $\G$.
If the coefficent functions $a_{\alpha \beta}$ of $\A$
are measurable and essentially bounded,
then the mapping $\DF_{\A}$ is continuous
and thus can continuously be extended to the product space
$\HN{m}{\G} \xtimes \HN{m}{\G}$.

\paragraph{} %- Garding's inequality
%------------------------------------------------------------------------------%
We are mainly interested in a $2m$-order linear differential operator $\A$
given in divergence form,
which together with the underlying domain $\G$ fulfills the already 
mentioned conditions (I), (I'), (II), (III) and (III').
\textsc{{\GARDING}}\textit{'s inequality} can be seen as a bridge
from the theory of linear-elliptic differential operators
to abstract functional analysis:

\begin{theorem}\label{GARDING_INEQUALITY_THEOREM}
Let the domain $\G$ fulfill (I),\,(I') 
and let the linear differential expression $\A$ of order $2m$ fulfill
(II),\,\,(III),\,(III').\\
Then there is $\lambda_0 \xeof \R$ and $a_0 {\,>\,} 0$ such that
for all $u \xeof \HN{m}{\G}$
\begin{equation}\label{GARDING_INEQUALITY}
\DF_{\A}(u,u) + \lambda_0 \norm{u}^2 =
\scalar{\,(\A+\lambda_0)\,u}{u} \geq a_0 \, {\abs{u} _m}^2.
\end{equation}
\end{theorem}

\subparagraph{}
In other words,
theorem \ref{GARDING_INEQUALITY_THEOREM} quotes that
the Dirichlet form associated with the operator $\A + \lambda_0$ is coercive,
hence the {Lax-Milgram} theorem may be applied.
It should be clear that this assertion also holds
for all $\lambda {\,>\,} \lambda_0$.

\paragraph{} %- Garding's inequality for second order operators
%------------------------------------------------------------------------------%
The most important class of linear partial differential operators
in divergence form {\wrt} applications is that of \textit{second order},
that is,
our operator $\A$ is of the form
\begin{equation}\label{GENERAL_SECOND_ORDER_DIFFERENTIAL_OPERATOR}
\A(x,\cdot) = 
-\sum_{i=1}^d \,
\partial_i \, (\sum_{j=1}^d a_{ij}(x)) 
+
\sum_{j=1}^d \, (b_j(x) \, \partial_j)
+
c(x),
\end{equation}
which is nothing but \eqref{LPDO_DIV} for $m \eq 1$.
The Dirichlet form $\DF_{\A}$ belonging to $\A$ is the sesquilinear form
\[
\DF_{\A}(u,v) \= \int_{\,\G} 
(\sum_{i,j=1}^d a_{ij} \, \partial_i u \, \partial_j \overline{v} \+ 
 \sum_{j=1}^d b_j \, \partial_j u \, \overline{v} \+
 c \,u \, \overline{v}
)
\, dx.
\]

\subparagraph{} %- coerciveness of the Dirichlet form of second-order operators
A short computation shows
that if the function $c$ is large enough in comparison
with the sum of partial derivatives of the functions $b_j$,
then $\DF_{\A}$ is coercive.
More precisely,
$\DF_{\A}$ proves to be coercive if
\begin{equation}\label{COERCIVITY_CONDITION}
c(x) \, \geq \, \frac{1}{2} \,
\sum_{k=1}^d \, \frac{\partial b_k}{\partial x_k}(x)
\end{equation}
holds in all points $x \xeof \G$.

\subparagraph{}
On the other hand,
if the matrix $A(x) \eq (a_{ij}(x))_{i,j=1}^d$ of second-order coefficients
in \eqref{GENERAL_SECOND_ORDER_DIFFERENTIAL_OPERATOR}
fails to be positive definite for all $x \xeof \G'$,
where $\G' \xsubset \G$ has positive measure,
then $\DF_{\A}$ cannot be coercive on $\HN{1}{\G}$.
Also,
if $A(x)$ is positive definite in all $x \xeof \G$
and condition \eqref{COERCIVITY_CONDITION} does not hold.
then the functions $b_j$ disturb the coercivity of $\DF_{\A}$.
But after all,
the form $\DF_{\A}$ proves at least to be quasi-coercive,
which is nothing but the assertion of \textsc{\GARDING}'s inequality.

\subparagraph{} %- Garding's inequality for expressions of second order
\textsc{\GARDING}'s inequality remains even true
for the Dirichlet forms of second-order differential operators,
if the basic conditions (I)' and (III') are dropped:

\begin{theorem}\label{GARDING_INEQUALITY_THEOREM_II}
Let $\G \xsubset \Rd$ fulfill (I) and
the differential operator $\A$ of second order and given in divergence form
fulfill (II) and (III).
Then there is $\lambda_0 \xeof \R$ and $a_0 {\,>\,} 0$ such that
\[
\DF_{\A}(u,u) \+ \lambda_0 \norm{u}^2
\=
\scalar{\,(\A {+} \lambda_0)\,u}{u}
\, \geq \,
a_0 \, {\abs{u} _1}^2
\]
for all $u \xeof \HN{1}{\G}$,
where $a_0$ is the uniform lower bound of the eigenvalues
of the matrices $A(x)$ with $x \xeof \G$.
\end{theorem}

%-- weak solutions of boundary and initial-boundary value problems
%------------------------------------------------------------------------------%
\subsubsection{V. \TITLEE}
The variational approach to stationary and evolutional boundary value problems
provides the concept of weak solution,
and in this part of the section we want to apply the abstract theory
boundary value problems governed by strictly elliptic differential operators
by quoting the related existence and uniqueness assertions.

\subparagraph{} %- the variational Dirichlet problem
We basically assume that

\begin{itemize}[leftmargin=15mm]
\item[(DP-1)]
$\A$ is a the differential operator of order $2m$,
strongly elliptic in a \textit{bounded} domain $\G \xsubset \Rd$
and given in divergence form;

\item[(DP-2)]
the coefficients $a_{\alpha \beta}$ belong to $L^{\infty}(\G)$
and the top-order coefficients are uniformly continuous in $\clG$.
\end{itemize}

\paragraph{} %- homogeneous Dirichlet problem
%------------------------------------------------------------------------------%
It is our first aim to solve the \textit{homogeneous Dirichlet problem}
\[
\begin{cases}[1.2]
\; (\A\,u)(x)\= f(x)& (x \xeof \G),\\
\; u(x) =
\partial u/ \partial \nu (x) =
\ldots = \partial u^{m{-}1} / \partial \nu^{m{-}1}(x) = 0 & (x \xeof \bndG).\\
\end{cases}
\]
Let $\vecV \eq \HN{m}{\G}$ and $\hilH \eq \Lh$,
hence the pair $(\vecV,\hilH)$ is a Hilbert triple.
By \textsc{{\GARDING}}'s inequality,
the Dirichlet form
\[
\DF_{\A}(u,v) \eqdef \scalar{\A \, u}{v} \quad \quad (u,v \in \vecV)
\]
is quasi-coercive with some $\lambda_0 \xeof \R$.
Then,
by means of the {Lax-Milgram} theorem,
we obtain the following \textit{well-posedness theorem}:

\begin{theorem}\label{DIRICHLET_EXISTENCE_UNIQUENESS_THEOREM}
For $\lambda \geq \lambda_0$, the variational equation
\[
\DF_{\A}(u,v) \+ \lambda \scalar{u}{v}
\= f(v) \quad (v \xeof \HN{m}{\G})
\]
owns for $f \xeof H^{-m}(\G) := \HN{m}{\G}^*$
a unique solution $u \xeof \HN{m}{\G}$ depending
continuously on the right-hand side $f$ and
called the \textit{weak solution} of the homogeneous Dirichlet problem.
\end{theorem}

%-- the Fredholm property of elliptic operators
%------------------------------------------------------------------------------%
\subsubsection{VI. \TITLEF}
This statement is properly not correct,
since it only proves to be true for properly elliptic differential operators
defined in a \textit{bounded} domain $\G \xsubset \Rd$
and endowed with certain boundary conditions,
for example those given by a homogeneous Dirichlet system.

\subparagraph{} %- Rellich's embedding theorem
We first want to present \textsc{Rellich}\textit{'s embedding theorem},
which holds for Sobolev spaces whose underlying domain is \textit{bounded}.
It can be seen as a key theorem in the analysis
of linear-elliptic boundary problems on such domains:

\begin{theorem}\label{RELLICH_EMBEDDING_THEOREM}
If $\G \xsubset \Rd$ is bounded and owns smooth boundary $\bndG$,
then the continuous embeddings
$\iota_k \: \HN{k}{\G} \hookrightarrow \Lh$ are compact for all $k \xeof \N$.
\end{theorem}

%\subparagraph{}
%A proof of this theorem can be found for example in \ldots.

\subparagraph{} %- elliptic operators as continuous mappings
The basic approach to investigate the Fredholm property
of a properly elliptic operator $\A$ in divergence form
and of order $2m$
consists of two small changes for the Dirichlet realization $\linA_D$
having domain of definition $\domain{\linA} \eq H^{2m}(\G) \cap \HN{m}{\G}$:
let us consider the linear operator $\linA_D$ as

\begin{itemize}[leftmargin=9mm]
\item[(F-1)]
a mapping $\linA_D \: \HN{m}{\G} \xto \Lh$,
which is continuous due to the essential boundedness
of the coefficients $a_{\alpha,\beta}$ in \eqref{LPDO_DIV};

\item[(F-2)]
the sum $\linA_1 + \linA_2$ of two further linear operators,
where we shall assume that
$\linA_1 \: \HN{m}{\G} \xto \Lh$ is the main part 
having weak derivatives of order $2m$ only,
whereas $\linA_2 \: \HN{m}{\G} \xto \Lh$ is the part of $\linA_D$
consisting of weak derivatives with order less than $2m$.
\end{itemize}

\subparagraph{}
Let $\lambda_0$ be the smallest parameter
for the Dirichlet form $\DF_{\linA+\lambda_0}$ to be coercive.
Since
\[
\linA_D \= (\linA_1 + \lambda_0) \+ (\linA_2 - \lambda_0)
\]
and $\linA_2- \lambda_0$ may be seen as a compact perturbation
of $\linA_1+\lambda_0$,
which maps $H^m(\G)$ one-to-one onto $\Lh$ and consequently is Fredholm 
with vanishing index,
the operator $\linA_D$ is Fredholm again with index $\chi(\linA_D) \eq 0$.

%-- spectral theory and sectorialness of strongly elliptic operators
%------------------------------------------------------------------------------%
\subsubsection{VII. \TITLEG}
Let us consider again a strongly elliptic differential operator $\A$
given in divergence form and its Dirichlet realization $\linA_D$
in a bounded domain $\G \xsubset \Rd$.
By {\GARDING}'s inequality,
the Dirichlet form $\DF_{\A}$ is quasi-coercive
{\wrt} the parameter $\lambda_0$,
hence $\linA_D + \lambda$ proves to be bijective
for all $\re \lambda {\,>\,} \lambda_0$.
Furthermore,
the homogeneous Dirichlet problem in $\Lh$
\begin{equation}\label{HOMOGENEOUS_DIRICHLET_PROBLEM}
\linA_D \, u  \+ \lambda u \= f \quad \quad (f \in \Lh)
\end{equation}
is well-posed,
which is a direct consequence of
theorem \ref{DIRICHLET_EXISTENCE_UNIQUENESS_THEOREM}.

\subparagraph{}
In the special case,
where $f \xeof \Lh$,
the solutions $u$ of \eqref{HOMOGENEOUS_DIRICHLET_PROBLEM}
belong to the Hilbert-Sobolev space $H^{2m}(\G)$ and are called \textit{strong}.
It should be clear that 
each strong solution of \eqref{HOMOGENEOUS_DIRICHLET_PROBLEM}
is also a weak solution.
But in order to show that these solutions $u$ are even classical,
that is,
$u \xeof C^{2m}(\G) \cap C^m(\clG)$,
methods of \textit{regularity theory} for elliptic problems have to be applied.

\paragraph{} %- spectrum of the Dirichlet operator
%------------------------------------------------------------------------------%
Let us turn to the spectrum of $\linA_D$.
Since $-\Delta_D + \lambda$ forms a bijective and closed operator
between the spaces $H^{2m}(\G) \cap \HN{m}{\G}$ and $\Lh$
for each $\lambda \xeof C$ with $\re \lambda \geq \lambda_0$,
we surely have
\[
\{ \zeta \in \C: \re \zeta \geq \lambda_0\} \xsubset \rho(\linA_D)
\and
\sigma(\linA_D) \xsubset \{ \zeta \in \C: \re \zeta < \lambda_0\}.
\]
Furthermore,
if $\G$ is bounded,
then $\linA_D$ is an operator with pure point spectrum.

\paragraph{} %- sectorialness of the Dirichlet operator
%------------------------------------------------------------------------------%
We can also make assertions on the accretivity and sectorialness
of the operator $\linA_D$.
As theorem \eqref{QUASI-COERCIVE_FORMS_AND_ASSOCIATED_OPERATORS}
(d),\,(e) shows,
$\linA_D + \lambda_0$ is the generator
of a semigroup of bounded linear operators
that is strongly continuous and analytic at the same time.
Hence the initial value problems for the heat and wave equation
governed by the operator $\linA_D + \lambda_0$ is well-posed.
By perturbation theory of semigroups,
this is also true for both problems governed by the operator $\linA_D$ itself.

\paragraph{} %- Gauss-Weierstrass semigroup
%------------------------------------------------------------------------------%
Let us finally consider the heat equation
for the negative Laplacian $-\Delta_D$ in the entire space $\Rd$,
which of course is a free space problem.
Theorem \ref{SEMIGROUP_CREATED_BY_SECTORIAL_OPERATOR_IS_ANALYTIC} ensures
that the realization $-\Delta_D$ in Hilbert space $L^2(\Rd)$
with domain of definition
\[
\domain{-\Delta_D} \= H^2(\Rd)
\]
is the infinitesimal generator
of a sectorially contractive analytic semigroup $\{W(t)\}_{t \geq 0}$,
denoted the \textit{Gauss-Weierstrass} \textit{semigroup}.
Hence there is a uniquely defined solution $u \xeof C^1(\J \times \Rd,\R)$
for the homogeneous Cauchy problem
\[
\begin{cases}[1.2]
\quad u_t - \Delta u \= 0,\\
\quad u(x,0)         \= u_0(x) \for x \in \Rd \\
\end{cases}
\]
with initial mapping $u_0 \xeof \HN{2}{\Rd}$.
Then one may represent the semigroup $\{W(t)\}_{t \geq 0}$
explicitly by means of the \textit{heat kernel} function $\Phi_d$,
that is,
for all $t > 0$ and $x \xeof \Rd$
\[
(W(t) \, u_0) \, (x)
=
\frac{1}{(4 \pi \, t)^{d/2}}
\int_{\,\Rd} e^{- \frac{\abs{x-y} ^2}{4t}} u_0(y) \, dy
=
(\Phi_d(t) \ast u_0)\,(x).
\]

%-- self-adjoint extensions of half-bounded operators
%------------------------------------------------------------------------------%
\subsubsection{VIII. \TITLEH}
Remember the definition of the differential expression $\DO^{\dag}$.
We then denote a linear differential expression $\DO$
\textit{formally self-adjoint} in case of $\DO^{\dag} \eq \DO$.
If $\DO$ has constant coefficients,
then $\A$ proves to be formally self-adjoint
iff
$c_{\alpha} \xeof \R$ for $\abs{\alpha} $ even
and $c_{\alpha} \xeof i \, \R$ for $\abs{\alpha} $ odd.

\paragraph{} %- examples of formally self-adjoint differential expressions
%------------------------------------------------------------------------------%
In the following let us collect a couple of
formally self-adjoint linear differential expressions
and describe some self-adjoint realizations of them:

\begin{itemize}[leftmargin=17pt]
\item[(a)] \textit{Sturm-Liouville operators.}
Our first example of a formally self-adjoint linear differential operator
(given in divergence form) is the \textit{Sturm-Liouville} \textit{operator}
\[
\A_{sl}(u,x) \eqdef -(p(x)\, u'(x))' \+ q(x)\, u(x),
\quad \quad x \xeof (a,b)
\]
with $p(x) {\,>\,} 0$ and $q(x) \geq 0$ for $x \xeof [a,b]$,
$-\infty {\,<\,} a {\,<\,} b {\,<\,} \infty$.
This operator determines the elongations of a vibrating string
being fixed at the endpoints of $[a,b]$.
Then the corresponding Dirichlet form reads as
\[
\DF_{\A_{sl}}(u,v) \=
\int_a^b \,
(p \, u'\, \overline{v}' \+ q \, u \, \overline{v}) \, dx,
\]
which shows that the expression $\A_{sl}$ is symmetric and half-bounded
on the domain of definition $\TF$.
The energetic space $H_{\A_{sl}}$ is the Hilbert-Sobolev space $\HN{1}{a,b}$
containing the restrictions of 
\textit{absolutely continuous functions} $u\: [a,b] \xto \R$
with boundary conditions $u(a) \eq u(b) \eq 0$ to the interval $(a,b)$.

\subparagraph{}
The self-adjoint extension $\linA_{sl,F}$
of the {Sturm-Liouville} operator in $L^2(a,b)$ has domain of definition
\[
\domain{\linA_{sl,F}} = H^2(a,b) \, \cap \, \HN{1}{a,b}.
\]
The operator $\linA_{sl,F}$ has pure point spectrum,
since the interval $(a,b)$ is assumed to be bounded.
Especially in case of $p(x) \eq 1$ and $q(x) \eq 0$ for $x \xeof [a,b]$, 
the eigen\-values and eigen\-functions
of the operator $\A_{sl,F} \eq - {d^2} / {dx^2}$ can exactly be determined.
In fact,
the eigenvalue distribution of $\A_{sl,F}$ is given by
\[
\lambda_n \= \frac{n^2 \pi^2}{(b-a)^2} \quad \quad (n \xeof \N),
\]
so the asymptotical behaviour of the eigen\-values is $\lambda_n \sim n^2$.
Each eigen\-value $\lambda_n$ owns a one-dimensional eigen\-space,
and the eigen\-function $u_n$ of $\lambda_n$ is given by
\[
u_n(x) \eqdef \sin \, (\frac{n \pi x}{b-a}) \quad \for x \in [a,b].
\]
Finally,
the sequence of the eigen\-functions $u_n$
constitutes after suitable renormations an ortho\-normal base in $L^2(a,b)$.

\item[(b)] \textit{(Negative) Laplacian in various domains $\G \xsubset \Rd$.}
In order to apply the theory of self-adjoint operators
to the investigation of the negative Laplacian,
we start with the realization $\linA_{-\Delta,0}$ of the expression $-\Delta$
in the space $\TF$ of test functions.

\subparagraph{} %- Green's formula for the negative Laplacian
Using \textsc{Green}'s second formula,
and due to the boundary condition $\phi(x) \eq 0$ for $x \xeof \partial \G$
and $\phi \xeof \TF$,
we compute for $\phi,\psi \xeof \TF$
\[
\int_{\G} -\Delta \phi \, \overline{\psi} \, dx
\=
\sum_{i=1}^n \int_{\G}
\frac{\partial^2 \phi}{\partial x_i^2} \, \overline{\psi} \, dx
\]
\[
\=
\sum_{i=1}^n \int_{\G}
u \, \overline{\frac{\partial^2 \phi}{\partial x_i^2}} \, dx 
\=
\int_{\G} \phi \, \overline{\Delta \psi} \, dx,
\]
which implies symmetry of the operator $\linA_{-\Delta,0}$.
Moreover,
we obtain for $\phi \xeof \TF$ the inequality
\[
\scalar{-\Delta \phi}{\phi} \=
\sum_{i=1}^n \int_{\G}
\frac{\partial^2 \phi}{\partial x_i^2} \, \overline{\phi} \, dx
\=
\sum_{i=1}^n \int_{\G} \abs{ \frac{\partial \phi}{\partial x_i}} ^2 \, dx 
\, \geq \, 0,
\]
hence $\linA_{-\Delta,0}$ is half-bounded with lower bound $0$
and fulfills the prerequisites
to apply the {Friedrichs} extension method.
Thus we may quote
that the operator $\linA_{-\Delta,0}$ owns at least one self-adjoint extension,
and this assertion is true actually for arbitrary domains $\G \xsubset \Rd$.

\paragraph{} % realizations of the Laplacian
%------------------------------------------------------------------------------%
In the following enumeration we consider various self-adjoint realizations
of the negative Laplacian in a domain $\G$,
all of them being extensions of the operator $\linA_{-\Delta,0}$:

\begin{itemize}[leftmargin=5mm]
\item[(i)]
\textit{(Negative) Laplacian in $\Rd$}.
The analysis of the operator $-\Delta$ in the entire Euclidean space $\Rd$
is easy,
since the differential operator proves to be unitarily equivalent to
the multiplication operator $\linA_m$ generated by the function
\[
m \: (x_1,\ldots,x_n) \mapsto (x_1^2 + \ldots + x_n^2) \= \abs{x} ^2
\quad (x \xeof \Rd).
\]
In fact,
the operator $\linA_m$ is real-valued and self-adjoint,
which is a consequence of
the \textit{Fourier-Plancherel transformation}\footnote{
\,We introduce the \textit{Fourier-Plancherel transformation} in $L^2(\Rd)$
by 
\[
\hat{u}(\xi) \eqdef
\frac{1}{(2 \pi)^{d/2}}
\int_{\,\Rd} \, e^{-2 \pi i \scalar{\xi}{x}} \, u(x) \, dx
\]
for all $u \xeof C_c^{\infty}(\Rd)$,
which forms a partial isometry $\FT \: u \mapsto \hat{u} $
from the space $C_c^{\infty}(\Rd)$ to the space $L^2(\Rd)$.
Then for functions $u,v \xeof L^2(\Rd)$ 
the inner product between $u$ and $v$
and the inner product of their Fourier-Plancherel transforms
$\hat{u}$ and $\hat{v}$ are related by \textsc{Plancherel}\textit{'s formula}
$\scalar{u}{v} \eq \scalar{ \hat{u} }{\hat{v} }$.
Since the mapping $\FT$ is linear and continuous on $C_c^{\infty}(\Rd)$,
which is a dense subspace of $L^2(\Rd)$,
it can be extended to an \textit{unitary},
that is,
a bijective and isometric linear operator
$\FT_P \: L^2(\Rd) \xto L^2(\Rd)$,
called the \textit{Fourier-Plancherel} \textit{operator}
or the \textit{$L^2$-Fourier transformation} in $L^2(\Rd)$.}
.
This implies
that the operator $\linA_{sa}$ with domain of definition
$\domain{\linA_{sa}} := H^2(\Rd)$ must be self-adjoint,
too.
The spectrum of $\linA_{sa}$ is that of the operator $\linA_m$,
that is,
it is purely continuous and $\sigma_c(\linA_{sa}) \eq [0,\infty)$.
Moreover,
since $\linA_{sa}$ is the closure of the operator $\linA_{-\Delta,0}$,
the latter proves to be essentially self-adjoint.

\item[(ii)]
\textit{Dirichlet-Laplacian}.
Let the domain $\G \xsubset \Rd$ be \textit{bounded}
and consider $-\Delta$ with domain of definition $\domain{-\Delta} := \TF$.
The related Dirichlet form $\DF_{-\Delta}$ is given by 
\[
\DF_{-\Delta}(\phi,\psi) \eqdef
\int_{\G} \nabla \phi \, \nabla \psi \, dx
\quad \quad
(\phi,\psi \xeof \TF),
\]
which can continuously be extended to the sesquilinear form
$\overline{\DF_{-\Delta}}$ with domain of definition
$\domain{\overline{\DF_{-\Delta}}} := \HN{1}{\G} \xtimes \HN{1}{\G}$.

\subparagraph{}
Due to theorem \ref{QUASI-COERCIVE_FORMS_AND_ASSOCIATED_OPERATORS},
the operator $\linA_D$ associated with the form $\DF_{-\Delta}$
and defined by
\[
\begin{cases}[1.2]
\quad \domain{\linA_D} \eqdef
\{u \xeof \HN{1}{\G} \: \Delta u \xeof L^2(\G)\} \\
\quad \linA_D u \eqdef - \Delta u \; \for u \xeof \domain{\linA_D} \\
\end{cases}
\] 
is self-adjoint,
called the \textit{negative Laplacian with Dirichlet boundary conditions}
or briefly the \textit{Dirichlet-Laplacian}.

\subparagraph{}
We obviously have
$\HN{1}{\G} \cap H^2(\G) \xsubset \domain{\linA_D}$,
and one can prove that
if $\G$ is bounded and owns a $C^2$-boundary or $\G \eq \R_+^n$,
then the reversed inclusion also holds,
hence $\domain{\linA_D} \eq \HN{1}{\G} \cap H^2(\G)$.
This equality is even true in the case,
where $\G$ is bounded and convex.

\subparagraph{}
Furthermore,
if $\G$ has the $1$-extension property,
then the space $\HN{1}{\G}$ is compactly embedded into the space $\Lh$
and $\linA_D$ is an operator with pure point spectrum.
In this case it is an important property of the Dirichlet-Laplacian
that the number $0$ does \textit{not} belong to the spectrum $\sigma(\linA_D)$,
so $\linA_D$ is invertible and the operator $\linA_D^{-1}$ proves to be bounded.

\item[(iii)]
\textit{Neumann-Laplacian}.
Let $\G \xsubset \Rd$ be a bounded domain with $C^1$-boundary $\partial \G$
and let the realization $\linA_N$ of the negative Laplacian in $\Lh$
be defined by
\[
\begin{cases}[1.2]
\quad \domain{\linA_N} \eqdef
\{ u \xeof H^2(\G) \: u_v(x)= 0, \; x \xeof \partial \G \} \\
\quad \linA_N u \eqdef - \Delta u \; \for u \xeof \domain{\linA_N}. \\
\end{cases}
\]
The realization $\linA_N$ is said to be the
\textit{(negative) Laplacian with Neumann boundary condition}
or briefly the \textit{Neumann-Laplacian}.
It owns $\lambda \eq 0$ as the minimal eigen\-value,
and all constant functions in $\domain{\linA_N}$ are eigen\-functions
of $\lambda \eq 0$.
\end{itemize}

\item[(c)] \textit{{\SCHROEDINGER} operators.}
By the \textit{{\SCHROEDINGER} operator} is meant
the linear partial differential expression given by
\[
\HAM_q(u,x) \eqdef (-\Delta u)(x) \+ q(x),
\]
defined in the entire Euclidean space $\Rd$.
In the special case,
where $q \eq 0$,
self-adjointness of the operator $-\Delta$ acting in $L^2(\Rd)$
and with domain of definition $H^2(\Rd)$ is well-known.

\subparagraph{} % potential functions and multiplication operators
The so-called \textit{potential function} $q \: \Rd \xto \R$
gives rise to a multiplication operator $M_q$ acting in $L^2(\Rd)$,
that is,
\[
\begin{cases}[1.2]
\quad \domain{M_q} \eqdef \{u \in L^2(\Rd) \: q {\cdot} u \in L^2(\Rd)\},\\
\quad M_q u        \eqdef q {\cdot} u \; \for u \in \domain{M_q}   ,\\
\end{cases}
\]
so the operator $\HAM_q$ can be seen as a perturbation
of the negative Laplacian by the operator $M_q$.

\paragraph{} %- Schroedinger equation
%------------------------------------------------------------------------------%
By the time-dependent \textit{{\SCHROEDINGER} equation} is meant
the initial value problem
\[
i \, \dot{u} \= -\Delta u \+ M_q u
\]
with prescribed initial value $u(0) \eq u_0 \xeof H^2(\Rd)$,
where it is supposed that $\HAM_q \eq -\Delta + M_q$
is a self-adjoint operator in $L^2(\Rd)$
with domain of definition $\domain{\HAM_q} \eq H^2(\Rd)$.
The existence and uniqueness of a solution for this problem is ensured
by the functional calculus for self-adjoint operators
using the spectral resolution of the operator $\HAM_q$
and the function $f:\lambda \xeof \R \mapsto e^{it \lambda}$.
The solution $u$ of the problem is then provided
by a group of unitary operators $\{U(t)\}_{t \in \R}$,
which implies $u(t) \eq U(t) u_0$ for every $t \xeof \R$
(see also section 5 for the treatment of the abstract {\SCHROEDINGER} equation).

\paragraph{} %- spectral analysis of Schroedinger operators
%------------------------------------------------------------------------------%
The study of \textit{time-independent solutions}
of the {\SCHROEDINGER} equation leads to the eigenvalue problem
\begin{equation}\label{SCHROEDINGER_EIGENVALUE_PROBLEM}
\HAM_q u \eqdef -\Delta u + M_q u \= \lambda u,
\end{equation}
which means to find those real numbers $\lambda$,
for which there are non-trivial solutions
of \eqref{SCHROEDINGER_EIGENVALUE_PROBLEM}.
The discrete spectrum $\sigma_d(\HAM_q)$ of the operator $\HAM_q$
is in fact of great interest,
since the eigen\-functions of eigen\-values $\lambda \xeof \sigma_d(\HAM_q)$
describe the stationary states of a quantum mechanical system.

\subparagraph{} %- spectrum of Schroedinger operators
The presence of a potential function $q$
can fairly change the spectrum of a {\SCHROEDINGER} operator 
like the example of the \textit{harmonic oscillator} shows:
the spectrum of the negative Laplacian is the connected set $[0,\infty)$,
but the spectrum of the operator $-\Delta {\,+\,} M_q$ 
with $q \: x \mapsto \frac{1}{2}x^2$ consists of discrete points only.

\subparagraph{} %- lower and upper bounds for the discrete spectrum
It is a further and interesting problem in quantum mechanics
to determine lower and upper bounds for the discrete spectrum
of the operator $\HAM_q$.
In the following we want to present three assertions,
which contain information on the bounds of the set $\sigma_d(\HAM_q)$,
namely
the \textit{Birman-Schwinger} \textit{bound}, 
the \textit{Lieb-Cwikel-Rozenblum} \textit{bound} and
    \textsc{Kato}\textit{'s theorem}:

\begin{itemize}[leftmargin=4mm]
\item[1.]
The \textit{Birman-Schwinger} \textit{bound.}
This assertion presents an upper bound for the set $\sigma_d(\HAM_q)$,
which is valid for certain potentials $q$ in case of $d \eq 3$.
Let $q \xeof L^{\infty}(\R^3)$ and 
\[
B_q \eqdef
\int_{\,\R^3} \int_{\,\R^3} \frac{q(x) \, q(y)}{\abs{x-y} ^2} \ dx \, dy
\, < \, + \infty.
\]
Then the total number $N(\HAM_q)$ of eigen\-values of the operator $\HAM_q$
can be estimated according to
\[
N(\HAM_q) \, \leq \, \frac{1}{16 \pi^2} \, B_q.
\]

\item[2.]
The \textit{Lieb-Cwikel-Rozenblum} \textit{theorem}
quotes an estimate on the number $N_(\HAM_q)$
of negative eigen\-values of the operator $\HAM_q$:
if $d > 3$ and $q \xeof L_{loc}^{\infty}(\Rd)$,
then
\[
N(\HAM_q) \, \leq \, c_d \int_{\,\Rd} \min \{q(x),0\}^{n/2} \, dx,
\]
where $c_d$ is a constant only depending on the dimension $d$.

\item[3.]
Our last result is often denoted as \textit{Kato}\textit{'s theorem}.
It does not make assertions on bounds for the discrete spectrum,
but quotes the absence of positive eigen\-values
under mild assumptions on the function $q$:
if $q \xeof L_{loc}^{\infty}(\Rd)$ and
\[
\lim_{\abs{x} \to \infty} \, \abs{x} \, q(x) \= 0,
\]
then the operator $\HAM_q$ does not own any positive eigen\-value.
\end{itemize}
\end{itemize}
%------------------------------------------------------------------------------%
%                                 BIBLIOGRAPHY                                 %
%------------------------------------------------------------------------------%
\newpage
\thispagestyle{empty}

\vspace{8mm}
\begin{thebibliography}{99}
\begin{small}

\bibitem{AT} W.\,Arendt, A.F.M. ter Elst:
\textit{From Forms to Semigroups}.
Operator Theory: Advances and Applications, Vol.\,221, pp.\,47-69.

\bibitem{Br1} F.\,Browder:
\textit{Functional Analysis and Partial Differential Equations I}.
Math. Annalen 138, pp.\,55-79, 1959.

\bibitem{Br2} F.\,Browder:
\textit{On the Spectral Theory of Elliptic Differential Operators I}.
Math. Annalen 142, pp.\,22-130, 1961.

\bibitem{EE}  D.\,E.\,Edmunds, W.\,D.\,Evans:
\textit{Spectral theory and differential operators}.
Oxford University Press, 1987.

\bibitem{EN}  K.-J.\,Engel, R.\,Nagel:
\textit{A Short Course on Operator Semigroups}.
Universitext, Springer, 2006.

\bibitem{GT}  D.\,Gilbarg, N.S.\,Trudinger:
\textit{Elliptic Partial Differential Operators of Second Order}.
Springer, 2001.

\bibitem{Go} J.\,Goldstein:
\textit{Semigroups of Linear Operators and Applications}.
Oxford University Press, 1985.

\bibitem{Ha}  M.\,Haase:
\textit{Lectures on Functional Calculus}.
21st International Internet Seminar, Kiel 2018.

\bibitem{Ka}  T.\,Kato:
\textit{{Perturbation theory of linear operators}}.
Springer, Grund\-lehren der mathematischen Wissenschaften, Band 132, 1966.

\bibitem{Kr}  S.\,Krein:
\textit{Linear Equations in Banach Spaces}.
Birk\-h\"{a}user, 1982.

\bibitem{Lu}  A.\,Lunardi:
\textit{Analytic Semigroups and Optimal Regularity in Parabolic Problems}.
Birk\-h\"{a}user, 1995.

\bibitem{Pa}  A.\,Pazy:
\textit{Semigroups of Linear Operators and
        Applications to Partial Differential Equations}.
Applied Mathematical Sciences Vol.\,44, Springer, 1983.

\bibitem{RR}  M.\,Renardy, R.\,Rogers:
\textit{An Introduction to Partial Differential Equations}.
Texts in Applied Mathematics 13, Springer, 2004.

\bibitem{RS1} M.\,Reed, B.\,Simon:
\textit{Modern Methods of Mathematical Physics I: Functional Analysis}.
Academic Press Inc., 1980.

\bibitem{Sc}  M.\,Schechter:
\textit{Fredholm Operators and the Essential Spectrum}.
New York University, Courant Institute of Mathematical Sciences,
January 1965.

% \bibitem{Sh}  R.\,Showalter:
% \textit{Hilbert Space Methods in Partial Differential Equations}.
% Dover Publications, 2010.

\bibitem{Wl}  J.\,Wloka:
\textit{Partial Differential Equations}.
Cambridge University Press, 1987.

\bibitem{Yo}  K.\,Yosida:
\textit{Functional Analysis}.
Springer, 1980.
\end{small}
\end{thebibliography}

%------------------------------------------------------------------------------%
%                               END OF DOCUMENT                                %
%------------------------------------------------------------------------------%
\end{document}

%------------------------------------------------------------------------------%
%                              TABLE OF CONTENTS                               %
%------------------------------------------------------------------------------%

 1. BANACH AND HILBERT SPACES
    -------------------------
    normed linear spaces
    sequentially complete normed spaces, Banach spaces
    inner products and Hilbert spaces
    the inner geometry of a Hilbert space

 2. BASIC THEORY OF LINEAR OPERATORS
    --------------------------------
    generalities on linear operators
    classification of linear operators by topological properties
      - continuous operators
      - compact operators
      - closed operators
    comparison of continuous and closed linear operators
    closed graph theorem

 3. ADJOINT SPACES AND OPERATORS
    ----------------------------
    classification of operators by choices of the range space
    Riesz's representation theorem
    adjoint operators
    the Hilbert adjoint
    normal operators in Hilbert space

 4. RESOLVENT AND SPECTRUM
    ----------------------
    resolvent set and spectrum
    resolvent functions
    various kinds of spectra
    spectrum and functional calculus for bounded operators
    projections in a Banach space

 5. NORMALLY SOLVABLE OPERATORS
    ---------------------------
    Banach's closed range theorem
    semi-Fredholm operators
    perturbation theory of semi-Fredholm operators
    Fredholm operators
    essential spectrum

 6. SYMMETRY AND SELF-ADJOINTNESS
    -----------------------------
    numerical range and external field of regularity
    symmetric operators
    spectrum of symmetric operators
    self-adjoint operators
    perturbation theory of self-adjoint operators
    multiplication operators

 7. SPECTRAL THEORY OF SELF-ADJOINT OPERATORS
    -----------------------------------------
    the spectrum of self-adjoint operators
    spectral resolutions
    spectral operators
    the spectral theorem
    the functional calculus for self-adjoint operators
    evolution equations governed by self-adjoint operators

 8. OPERATORS WITH PURE POINT SPECTRUM
    ----------------------------------
    discrete spectrum
    Riesz-Schauder theory for compact linear operators
    operators with compact resolvents
    the discrete spectral theorem
    concepts of solvability for linear equations

 9. INFINITESIMAL GENERATORS OF SEMIGROUPS
    --------------------------------------
    strongly continuous semigroups of bounded linear operators
    infinitesimal generators
    well-posedness of first order Cauchy problems
    generator theorems for strongly continuous semigroups

10. ACCRETIVE OPERATORS
    -------------------
    concepts of accretivity
    Lumer-Phillips theorem
    well-posedness of second order Cauchy problems
    regularly accretive operators

11. SECTORIAL OPERATORS AND ANALYTIC SEMIGROUPS
    -------------------------------------------
    definition of sectorial operators
    the functional calculus for sectorial operators
    analytic semigroups
    the generator theorem for analytic semigroups
    analytic semigroups created by quasi-coercive forms

12. APPLICATIONS TO LINEAR DIFFERENTIAL OPERATORS
    ---------------------------------------------
    linear partial differential operators
    free space and boundary value problems
    operators of elliptic, parabolic and hyperbolic type
    realizations in the spaces C^{k,\gamma}(G) and L^2(G)
    the closeness of strongly elliptic Dirichlet operators
    formally adjoint expressions and adjointness theorems
    Gardings's inequality and quasi-coercive Dirichlet forms
    weak solutions of boundary and initial-boundary value problems
    the Fredholm property of elliptic operators
    spectral theory and sectorialness of strongly elliptic operators
    self-adjoint extensions of half-bounded operators
